\documentclass{tufte-book}

\hypersetup{colorlinks}

\title{Modern Analysis}
\author[The Tufte-LaTeX Developers]{Claire A. Saint-Donat}
\publisher{Columbia University, Fall 2013}
\usepackage{lipsum}


\usepackage{booktabs}
\usepackage{amsmath, amsthm, amssymb}
\usepackage{amsmath}
\usepackage{amssymb}
\usepackage{amsfonts}
\usepackage{url}
\usepackage{xy}
\usepackage{mathrsfs}
\usepackage{hyperref}
\usepackage{verbatim}
\usepackage{enumerate}
\usepackage{multicol}
\usepackage{manfnt}
\usepackage{amsthm}
\usepackage{titletoc}
\usepackage{amssymb}
\usepackage{graphicx}
\setkeys{Gin}{width=\linewidth,totalheight=\textheight,keepaspectratio}
\graphicspath{{graphics/}}


\usepackage{fancyvrb}
\fvset{fontsize=\normalsize}


\setcounter{chapter}{1}
\setcounter{secnumdepth}{2}


\newtheorem{theorem}{Theorem}[chapter]
\newtheorem{lemma}[theorem]{Lemma}
\newtheorem{corollary}[theorem]{Corollary}
\newtheorem{proposition}{Proposition}[chapter]
\theoremstyle{definition}
\newtheorem{definition}{Definition}[chapter]
\newtheorem*{remark}{Remark}
\newtheorem*{example}{Example}
\newtheorem*{imp}{Important point}
\newtheorem*{dangerousbend}{\textdbend Caution}
\newtheorem*{summary}{Summary}
\renewcommand{\thesection}{\thechapter.\arabic{section}}
\numberwithin{section}{chapter}


\newcommand{\hangp}[1]{\makebox[0pt][r]{(}#1\makebox[0pt][l]{)}}

\newcommand{\hangstar}{\makebox[0pt][l]{*}}

\usepackage{xspace}

\newcommand{\vdqi}{\textit{VDQI}\xspace}
\newcommand{\ei}{\textit{EI}\xspace}
\newcommand{\ve}{\textit{VE}\xspace}
\newcommand{\be}{\textit{BE}\xspace}
\newcommand{\VDQI}{\textit{The Visual Display of Quantitative Information}\xspace}
\newcommand{\EI}{\textit{Envisioning Information}\xspace}
\newcommand{\VE}{\textit{Visual Explanations}\xspace}
\newcommand{\BE}{\textit{Beautiful Evidence}\xspace}

\newcommand{\TL}{Tufte-\LaTeX\xspace}

\newcommand{\monthyear}{%
  \ifcase\month\or January\or February\or March\or April\or May\or June\or
  July\or August\or September\or October\or November\or
  December\fi\space\number\year
}


\newcommand{\openepigraph}[2]{%
  %\sffamily\fontsize{14}{16}\selectfont
  \begin{fullwidth}
  \sffamily\large
  \begin{doublespace}
  \noindent\allcaps{#1}\\% epigraph
  \noindent\allcaps{#2}% author
  \end{doublespace}
  \end{fullwidth}
}


\newcommand{\blankpage}{\newpage\hbox{}\thispagestyle{empty}\newpage}

\usepackage{units}

% Typesets the font size, leading, and measure in the form of 10/12x26 pc.
\newcommand{\measure}[3]{#1/#2$\times$\unit[#3]{pc}}

% Macros for typesetting the documentation
\newcommand{\hlred}[1]{\textcolor{Maroon}{#1}}% prints in red
\newcommand{\hangleft}[1]{\makebox[0pt][r]{#1}}
\newcommand{\hairsp}{\hspace{1pt}}% hair space
\newcommand{\hquad}{\hskip0.5em\relax}% half quad space
\newcommand{\TODO}{\textcolor{red}{\bf TODO!}\xspace}
\newcommand{\ie}{\textit{i.\hairsp{}e.}\xspace}
\newcommand{\eg}{\textit{e.\hairsp{}g.}\xspace}
\newcommand{\na}{\quad--}% used in tables for N/A cells
\providecommand{\XeLaTeX}{X\lower.5ex\hbox{\kern-0.15em\reflectbox{E}}\kern-0.1em\LaTeX}
\newcommand{\tXeLaTeX}{\XeLaTeX\index{XeLaTeX@\protect\XeLaTeX}}
% \index{\texttt{\textbackslash xyz}@\hangleft{\texttt{\textbackslash}}\texttt{xyz}}
\newcommand{\tuftebs}{\symbol{'134}}% a backslash in tt type in OT1/T1
\newcommand{\doccmdnoindex}[2][]{\texttt{\tuftebs#2}}% command name -- adds backslash automatically (and doesn't add cmd to the index)
\newcommand{\doccmddef}[2][]{%
  \hlred{\texttt{\tuftebs#2}}\label{cmd:#2}%
  \ifthenelse{\isempty{#1}}%
    {% add the command to the index
      \index{#2 command@\protect\hangleft{\texttt{\tuftebs}}\texttt{#2}}% command name
    }%
    {% add the command and package to the index
      \index{#2 command@\protect\hangleft{\texttt{\tuftebs}}\texttt{#2} (\texttt{#1} package)}% command name
      \index{#1 package@\texttt{#1} package}\index{packages!#1@\texttt{#1}}% package name
    }%
}% command name -- adds backslash automatically
\newcommand{\doccmd}[2][]{%
  \texttt{\tuftebs#2}%
  \ifthenelse{\isempty{#1}}%
    {% add the command to the index
      \index{#2 command@\protect\hangleft{\texttt{\tuftebs}}\texttt{#2}}% command name
    }%
    {% add the command and package to the index
      \index{#2 command@\protect\hangleft{\texttt{\tuftebs}}\texttt{#2} (\texttt{#1} package)}% command name
      \index{#1 package@\texttt{#1} package}\index{packages!#1@\texttt{#1}}% package name
    }%
}% command name -- adds backslash automatically
\newcommand{\docopt}[1]{\ensuremath{\langle}\textrm{\textit{#1}}\ensuremath{\rangle}}% optional command argument
\newcommand{\docarg}[1]{\textrm{\textit{#1}}}% (required) command argument
\newenvironment{docspec}{\begin{quotation}\ttfamily\parskip0pt\parindent0pt\ignorespaces}{\end{quotation}}% command specification environment
\newcommand{\docenv}[1]{\texttt{#1}\index{#1 environment@\texttt{#1} environment}\index{environments!#1@\texttt{#1}}}% environment name
\newcommand{\docenvdef}[1]{\hlred{\texttt{#1}}\label{env:#1}\index{#1 environment@\texttt{#1} environment}\index{environments!#1@\texttt{#1}}}% environment name
\newcommand{\docpkg}[1]{\texttt{#1}\index{#1 package@\texttt{#1} package}\index{packages!#1@\texttt{#1}}}% package name
\newcommand{\doccls}[1]{\texttt{#1}}% document class name
\newcommand{\docclsopt}[1]{\texttt{#1}\index{#1 class option@\texttt{#1} class option}\index{class options!#1@\texttt{#1}}}% document class option name
\newcommand{\docclsoptdef}[1]{\hlred{\texttt{#1}}\label{clsopt:#1}\index{#1 class option@\texttt{#1} class option}\index{class options!#1@\texttt{#1}}}% document class option name defined
\newcommand{\docmsg}[2]{\bigskip\begin{fullwidth}\noindent\ttfamily#1\end{fullwidth}\medskip\par\noindent#2}
\newcommand{\docfilehook}[2]{\texttt{#1}\index{file hooks!#2}\index{#1@\texttt{#1}}}
\newcommand{\doccounter}[1]{\texttt{#1}\index{#1 counter@\texttt{#1} counter}}

% Generates the index
\usepackage{makeidx}
\makeindex

\begin{document}

\frontmatter

\maketitle

\newpage\thispagestyle{empty}
\openepigraph{%
Real mathematics must be justified as art, if it can be justified at all.
}{G.H. Hardy
}
\vfill
\openepigraph{%
The natural numbers are the work of God.  All the rest is the work of mankind.
}{Leopold Kronecker}
\vfill
\openepigraph{%
\ldots It has been our experience that a mathematical question resolved by a counterexample has the pungency of good drama.  Many of the most elegant and artistic contributions to mathematics belong to this genre.
}{Bernard R. Gelbaum}


\vfill
\openepigraph{%
If people do not believe that mathematics is simple, it is only because they do not realize how complicated life is.}{ John von Neumann}

\vfill
\openepigraph{%
God exists since mathematics is consistent, and the Devil exists since we cannot prove it.}{Andr\'{e} Weil}

%\vfill
%\openepigraph{%A mathematician is a person who can find analogies between theorems; a better mathematician is one who can see analogies between proofs and the best mathematician can notice analogies between theories. One can imagine that the ultimate mathematician is one who can see analogies between analogies.}{Stefan Banach}


% r.3 full title page
%\maketitle


% r.5 contents
\tableofcontents

%\listoffigures

%\listoftables

% r.7 dedication
%\cleardoublepage
%~\vfill
%\begin{doublespace}
%\noindent\fontsize{18}{22}\selectfont\itshape
%\nohyphenation
%Dedicated to those who appreciate \LaTeX{} 
%and the work of \mbox{Edward R.~Tufte} 
%and \mbox{Donald E.~Knuth}.
%\end{doublespace}
%\vfill
%\vfill
%
%
%
%\cleardoublepage




%\maketitle

\chapter*{Introduction}
\noindent This document is written in the Tufte handout \LaTeX\ document style.
It provides examples and definitions from Walter Rudin's famous text on Mathematical Analysis which I read in the Fall of 2013 at Columbia University.  What follows are my notes and observations from this critical mathematical work.




\bookchapter{One}{The Real and Complex Number Systems}

In order to have a fruitful discussion on the main concepts of analysis (namely, convergence, continuity, differentiation, and integration) we must first begin with an accurately-defined \textsc{number} concept.

We shall assume all of the axioms that govern arithmetic of the integers and assume familiarity with the rational numbers.  

We will demonstrate how the rational number system is inadequate for many purposes, both as a \textsc{field} and as an \textsc{ordered set}.   (These terms will be defined shortly).   This leads to a natural introduction of the "irrational numbers."

\newthought{The purpose of our discussion of the real numbers}, will be to show that the rational number system has certain gaps, in spite of the fact that between any two rationals there is another. \marginnote{If $r<s$ then $r< \frac{(r+s)}{2} < s$.}  The real number system on the other hand, fills these gaps.  This is the principle reason for the fundamental role it plays in analysis. 

\bigskip

\section{The Axiom of Completeness}
\subsection{Ordered Sets}
%\marginnote{Here are some standard set-theory terms that will be used throughout:  If $A$ is any set (whose elements may be numbers or any other objects), we write $x \in A$ to indicate that $x$ is a member of $A$.   If $x$ is not a member of $A$, we write that $x \notin A$. \\  The set which contains no element will be called the \textsc{empty set}.   If a set has at least one element, it is called nonempty. \\  If $A$ and $B$ are sets, and if every element of $A$ is an element of $B$, we say that $A$ is a subset of $B$ and write $A \subset B$ or $B \supset A$.  If, in addition, there is another element of $B$ not in $A$ then we say that $A$ is a proper subset of $B$.   Note that $A \subset A$ for every set $A$ \\ If $A \subset B$ and $B\subset A$ , we write that $A=B$.  Otherwise $A \neq B$.  }
In order to elucidate the structure of the real numbers as well as that of the complex numbers, we must start with a brief discussion of the general concepts of ordered set and field.



\begin{definition}Let $S$ be a set.  An \textsc{order} on $S$ is a relation, denoted by $<$, with the following two properties

\begin{itemize}
	\item  If $x \in S$ and $y \in S$ then one and only one\footnote{Trichotomy: Between any two elements in the set, only one of the relations can hold.  There is a unique relationship between any two elements.} of the following statements is true:	\[	x< y \qquad x=y \qquad y<x\]  
	\item If $x,y,z \in S$ and, if $x<y$ and $y<z$ then $x<z$
	\footnote{Transitivity}
\end{itemize}
\end{definition}


The statement "$x<y$" may be read as "x is less than y" or "x is smaller than y" or even "x precedes y."

The notation $x \leq y$ indicates that $x<y$ or $x=y$ without specifying which of these two is in fact the case.  In other words $x \leq y$ is the negation of $y<x$.


\begin{definition}An {ordered set} is a set $S$ in which an order is defined.  For example, Q is an ordered set if $r<s$ is defined to mean that $s-r$ is a positive rational number. \end{definition}



\begin{definition}Suppose that $S$ is an ordered set and that $E \subset S$.  If there exists a $\beta \in S$ such that $x \leq \beta$ for every $x \in E$ then we say that $E$ is \textsc{bounded above} and call that $\beta$ an \textsc{upper bound} of $E$.  Lower bounds are similarly defined. \end{definition}


\begin{definition} Suppose that $S$ \marginnote{An alternative and often useful definition of supremum is as follows:  Assume $\alpha \in S$ and is an upper bound for a set $E \subset S$.  Then $\alpha= \sup E$ iff $\forall \quad \varepsilon>0 $ there exists an element $x \in E$ such that $\alpha-\varepsilon < x.$}
is an ordered set, $E \subset S$ and $E$ is bounded above.    Suppose that there exists an $\alpha \in S$ with the following properties:
\begin{itemize}
	\item $\alpha$ is an upper bound of $E$
	\item If $\gamma< \alpha$ then $\gamma$ is not an upper bound of $E$
\end{itemize}
The $\alpha$ is called the \textsc{least upper bound} \footnote{that there is a most one such $\alpha$ is clear from the second premise} or the \textsc{supremum of $E$} and we write
\[\alpha = \sup E\]
\end{definition}

The greatest lower bound of infimum of a set which is bounded below if defined in a similar manner.




\begin{figure}
  \includegraphics{SupInf.pdf}
%  \checkparity This is an \pageparity\ page.%
  \caption{An example of Supremums and Infimums.}
  %\label{fig:textfig}
  %\zsavepos{pos:textfig}
  \setfloatalignment{b}
\end{figure}

\section{Examples}
Let us now show that $\sqrt{2}$ is irrational.

\begin{proposition} There is no such rational $p$ such that $p^2 = 2$. \end{proposition}

\begin{proof}  If there were such a $p$, we could write $p=m/n$ where $m$ and $n$ are both integers that are not both even (alternatively that $\gcd(m,n) = 1$).  Let us assume this is done.\\
Then we have that $m^2 = 2n^2$.  Hence $m$ is even. \marginnote{$k^2$ is even $\iff$ $k$ is even for all $k \in \mathbb{Z}$}.  Then $m=2k$ for some $k$ and $m^2 = (2k)^2 = 4k^2 = 2n^2$.  Therefore $n^2$ is also even and so therefore is $n$.\\
This conclusion is contrary to our choice of $m$ and $n$.  Hence there exists no such rational $p$. \end{proof}

Let's now examine the situation a little more closely.   Let $A$ be the set of all positive rationals such that $p^2 < 2$ and let $B$ consist of all the positive rationals such that $q^2 > 2$.  We can see that $A$ contains no largest member and $B$ contains no smallest, despite the fact that they are bounded above and below respectively.  Therefore $A$ has no least upper bound in $\mathbb{Q}$ and $B$ has no greatest lower bound in $\mathbb{Q}$.  \\

\newthought{Let $E$} consist of all numbers $1/n$ where $n=1,2,3,\ldots$.  Then, $\sup E = 1 \in E$ and $\inf E =0$ which is not in $E$.  This is to demonstrate that infimum and supremum may or may not be included in the sets which they bound.


\section{Least Upper Bound Property}
An ordered set $S$ is said to have the \textsc{least-upper-bound property} is the following is true:

\begin{fullwidth}
If $E \subset S$, $E$ is not empty, and $E$ is bounded above, then $\sup E $ exists in $S$.
\end{fullwidth}

Abbott calls this the Axiom of Completeness and purports that it is the final, most distinctive assumption about the real number system.

For example, we have already shown by a previous example that $\mathbb{Q}$ does not have the least upper bound property.


We shall now show with a central theorem that there is a close relation between greatest lower bounds and least upper bounds, and that every ordered set with the least upper bound property also has the greatest lower bound property.

\begin{theorem} \marginnote{Every ordered set with the least upper bound property, also has the greatest lower bound property.}
Suppose that S is an ordered set with the least-upper bound property, $B \subset S$, $B$ is not empty and $B$ is bounded below.

Let $L$ be the set of all lower bounds of $B$.  Then
\[\alpha = \sup L\]
exists in $S$ and $\alpha = \inf B$.

In particular, $\inf B$ exists in $S$.
\end{theorem}

\begin{proof}
Since $B$ is bounded below, $L$ is not empty.  Since $L$ consists of exactly those $y \in S$ which satisfy the inequality $y \leq x$ for every $x \in B$, we see that every $x \in B$ is an upper bound of $L$.  Thus $L$ is bounded above.  Our hypothesis about $S$ implies therefore that $L$ has a supremum in $S$; call it $\alpha$.

If $\gamma < \alpha$ then $\gamma$ is not an upper bound of $L$, hence $\gamma \notin B$.  It follows that $\alpha \leq x$ for every $x \in B$.  Thus $\alpha \in L$.    

If $\alpha < \beta$ then $\beta \notin L$,  since $\alpha$ is an upper bound of $L$.

We have shown that $\alpha \in L$ but $\beta \notin L$ if $\beta> \alpha$.  In other words, $\alpha$ is a lower bound of $B$, but $\beta$ is not if $\beta>\alpha$.  This means that $\alpha = \inf B$.
\end{proof}

\section{Consequences of Completeness}

Complete metric spaces often have certain nice properties as a result of having least upper bounds.  What follows are some of the consequences of the Completeness Axiom.

For simplicity, we often assume that complete metric space we are working in is the real number line, but be advised that these properties will hold in any metric space with similar properties.

\begin{theorem}\marginnote{Nested Interval Property}  For each $n \in N$ assume we are given a closed interval $I_n = [a_n , b_n] = \{ x \in R : a_n \leq x \leq b_n\}$.  Assume also that each $I_n$ contains $I_{n+1}$.  Then, the resulting sequence of nested intervals, 

\[I_1 \supseteq I_2 \supseteq I_3 \supseteq I_4 \supseteq \cdots \]

has a nonempty intersection; that is, $\cap_{n=1}^{\infty} I_n \neq \emptyset$.
\end{theorem}

\begin{proof}
Knowing that we want to use the AoC to produce a single real value  $x$ satisfying $x \in I_n$ for all $N$, we should consider a bounded set such that its supremum might be this $x$.  Let us consider the set 
\[ A = \{a_n : n \in N\}\]
of the left-handed endpoints of the intervals.

\begin{marginfigure}
  \includegraphics{NIP.pdf}
  \caption{An example of the Nested Interval Property.}
 % \label{fig:textfig}
\end{marginfigure}


Because the intervals are nested, we see that every $b_n$ serves as an upper bound for $A$.  Thus, we are justified in setting
\[ x = \sup A.\]

Now, consider a particular $I_n = [a_n , b_n]$.  Because $x$ is an upper bound for $A$, we have $a_n \leq x$.  The fact that each $b_n$ is an upper bound for $A$ and that $x$ is the least upper bound implies $x\leq b_n$.

Altogether then, we have $a_n \leq x \leq b_n$, which means $x \in I_n$ for every choice of $n \in N$.  Hence, $x \in \cap_{n=1}^{\infty} I_n$ and the intersection is not empty.
\end{proof}


\section{The Density of Q and C in R}
Again here, we are restricting our discussion to the real numbers.  

\begin{theorem}\marginnote{Archimedean Property}  \begin{itemize}[(i)] 
\item  Given any number $x \in R$ there exists an $n \in N$ such that $x<n$.  \item Given any real number $y>0$, there exists an $ n \in N$ such that $\frac{1}{n}< y$. \end{itemize} \end{theorem}

\begin{proof}   Part (i) of the proposition states that N is not bounded above.  This might seem obvious and not in need of a proof yet it is important to note that a set might posses the Archimedean property without being Complete- take Q for example- but a demonstration of this fact requires an investigation into the axiomatic construction of the ordered field in question. In the case of R, the AoC furnishes us with a very short argument.

Assume that N is bounded above.  By the Axiom of Completeness, N should then have a least upper bound property, and we can set $\alpha = \sup N$.  If we were to consider $\alpha -1$, then we no longer have an upper bound and therefore there exists an $n \in N$ such that $\alpha - 1 < n$.  But this is equivalent to $\alpha < n+1$. But because $n+1 \in N$ we have a contradiction to the fact that $\alpha $ is meant to be  an upper bound for $N$.  \marginnote{The contradiction here rests on the AoC and the fact that N is closed under addition.}

For part (ii) just let $x = 1/y$ from above.
Note that $\frac{1}{n} < y$ implies that $\frac{1}{y} < n$.  Therefore this proposition again states that N is not bounded above and the proof is exactly as before.
\end{proof}

\begin{theorem}  \marginnote{Density of C and Q in R}  

\begin{itemize}[(i)] 
\item    For every two real number $a$ and $b$ with $a<b$, there exists a rational number $r$ satisfying $a<r<b$
\item Given any two real numbers $a$ and $b$ with $a<b$. there exists an irrational number $t$ satisfying $a<t<b$
 \end{itemize} \end{theorem}
 
 \begin{proof}
 \begin{itemize}[(i)]
 	\item WLOG assume that  $0 \leq a \leq b$.  A rational number is a quotient, therefore there exists $n,m \in N$ such that 
	\[	a  < \frac{m}{n} < b	\]
 The first step is to choose a denominator large enough so that consecutive increments of size $1/n$ are too close together to "step over" the interval $(a,b)$.
 
 Using the Archimedean property, we pick $n \in N$ so that 
 \[		\frac{1}{n} < b-a\]
The first inequality gives us $na < m <nb$.  With $n$ already chosen, we choose $m$ to be the smallest integer great that $na$.  In other words.
$$m-1 \leq na < m$$ 
which yields $a < m/n$.  The second equality is also equivalent to $a< b - 1/n$ and we can use the third to write,
$$  m \leq na +1$$
$$<  n \big( b - \frac{1}{n} \big) + 1 $$
$$ = nb $$

Because $m<nb$ implies that $m/n < b$, and we have that $a<m/n <b$ as desired.

\item  We want to prove the existence of an irrational number between any two real numbers $a$ and $b$.  Using the density of $Q$ in $R$ on the real number $a - \sqrt{2}$ and $b - \sqrt{2}$ we can find a rational number $r$ such that $a - \sqrt{2} < r < b - \sqrt{2}$.  This then implies that $a< r + \sqrt{2} < b$.  Since an irrational plus a rational is again an irrational (a result proved in the exercises) it follows that $r + \sqrt{2}$ is irrational.

 
 \end{itemize}
 \end{proof}





\section{Fields}
%Include the field axioms
\subsection{The Real Field}

We now state the existence theorem which is the core of this chapter.

\begin{theorem}  \marginnote{Here, we only state the existence of $\mathbb{R}$ but we do not yet provide a proof of its existence.}There exists an ordered field $R$ which has the least-upper bound property.   Moreover, $R$ contains $Q$ as a subfield. \end{theorem}


We omit the proof since it is a little tedious.



\section{Exercises}

\allcaps{1.}  If $r$ is nonzero and rational and $x$ is irrational, show that $rx$ and $r+x$ are also irrational. 

\begin{proof}
	i) Assume that $rx$ is rational for contradiction, where $r$ and $x$ are rational and irrational respectively.  Then assume $r = \frac{a}{b}$
 for some $a,b \in Z$, $b \neq 0$
 \[x \cdot \frac{a}{b} = \frac{c}{d} \leftrightarrow x = \frac{bc}{ad}\]
 contradicting our assumption that $x$ is irrational.
 
 ii) Assume for contradiction that $r+x$ is rational. Therefore $x = (r+x) - r$ is the difference between two rational numbers and therefore is rational itself.  This contradicts our assumption that $x$ is irrational.
 \end{proof}
 
\allcaps{4.}  Let E be a nonempty subset of an ordered set; suppose that $\alpha$ is a lower bound of E and that $\beta$ is an upper bound of E.  Prove that $\alpha \leq \beta$.

\begin{proof}
%Include Proof
\end{proof}


\allcaps{5.}  Let $A$ be a nonempty set of real numbers which is bounded below.  Let $-A$ be the set of all numbers $-x$ where $x$ is in $A$.  Prove that $\inf A = -\sup(-A)$.


\begin{proof}
Using the fact that $R$ is an ordered field with the least upper bound property, we can assume that a nonempty subset $A \subseteq R$ that is bounded below has a greatest lower bound or infimum in $R$.  Call it $\alpha$.

Taking an arbitrary member of this subset $x \in A$, we know that $-x \in -A$ by definition of membership in $-A$.  Furthermore by the definition of infimum, we know that $\alpha \leq x$ for all $x\in A$.  As a consequence of Proposition 1.18 (a) we can deduce that $-\alpha \geq -x$ for all $-x \in -A$.

Using an arbitrary element in $-A$, we have shown that $-\alpha \geq -x$ for all $-x \in -A$ which suggests that $-\alpha$ is an upper bound for $-A$.  

What remains to be shown is that $-\alpha$ is a \emph{least} upper bound.  

Choosing some arbitrary element $\beta < -\alpha$, using again Proposition 1.18 (a), we can be sure that $-\beta > \alpha$.  Since $\alpha $, we already know, is the greatest lower bound of $A$, an element greater than $\alpha$ is not a lower bound but a member of $A$.  By the density of the real numbers, there must exist some element of $x \in A$ such that  $x < -\beta$.  Using yet again Proposition 1.18 (a), there exists some element $-x \in -A$ such that $-x > \beta$.  This states strictly then that $\beta$ \emph{cannot} be an upper bound of $-A$.

Therefore, we have shown that any element LESS than $-\alpha$ cannot be an upper bound for $-A$.  This together with the proof that $-\alpha$ is an upper bound for $-A$ shows that $-\alpha$ is the Supremum of $-A$.


By Proposition 1.14 (d), we can be sure that 
$$ -\sup(-A) = -(-\alpha) = \alpha =  \inf(A) $$
as desired.
\end{proof}



\section{Cardinality \& Countability}
\begin{definition}\textsc{Function:}  Consider two sets $A$ and $B$ and suppose that with each element $x$ of $A$ there is associated, in some manner, an element of $B$, which we denote by $f(x)$.  Then $f$ is said to be a \textsc{function} from $A$ to $B$ (or a mapping of $A$ into $B$).  The set $A$ is called the \textsc{domain} of $f$  (we also say that $f$ is defined on $A$) and the elements $f(x)$ are called the values of $f$. The set of all values of $f$ is called the \textsc{range} of $f$.\end{definition}
\bigskip

\begin{definition}\textsc{Image:}  Let $A$ and $B$ be two sets and let $f$ be a mapping of $A$ into $B$.   If $E \subset A$,  $f(E)$ is defined to be the set of all elements $f(x)$ for $x \in E$.   We call $f(E)$ the \textsc{image} of E under f.  \\In this notation, f(A) is the range of f.  It is clear that $f(A) \subset B$.  \\If $f(A) = B$, then we say that f maps \emph{onto} B.  That is, the function is \textsc{surjective}.\end{definition}
\bigskip


\begin{definition}\textsc{Inverse Image:}  If $E \subset B, f^{-1}(E)$ denotes the set of all $x \in A$ such that $f(x) \in E$.  We call $f^{-1}(E)$ the \textsc{inverse image} of $E$ under $f$.  \\
\marginnote{Alternatively, the concept of 1-1 can be expressed as follows:  $f$ is a 1-1 mapping of A into B provided that $f(x_{1}) \neq f(x_{2})$ whenever $x_{1} \neq x_{2}, x_{1} \in A, x_{2} \in A $}


If $y \in B, f^{-1}(y)$ is the set of all $x\in A$ such that $f(x) = y$.  If for each $y \in B, f^{-1}(y)$ consists of at most one element of $A$, then $f$ is said to be a \emph{1-1 (one-to-one) mapping} of A into B. \end{definition}
\bigskip


\begin{definition}\textsc{Cardinality:}  If there exists a 1-1 mapping of A onto B, we say that A and B can be put into 1-1 correspondence, or that A and B have the same \textsc{cardinal number}, or, briefly that A and B are \emph{equivalent}, and we write that $A \sim B $.  This relation is also an equivalence relation and therefore has the following properties:\marginnote{For two finite set A and B, we have $A \sim B$ iff A and B contain the same number of elements. \\
	For infinite sets, however this is a useful definition since the notion of 1-1 correspondence often maintains its clarity.}
	\begin{itemize}
	\item It is reflexive:  $A \sim A$
	\item It is symmetric:  If $A \sim B$ then $B \sim A$
	\item It is transitive:  If $A \sim B$ and $B\sim C$ then $A \sim C$
\end{itemize}\end{definition}
\bigskip


\begin{definition}\textsc{Countability:}  For a set A
\begin{itemize}
	\item A is \textsc{finite} if $A\sim [n]$ for some n \marginnote{Note: the empty set is also considered finite! \quad $A$ is infinite if $A$ is equivalent to one of its proper subsets.}
	\item A is \textsc{infinite} if $A$ is not finite
	\item A is \textsc{countable} if $A \sim N$
	\item A is \textsc{uncountable} if A is neither finite nor countable
	\item A is \textsc{at most countable} if A is either finite or countable
	
\end{itemize}\end{definition}

Notice that $\mathbb{Z}$ is countable.  (We demonstrate this by defining the bijection). 
\marginnote{An example of a 1-1 function from $N \to Z$:
\begin{displaymath}
   f(x) = \left\{
     \begin{array}{lr}
       (n-1)/2 & : n \quad \text{is odd}\\
       -n/2 & : n \quad \text{is even}
     \end{array}
   \right.
\end{displaymath}} 
Also notice that $\mathbb{Q}$ is countable. (We demonstrate this using the Cantor construction).  Likewise, the set of even natural numbers is countable.  We can show this by inventing an appropriate function.


Also notice,  the function $f(x) = \frac{x}{x^2 - 1}$ takes the interval $(-1,1)$ onto $R$ in a 1-1 fashion.   Thus $(-1,1) \sim R$.   In fact, $(a,b) \sim R$ for any interval $(a,b)$.      Therefore, intervals on the real line are uncountable.



A finite set cannot be equivalent to one of its proper subsets.   That is, however possible for infinite sets.  In fact we could replace the definition for infinite with this statement:  A is infinite if A is equivalent to one of its proper subsets.
\bigskip


\noindent \begin{definition}\textsc{Sequence:} \marginnote{Simply put, a sequence is a function whose domain is $\mathbb{N}$} \end{definition}By a sequence, we mean a function $f$ defined on the set of natural numbers.   If $f(n) = x_{n}$ for $n \in \mathbb{N}$, it is customary to denote the sequence f by the symbol $\{x_{n}\}$, or sometimes by $x_{1}, x_{2}, \ldots$. The values of $f$, that is, the elements of the sequence are called terms.  If $A$ is a set and if $x_{n} \in A$ for all $n \in J$, then $\{x_{n}\}$ is said to be a sequence in A, or a sequence of elements of A.\\

Note that the terms of a sequence need not be distinct.\\


\smallskip
Since every countable set is the range of a bijection defined on the set of natural numbers, we can regard every countable set as the range of a sequence of distinct terms.  In other words, the elements of a countable set can be "arranged in a sequence."


\smallskip

\begin{theorem}  If $A$ is countable and $E \subset A$ then $E$ is at most countable, that is finite or countable.  \marginnote{Alternatively, Every infinite subset of a countable set is countable} \end{theorem}

\begin{proof}
Let $E \subset A$.
\indent If $E$ is finite, then we are done. \\
\indent On the other hand assume that $E$ is infinite.  Then let $A = \{x_{1}, x_{2}, \ldots\}$ which we can do since A is countable.  Then we construct a sequence $\{n_{k}\}$ as follows:  Let $n_{1}$ be the smallest natural number such that $x_{n_{1}} \in E$.   Having chosen, $n_{1}, \ldots, n_{k-1}$, let $n_{k}$ be the smallest natural number greater than $n_{k-1}$ such that $x_{n_{k}} \in E$.   Then set $f(k) = x_{n_{k}} \in E$.  This function is a bijection by construction and therefore $E$ is countable.

\end{proof}

This theorem shows that roughly speaking, no uncountable set can be a subset of a countable set.

\begin{theorem}  The union of countably many countable sets is countable.   Let $\{E_{n}\} ,  n=1,2,3,\ldots$ be a sequence of countable sets.  Then the set $S$,
\[S = \bigcup_{n=1}^{\infty} E_{n} \] is countable. \end{theorem}

\begin{proof}
Let every set $E_{n}$ be arranged in a sequence $\{x_{n_k}\}, k = 1,2,3, \ldots$ and consider the infinite array in which the elements of $E_n$ form the nth row of the array.  The array contains all of the elements of $S$.  As in the Cantor construction, these elements can be arranged in a sequence (denoted by the arrows) of the form  
\begin{marginfigure}
  \includegraphics{rationals-countable.jpg}
  \caption{An example of the Cantor Construction for countability proofs.}
 % \label{fig:textfig}
\end{marginfigure}

\[x_{11}; x_{21}, x_{12}; x_{31}, x_{22}, x_{13}; x_{41}, x_{32}, x_{23}, x_{14}; \ldots\]

If any two sets have an element in common these will appear twice in the about sequence.  Therefore there is a subset of the natural numbers, call it T, such that $S \sim T$ which shows that S is at most countable.

However, since $E_{1}\subset S$ and $E_{1}$ is infinite, then S is infinite and thus countable.
\end{proof}

\begin{corollary} Suppose $A$ is at most countable and that for every $\alpha \in A$, there is associated a set $B_{\alpha}$ that is at most countable.  Then the set T
\[T = \bigcup_{\alpha \in A} B_{\alpha}\]
is at most countable. \end{corollary}



\begin{theorem}  A \emph{finite} product of countable sets is countable.  For any n,  $E_{1} \times E_{2} \times \ldots \times E_{n}$ is countable for $E_{i}$ all countable. \end{theorem}

\marginnote{Note that the fact that the product is \emph{finite} is key.  \\Take for example, \\$E = \{0,1\},\quad E \times E \times E \times E \ldots$ \\which is not countable.  (Exercise!)}


\begin{proof}
We complete this proof by induction on n.  For $n=1$, the assertion holds trivially.  Suppose that $E_{1} \times E_{2} \times \ldots \times E_{n-1}$ is countable.  To show this for $E_{1} \times E_{2} \times \ldots E_{n-1}\times E_{n}$, we notice that since $E_{n}$ is countable, we can list its elements as a sequence;  $E_{n}= {y_1, y_2, \ldots}$.  Then

\[E_{1} \times E_{2} \times \ldots E_{n-1}\times E_{n} = \]
\[  \big( E_{1} \times \ldots \times E_{n-1} \times \{y_1\} \big) \cup  \big( E_{1} \times  \ldots \times E_{n-1} \times \{y_2\}\big) \cup  \ldots   \] 

\center{And}
\[E_{1} \times E_{2} \times \ldots E_{n-1}\times E_{n} = \bigcup_{i=1}^{\infty} (E_{1} \times  \ldots \times E_{n-1}) \times \{y_i\}  \] 

Therefore,  $E_{1} \times E_{2} \times \ldots E_{n-1}\times E_{n} $ is equivalent to the countable union of countable sets and, by the theorem above, is also countable.
\end{proof}


\begin{corollary}The set of all rational numbers is countable. \end{corollary}

By the Theorem above, we take $n=2$ noting that every rational is of the form $b/a$ where $a$ and $b$ are both integers.  The set of pairs $(a,b)$ and therefore the set of fractions $b/a$ is countable.  
\medskip

In fact, even the set of all algebraic numbers is countable.

\begin{corollary} 
A complex number is said to be \emph{algebraic} if there are integers $a_{0}, \ldots, a_{n}$ not all zero such that 
$$ a_{0}z^n + a_{1}z^{n-1} + \cdots + a_{n-1}z + a_{n} = 0. $$

Prove that the set of algebraic numbers is countable.
\end{corollary}
 
\begin{proof} 
For each fixed integer $n \geq 0$, let us call the set $P_{n}$ the set of all polynomials with integer coefficients and degree of at most $n$.  We can be sure that such a set is countable since it is equivalent to set of all $n+1$-tuples $(a_{0}, \ldots, a_{n})$, where $a_{i} \in Z$ which is countable, as stated in Theorem 2.13.  

The set of ALL polynomials of every degree with integer coefficients can therefore be expressed as
$$   P = \bigcup^{\infty}_{n=0} P_{n} \qquad \text{where} \ n \in N.$$


By Theorem 2.12, the countable union of countable sets is also countable and therefore P, the set of all polynomials with integer coefficients, is countable.

By the Fundamental Theorem of Algebra, each polynomial of degree $n$ has only finitely many roots (at most $n$). Therefore the set of all possible roots of all polynomials with integer coefficients is a countable union of finite sets, hence countable.
\end{proof}


But note that not all infinite sets are countable as shown in the next theorem.

\begin{theorem} Let A be the set of sequences whose elements are 0's and 1's.  This set A is uncountable. \end{theorem}

\begin{proof}
Let $E$ be a countable subset of A and let E consists of the sequences $s_1, s_2, s_3, \ldots$.   We construct a sequence $s$ as follows: If the nth digit in $s_n$ is 1, we let the nth digit of $s$ be 0, and vice versa.  Then the sequence s differs from every member of E in at least one place and hence, $s \notin E$.  Clearly $s \in A$ however so that E is a proper subset of A.\\

We have shown that every countable subset of A is a proper subset of A.  It follow that A is uncountable (otherwise A would be a proper subset of A which is absurd).
\end{proof}

\begin{proof} \marginnote{Alternative proof presented by Andre}

We proceed to prove by contradiction.  Say A is countable and we can list its elements in an array:  $A = \{x_1, x_2, \ldots \}$.  Let $y = (y_1, y_2, y_3, \ldots)$ be the sequence of 0's and 1's listed along the diagonal in the array.   Define the sequence $z \in A$ as the sequence that exchanges all of the 1's in y for 0's and vice versa.  But then we have that z differs from every $x_i \in A$ in at least one place and hence $z$ cannot be in $A$.  Contradiction.

\end{proof}

\begin{theorem} The open interval $(0,1) \subset \mathbb{R}$ is uncountable.  \marginnote{This theorem is equivalent to the one below and to the theorem which states that R is uncountable.  (0,1) is uncountable iff [0,1] is uncountable iff R is uncountable.  See why this is true.}  \end{theorem}


\begin{proof}
We proceed to prove by contradiction.  Assume that there exists a function $f: \mathbb{N} \to (0,1)$ that is bijective.  Then for each $m \in \mathbb{N}$, $f(m)$ is a real number between 0 and 1, and we represent it using decimal notation

\[f(m) = .a_{m1}a_{m2}a_{m3}a_{m4}a_{m5} \ldots\]


\marginnote{This theorem can also be proved by the production of a bijection from $R \to (0,1)$.  Take for example $$ f(x) = \frac{(x - 1/2)}{(x-x^2)} $$  Since $R$ is uncountable, it must follow that $(0,1)$ is also uncountable.}
What is meant here is that for each $m,n \in \mathbb{N}, a_{mn}$ is the digit from the set $\{0,1,2, \ldots, 9\}$ that represents the nth digit in the decimal expansion of $f(m)$.  Therefore we can index the array of $f(m)'s$ by the digits of their decimal expansion.   We assume that \emph{every} real number between 0 and 1 appears somewhere on this list.\\

Now for the pearl of the argument.   Define a real number $x \in (0,1)$  with the decimal expansion $x=.b_1 b_2 b_3 b_4 \ldots$ using an arbitrary rule such as 

 \begin{displaymath}
   f(x) = \left\{
     \begin{array}{lr}
       2 & : a_{nn} \neq 2 \\
       3 & : a_{nn} = 2
     \end{array}
   \right.
 \end{displaymath} 
   
   
   Since, by construction, this real number x differs from every other number listed, x cannot be a member of $A$ as defined.  Therefore, we have a contradiction.
   
\end{proof}

\begin{theorem} The closed interval $[0,1] \subset \mathbb{R}$ is uncountable.  \end{theorem}
\marginnote{Note that this is true iff R is uncountable.  Also note that this theorem is equivalent to the one below.  From this we can conclude that R is uncountable.}  (Hint: express $x \in [0,1]$ in binary notation)   From this conclusion, we can determine that $\mathbb{R}$ is uncountable.  


\begin{theorem}  (i) The set $\mathbb{Q}$ is countable\\ (ii) The set $\mathbb{R}$ is uncountable.  \marginnote{There are 2 proofs for the countability of $Q$.  \begin{itemize}
\item Function using the Cantor Construction  \item $Q = N \times N$ and is a finite product of countable sets. Therefore countable. \end{itemize}}  \marginnote{There are 4 proofs for the uncountability of $R$:
\begin{itemize}
	\item NIP
	\item Using the binary representations of real numbers and the fact that the sequences of $0s$ and $1s$ is uncountable
	\item There are infinite subsets of $R$ (namely $(0,1)$ and $[1,1]$ ) that are uncountable.
	\item Max Lipyanskiy's Senior Thesis at Columbia
\end{itemize}}
\end{theorem}

\begin{proof}
(i) For each $n \in \mathbb{N}$, let $A_n$ ne the set given by:
\[A_n = \{\pm \frac{p}{q}: \text{where} \quad p,q \in \mathbb{N} \}\]
Also note that we stipulate that p and q are in the lowest possible terms with $p+q=n$.  \\
What's crucial to observe is that every $A_n$ is \emph{finite} and every rational number appears in exactly one of these sets.  Our 1-1 correspondence with the natural numbers is then achieved by listing each of the elements of Q.  \\
(ii) We proceed with a proof by contradiction.  Assume that there exists a bijection $f: \mathbb{N} \to \mathbb{R}$. This suggests that is it possible to enumerate the elements of $\mathbb{R}$.  We can write
\begin{equation}\mathbb{R} = \{x_1, x_2, x_3, x_4, \ldots\}\end{equation}
and be confident that every real number appears somewhere on the list.  We now use something called the Nested Interval Property\marginnote{Nested Interval Property: For each $n \in \mathbb{N}$, assume we are given a closed interval $[a_n, b_n] = \{x \in \mathbb{R}: a_n \leq x \leq b_n\}$.  Assume also that for each $I_n \supset I_{n+1}$.  Then the resulting nested sequence of closed intervals has a nonempty intersection. } to produce a real number that is not there on our list.\\  Let $I_1$ be a closed interval that \emph{does not} contain $x_1$.  Next, let $I_2$ be a closed interval, contained in $I_1$ such that $x_2 \notin I_2$.  The existence of $I_2$ is easy to verify.  Certainly, $I_1$ contains two smaller \emph{disjoint} closed intervals and $x_2$ can only be in one of these.  In general, given an interval $I_n$, construct $I_n+1$ to satisfy \begin{enumerate}
	\item $I_n+1 \subset I_n$ and
	\item $x_{n+1} \notin I_{n+1}$
\end{enumerate}


We now consider the intersection \[ \bigcap_{n=1}^{\infty}I_n \]  

\begin{marginfigure}
  \includegraphics{NIP_RealsUncount.pdf}
 % \label{fig:textfig}
\end{marginfigure}

If $x_{n_0}$ is some real number from the list (1), then we have that $x_{n_0} \notin I_{n_0}$ and therefore not in the intersection of the intervals.\\  Now assuming that (1) contains every real number, this leads to the conclusion that 
\[ \bigcap_{n=1}^{\infty}I_n = \emptyset \]

However, according to the Nested Interval Property this is impossible.  By NIP there is a least one x in the intersection that consequently cannot be on the list in (1).  This contradiction means that such an enumeration of $\mathbb{R}$ is impossible and we can conclude that $\mathbb{R}$ is uncountable.



\end{proof}

\section{Properties of Norms}


Suppose that $\mathbf{x,y,z} \in R^k$ and $\alpha$ is real. Then
\begin{itemize}[$\diamond$]
	\item $|\mathbf{x}| \geq 0$
	\item $|\mathbf{x}| = 0$ iff $\mathbf{x} = 0$
	\item $|\alpha \mathbf{x}| = |\alpha||\mathbf{x}|$
	\item $|\mathbf{x} \cdot \mathbf{y}| \leq |\mathbf{x}| |\mathbf{y}| $
	\marginnote{The Cauchy-Schwarz Inequality}
	\item $|\mathbf{x+y}| \leq |\mathbf{x}|  + |\mathbf{y}|$
	\marginnote{The Triangle Inequality}
	\item $|\mathbf{x}- \mathbf{z}| \leq |\mathbf{x-y}| + |\mathbf{y-z}|$     
\end{itemize}




\bookchapter{Two}{Sequences and Series}

\section{Sequences}

Most of the concepts in analysis can be reduced to statements about the behavior of sequences.  Therefore it is important to have a clear understanding of them before moving forward.




\begin{definition}\textsc{Sequences} are functions whose domain is $\mathbb{N}$.  \end{definition}

The formal definition leads immediately to the familiar depiction of a sequence as an ordered list of real numbers.  Given a function $f: \mathbb{N} \to \mathbb{R}$,   $f(n)$ is just the $n$th term on the list.  The notations for sequences reinforce this understanding.

What is crucial to understand is that a sequence is an \emph{infinite} list of real numbers.  What happens at the beginning is of little importance since the business of analysis is concerned with the behavior of the infinite tail of a given sequence.  

\begin{definition}Convergence of a Sequence:  A sequence $(p_n)$ converges to a real number $p$ if, \marginnote{It might be well to point out that our definition of convergent sequence depends not only on the sequence in question but also on the metric space $X$.}for every positive number $\varepsilon$, there exists an $N \in \mathbb{N}$ such that whenever $n \geq N$ it follows that $|p_n - p|< \varepsilon$. \end{definition}


For Rudin, a sequence $\{p_n\}$ in a metric space $X$ is said to converge if there is a point $p \in X$ with the following property:  For every $\varepsilon > 0$ there is an integer $N$ such that $n \geq N$ implies that $d(p_n, p) < \varepsilon$.

In this case, we also say that  $\{p_n\}$ converges to $p$, or that $p$ is the limit of  $\{p_n\}$ and we write that 
\[\lim_{n \to \infty} p_n = p   \qquad (p_n) \to p\]


If a sequence does not converge, it is said to diverge.  \marginnote{Note: the definition of divergence depends on the negation of quantifiers.  A sequence diverges if, for \emph{every} candidate limit $p$, there exists an $\varepsilon>0$ for which no $N$ works: however large $N$ is, some $n \geq N$ has $d(p_n, p) \geq \varepsilon$.}

\begin{definition}
We recall that the set of all points $p_n  (n=1,2,3, \ldots)$  is the \textsc{range} of the sequence.  The range can be infinite or finite.  


 The sequence $\{p_n\}$ is said to be \textsc{bounded} if its range is bounded.  That is, a sequence is bounded if there exists an $M > 0$ such that $|p_n|\leq M$ for all $n \in \mathbb{N}$.\end{definition}


\bigskip
In an effort to decipher this complicated definition, it helps to consider the phrase ``$d(p_n, p) < \varepsilon$'' and think about the points that satisfy an inequality of this type.


\begin{definition}$\varepsilon$-neighborhoods:  Given a point $p \in X$ and a positive number $\varepsilon > 0$,  the set
\[V_{\varepsilon}(p) = \{x \in X: d(x,p)< \varepsilon\}\]
is called the $\varepsilon$-neighborhood of $p$. \end{definition}

Notice that $V_{\varepsilon}(p)$ consists of all those points whose distance from p is less than $\varepsilon$.  Recasting the definition of convergence in terms of $\varepsilon$-neighborhoods gives us a more geometric impression of what is being described.

\begin{definition}\newthought{Convergence of a Sequence (Topological Version):}  a sequence $\{p_n\}$ in a metric space $X$ converges to $p$, if, given any $\varepsilon$-neighborhood  $V_{\varepsilon}(p)$ of $p$, there exists a point in the sequence after which all the terms are in $V_{\varepsilon}(p)$.  In other words, every $\varepsilon$-neighborhood centered at $p$ contains all but a finite number of the terms of $\{p_n\}$.
\end{definition}

Notice that these two definitions of convergence say exactly the same thing.  The natural number $N$ from before is the point where the sequence enters $V_{\varepsilon}(p)$ and never leaves.  It should be apparent that \emph{the value of $N$ depends on the choice of} $\varepsilon$.   

The smaller the $\varepsilon$-neighborhood, the larger $N$ may have to be.


\noindent \emph{Example:}  We claim that $$ \lim \bigg( \frac{1}{\sqrt{n}} \bigg) = 0 $$

\begin{proof}
Let $\varepsilon >0$ be an arbitrary positive number.  Choose a natural number $N$ such that 

$$	\frac{1}{\sqrt{N}} < \varepsilon, \quad \text{that is,} \quad N > \frac{1}{{\varepsilon}^2}.	$$

We now verify that this choice of $N$ has the desired property.  Let $n \geq N$.  Then,

$$	n> \frac{1}{\varepsilon^2}	  \quad \text{implies} \quad \frac{1}{\sqrt{n}} < \varepsilon \quad \text{and hence} \quad |a_{n} -0|< \varepsilon$$

\end{proof}

\section{Some Examples:}\label{sec:sequence-examples}

\begin{enumerate}
	\item If $s_n = 1/n$ then $\lim_{n \to \infty} s_n = 0$.  The range is infinite and the sequence is bounded.
	\item If $s_n = n^2$ then the sequence is divergent, unbounded and has infinite range.
	\item If $s_n = 1 + [{(-1)^n}/n]$ then the sequence $\{s_n\}$ converges to 1, is bounded, and has infinite range.
	\item If $s_n = i^n$, then the sequence is divergent, is bounded, and has finite range.
	\item If $s_n = 1$ then the sequence converges to 1, is bounded, and has finite range.
	\item  Now take the sequence $a_n = 1/\sqrt{n}$.
	
	Our intuitive understanding tells us that the sequence converges to 0 but let's first explore the relationship between $\varepsilon$ and $N$ in the definition of convergence.
	
	Let's define a sort of ``target zone''  and say that we take $\varepsilon = 1/10$.  By claiming that the limit of the sequence is 0, we are saying that the terms of the sequence eventually get arbitrarily close to 0.  
	
	We have set $\varepsilon = 1/10$ as our standard of ``closeness,'' which leads to the neighborhood $(-1/10, 1/10)$ around the limit 0.  How far out into the sequence before the terms fall into this interval?  The 100th term $a_{100} = 1/10$ puts us right on the boundary so we can say that,  if $n>100$ then $a_n \in (-1/10, 1/10)$.  Thus for $\varepsilon = 1/10$ we choose $N=101$ (or anything larger) as our response.
	
	But our choice of $\varepsilon $ was maybe rather whimsical.  How can we be sure that for every choice of $\varepsilon$ we can find an appropriate $N$?  For this problem, the desired relationship is implicit in the algebra carried out to compute our previous answer.  Whatever $\varepsilon $ happens to be, we want 
	\[\frac{1}{\sqrt{n}} < \varepsilon\]
	which is equivalent to 
	\[n > \frac{1}{\varepsilon ^ 2}\]
	
	Now we are ready to write a formal argument.

We claim that \[\lim_{n \to \infty} \frac{1}{\sqrt{n}} = 0\]

\begin{proof}
Let $\varepsilon> 0 $ be an arbitrary positive number.  Choose a natural number $N$ satisfying $N > \frac{1}{\varepsilon ^ 2}$.  

We now verify that this choice of $N$ has the desired property.  Let $n \geq N$.  Then,
\[n > \frac{1}{\varepsilon^2} \implies \frac{1}{\sqrt{n}}< \varepsilon\]

Then \[\left|\frac{1}{\sqrt{n}} - 0\right| < \varepsilon. \]  \end{proof}
\end{enumerate}



The  definition of convergence given earlier is the result of hundreds of years of refining the intuitive notion of a limit into a mathematically rigorous statement.  The logic of such a statement is intimately tied to the use of the quantifiers ``for all'' and ``there exists.''  The ability to write proofs of convergence depends heavily on a firm understanding of how to manipulate such statements.

\newthought{Template for Proof that $(x_{n}) \to x$}
	\begin{enumerate}
			\item Let $\varepsilon>0$ be arbitrary
			\item  Demonstrate a choice for $N \in \mathbb{N}$.  This usually requires the most work, which should be done before actually writing the proof (less messy)
			\item Now, show that $N$ actually works
				\begin{enumerate}
						\item Assume $n \geq N$
						\item With $N$ well chosen, it should be possible to derive the inequality that $|x_{n} - x| < \varepsilon$.
				\end{enumerate}
	\end{enumerate}


\noindent \emph{Example:}  We claim that $$ \lim \bigg( \frac{n+1}{n} \bigg) = 1 $$

\begin{proof}
Let $\varepsilon >0$ be an arbitrary positive number.  Choose a natural number $N$ such that $ N > \frac{1}{\varepsilon}	$.


We now verify that this choice of $N$ is appropriate.  Let $n \geq N$.  Then,

$$	n> \frac{1}{\varepsilon}	  \quad \text{implies} \quad \frac{1}{n} < \varepsilon \quad \text{and hence} \quad \bigg|  \frac{n+1}{n}  -1  \bigg|  = \bigg|  \frac{1}{n} \bigg|< \varepsilon$$

as desired.
\end{proof}



Significant insight into the role of the quantifiers in the definition of convergence
can be gained by studying an example of a sequence that does not have a limit.

\begin{definition}Divergence:  A sequence that does not converge is said to diverge. \end{definition}


To prove that a particular number $x$ is not the limit of a sequence $(x_{n})$ we must produce a single value of $\varepsilon$ for which no $N \in \mathbb{N}$ works.
\medskip

\noindent
\emph{Exercise 2.2.1.} Verify, using the definition of convergence of a sequence, that
the following sequences converge to the proposed limit.
\begin{enumerate}[(a)]
\item $\lim \frac{1}{(6n^{2}+1)} = 0$.
\item $\lim \frac{(3n+1)}{(2n+5)} = \frac{3}{2}$.
\item $\lim \frac{ 2}{\sqrt{n+3}} = 0$.
\end{enumerate}

\begin{proof}
	\begin{enumerate}[(a)]
			\item Let $\varepsilon > 0$ be arbitrary. We must show that there exists an
$N \in \mathbb{N}$ such that $n\geq N$ implies $\big|\frac{1}{(6n^{2}+1)} - 0 \big|  < \varepsilon$.  Well, manipulation of the statement tells us that 

$$\bigg|\frac{1}{(6n^{2}+1)} \bigg|  = \frac{1}{(6n^{2}+1)}  \quad \text{and therefore that}$$

$$ \frac{1}{(6N^{2}+1)}  < \varepsilon \quad \text{which holds when} \quad N > \sqrt{\frac{1- \varepsilon}{6\varepsilon}}$$

(for $\varepsilon < 1$; any $N$ works when $\varepsilon \geq 1$).  However, we can avoid the square root entirely.  Since $6n^2 + 1 > 6n$, we have $\frac{1}{6n^{2}+1} < \frac{1}{6n}$, so it is enough to choose $N > \frac{1}{6\varepsilon}$: then $n \geq N$ gives $\frac{1}{6n^2+1} < \frac{1}{6n} \leq \frac{1}{6N} < \varepsilon$.  

It follows that $n \geq N$  implies that $\frac{1}{(6n^{2}+1)}  < \varepsilon$ as desired.


\item  Let $\varepsilon>0$.  Produce $N$ such that $n\geq N$ implies
$$   \bigg|\frac{(3n+1)}{(2n+5)} -\frac{3}{2}\bigg| <  \varepsilon $$

A little algebra shows that the left-hand side equals $\frac{13}{4n+10}$, and that $n>\frac{13-10\varepsilon}{4\varepsilon}$ makes this smaller than $\varepsilon$.  So we pick $N >\frac{13-10\varepsilon}{4\varepsilon}$; then for $n \geq N$ it follows that $\frac{13}{4n+10} < \varepsilon$.


  \item Let $\varepsilon>0$.  We must produce an $N$ so that $n \geq N$ implies that $|\frac{2}{\sqrt{n+3}} - 0|< \varepsilon$.  Well,
  $$ \bigg| \frac{2}{\sqrt{n+3}} - 0 \bigg| = \frac{2}{\sqrt{n+3}} 	$$
  
  Pick $N$ such that $N > 4/{\varepsilon^2} - 3$.  It then follows that when $n \geq N$ we get $\frac{2}{\sqrt{n+3}} <\varepsilon$ as desired.
	\end{enumerate}
\end{proof}


\noindent
\emph{Exercise 2.2.4.} Argue that the sequence $$1,0,1,0,0,1,0,0,0,1,0,0,0,0,1,(\text{5 zeros}),1,...$$
does not converge to zero. For what values of $\varepsilon > 0$ does there exist a response $N$. For which values of $\varepsilon > 0$ is there no suitable response?

\begin{proof}
For any $\varepsilon$ that is greater than $1$, there exists a response $N$. In this case, $N$ can be any natural number.

For any $\varepsilon$ that is less than or equal to $1$, there exists no suitable response. This is because, although the 1s in the sequence occur less and less frequently as we go out the sequence, there is still no point in the sequence where the sequence enters the neighborhood $(-\varepsilon, \varepsilon)$ and never leaves.

\end{proof}


\emph{Exercise 2.2.6.} Suppose that for a particular $\varepsilon > 0$ we have found a suitable value of $N$ that ``works'' for a given sequence in the sense of Definition 2.2.
\begin{enumerate}[(a)]	
\item Then, any larger/smaller (pick one) $N$ will also work for the same $\varepsilon > 0$. 
\item Then, this same $N$ will also work for any larger/smaller value of $\varepsilon$.
\end{enumerate}

\begin{proof}
(a) Any larger $N$ will also work for the same $\varepsilon > 0$. (b) This same $N$ will also work for any larger value of $\varepsilon$.
\end{proof}

We now summarize some important properties of convergent sequences in metric spaces.

\begin{theorem}  Let $\{p_n\}$ be a sequence in a metric space $X$.
\begin{enumerate}[(a)]
	\item $\{p_n\}$  converges to $p$ iff every neighborhood of $p$ contains all but finitely many of the terms of $\{p_n\}$
	\item If $p \in X, p' \in X$ and if $\{p_n\}$ converges to $p$ and to $p'$, then $ p=p'$
	\item If $\{p_n\}$ converges, then $\{p_n\}$ is bounded.
	\item If $E \subset X$ and if $p$ is a limit point of $E$, then there is a sequence $\{p_n\}$ in $E$ such that $\{p_n\}$ converges to $p$.
\end{enumerate}
\end{theorem}

\begin{proof}
(a) Suppose that $p_n \to p$ and let $V$ be a neighborhood of $p$.  For some $\varepsilon > 0 $, the conditions $d(p,q)< \varepsilon, \quad q \in X$ imply $q \in V$.  Corresponding to this $\varepsilon$ there exists $N$ such that $n \geq N$ implies $d(p_n, p)< \varepsilon$.  Thus $n \geq N$ implies $p_n \in V$.

Conversely, suppose that every neighborhood of $p$ contains all but finitely many elements of $\{p_n\}$.  Fix $\varepsilon>0$ and let $V$ be the set of all $q \in X$ such that $d(p,q)< \varepsilon$.  By assumption, there exists an $N$ (corresponding to this $V$) such that $p_n \in V$ if $n \geq N$.  Thus $d(p_n, p) < \varepsilon$ if $n \geq N$; hence $p_n \to p$.

(b)  Let $\varepsilon > 0$ be given.  There exist integers $N, N'$ such that
\[n\geq N \implies d(p_n, p) < \frac{\varepsilon}{2}\]
\[n \geq N' \implies d(p_n, p') < \frac{\varepsilon}{2}\]

If we choose $n\geq \max(N, N')$, we have that 
\[d(p,p') \leq d(p_n,p) + d(p_n, p') < \varepsilon\]
Since $\varepsilon$ was an arbitrary choice, we can therefore conclude that $d(p',p) =0$ and that $p=p'$.

(c)  Suppose that $p_n \to p$.  Then there exists an integer $N$ such that $n\geq N$ implies that $d(p_n, p) < 1$.  Now suppose that we set $r$ to
\[r = \max \{1, d(p_1, p), \ldots, d(p_N, p)\}\]

Then $d(p_n, p) \leq r$ for all $n =1, 2, 3, \ldots$.  Therefore the sequence is bounded.

(d) Let $p$ be a limit point of $E \subset X$. \marginnote{Limit point:  A point $p \in X$ is a limit point of the set $E$ if every neighborhood $V_{\varepsilon}(p)$ contains a point $q \neq p$ such that $q \in E$} Then \emph{every} neighborhood of $p$ contains points $p_n\in E$ with $p_n\neq p$. In other words, every open ball around $p$, contains other points of $E$.  

We now build a sequence $\{p_n\}$ in $E$.  We define it as follows:  for each positive integer $n$,  let $p_n$ be a point of $E$, different from $p$, such that $d(p_n, p) < 1/n$.  

Given $\varepsilon > 0$, choose $N$ such that $N > \frac{1}{\varepsilon}$ (which is equivalent to $N\varepsilon>1$).  If $n > N$, it follows that $d(p_n, p)< \frac{1}{n} < \frac{1}{N} < \varepsilon$.  Hence $p_n \to p$.
\end{proof}


\section{The Algebraic and Order Limit Theorems}

The point of having such a logically tight description of convergence is so that we can confidently state and prove statements about convergent sequences in general.  We are ultimately trying to resolve arguments about what is and is not true regarding the behavior of limits with respect to the mathematical manipulations we intend to inflict on them.

\begin{definition}\newthought{Bounded:} A sequence $(x_{n})$ is bounded if there exists a number $M>0$ such that $|x_{n}|\leq M$ for all $n \in \mathbb{N}$. \end{definition}

Geometrically this implies that we can find an interval on the real number line $[-M, M]$ that contains every term of the sequence $(x_{n})$.

\begin{theorem} Every convergent sequence is bounded. \end{theorem}

\begin{proof}
Assume $(x_{n})$ converges to a limit $l$. This means that given a particular value of $\varepsilon$, we know that there exists an $N \in \mathbb{N}$ such that if $n\geq N$, then $x_{n}$ is in the interval $(l-\varepsilon, l+\varepsilon)$.  Without knowing whether or not $l$ is positive or negative, we can nonetheless be sure that \marginnote{ Note that $|x_n| = |x_n + l - l| \leq |x_n - l| + |l| < \varepsilon + |l|$}  $$ |x_{n}| < |l| + \varepsilon $$ for all $n\geq N$.

We still need to worry (just a little) about the terms that come before the Nth term.  Because there are only a finite number of these terms, we can let $$ M = \max \{|x_{1}|,|x_{2}|, \ldots, |x_{N-1}|,|l|+\varepsilon \}$$
It follows that $|x_{n}|\leq M$ for all $n \in \mathbb{N}$ as desired.
\end{proof}

The following theorem  illustrates that sequences behave very well with respect to each other.  This should be easily done using the convergence definition and the properties of norms from earlier.

\begin{theorem}[ Algebraic Limit Theorem] Suppose that $\{s_n\}$ and $\{t_n\}$ are complex sequences, and that they converge to $s$ and $t$ respectively.  Then:
\begin{enumerate}[(a)]
	\item $\{s_n + t_n\} \to s+t$
	\item $\{cs_n\} \to cs$ and $\{c+s_n\} \to c+s$, for any number $c$
	\item $\{s_n t_n\} \to st$
	\item $\{1/s_n\} \to 1/s$, provided that $s_n \neq 0$ for all $n$ and $s \neq 0$
\end{enumerate}
\end{theorem}

\begin{proof}
(a)  Given $\varepsilon > 0$, there exist integers $N_1, N_2$ such that 
\[n \geq N_1 \implies |s_n - s|< \frac{\varepsilon}{2}\]
\[n \geq N_2 \implies |t_n - t|< \frac{\varepsilon}{2}\]

If $N = \max(N_1, N_2)$ then $n\geq N$ implies
\[|(s_n + t_n)-(s + t)|\leq |s_n - s|+|t_n - t|<\varepsilon\] proving that $\{s_n + t_n\} \to s+t$.

(b)  If $c = 0$ there is nothing to prove.  Otherwise, given $\varepsilon>0$, there exists $N$ such that $n\geq N \implies |s_n-s|<\varepsilon/ |c|$.  That is, whenever $n\geq N$

\[|cs_n -cs |= |c||s_n - s|< \varepsilon\]  proving that $\{cs_n\} \to cs$.

Similarly, since $|(c+s_n) - (c+s)| = |s_n - s|$, any $N$ that works for $\{s_n\}$ also shows that $\{c+s_n\} \to c+s$.

(c) Given $\varepsilon > 0$, there exist integers $N_1, N_2$ such that 
\[n \geq N_1 \implies |s_n - s|< \sqrt{\varepsilon}\]
\[n \geq N_2 \implies |t_n - t|< \sqrt{\varepsilon}\]

If we take $N = \max(N_1, N_2), n\geq N$ implies that 
\[|(s_n -s )(t_n -t)| \leq |s_n - s| |t_n -t| < \varepsilon\]

which is to say that the sequence $\{(s_n - s)(t_n - t)\}$ converges to 0.

To show that $\{s_n t_n \} \to st$ we write 


$\begin{array} {lcl} 
|s_n t_n  - st|& = &|(s_n - s) (t_n - t) + s(t_n -t) + t(s_n - s)|\\ 


& \leq & |(s_n - s) (t_n - t)| + |s||(t_n -t) |+ |t||(s_n - s)|    \end{array}$

Therefore for $n  \geq N$ as before, we have that 

\[|(s_n - s) (t_n - t)| + |s||(t_n -t) |+ |t||(s_n - s)| < \varepsilon + |s| \sqrt{\varepsilon} + |t| \sqrt{\varepsilon}.\]
Since $\varepsilon$ was arbitrary, the right-hand side can be made as small as we please, and so $\{s_n t_n\} \to st$.



\emph{Alternative Proof:}  Begin by observing that 


$\begin{array} {lcl} 
|s_n t_n  - st|& = &|s_n t_n - st_n + st_n  -st|\\ 
& \leq & |s_n t_n - st_n| + |st_n  -st|\\
&=	& |t_n||s_n - s| + |s||t_n - t| 
    \end{array}$
    
Letting $\varepsilon$ be arbitrary, we proceed with the strategy of making each piece less than $\varepsilon/2$.  Using the fact that convergent sequences are bounded (proved previously) we know that there exists a bound $M>0$ such that $|t_n|\leq M$ for all $n \in \mathbb{N}$.  Therefore for any $\varepsilon >0$ there exists $N_1, N_2$ such that 
\[n\geq N_1 \implies |s_n - s| <  \frac{1}{M}\frac{\varepsilon}{2}\]
\[n\geq N_2 \implies |t_n - t| <  \frac{1}{|s|}\frac{\varepsilon}{2}\]

(If $s = 0$ the second term vanishes and $N_2$ is not needed.)  Pick $N = \max(N_1, N_2)$ and the calculation shows that 

\smallskip

$\begin{array} {lcl} 
|s_n t_n  - st|& \leq &|t_n||s_n - s| + |s||t_n - t| \\ 
& \leq & M|s_n - s| + |s||t_n - t|\\
& <	&  M\frac{1}{M}\frac{\varepsilon}{2} + |s|\frac{1}{|s|}\frac{\varepsilon}{2}\\
&=& \varepsilon
    \end{array}$

\smallskip



(d) To show that $\{1/s_n\} \to 1/s$ we first observe that 
\[\left\lvert \frac{1}{s_n} - \frac{1}{s} \right\rvert = \frac{|s-s_n|} {|s||s_n|}.\]

Since $\{s_n\} \to s$ we can make the preceding numerator as small as we like by choosing $n$ large.  The problem is that we also need an estimate on the size of $\frac{1} {|s||s_n|}$. The trick is to look far enough out into the sequence $\{s_n\}$ that the terms are closer to $s$ than they are to 0.

Consider the particular value $\varepsilon_0 = \frac{|s|}{2}$.  Because $\{s_n\} \to s$, there exists $N_1$ such that 
\[n \geq N_1 \implies |s_n - s|<\frac{|s|}{2}, \quad \text{and therefore} \quad |s_n| > \frac{|s|}{2}.\]
Next, we choose $N_2$ such that 
\[n \geq N_2 \implies |s_n -s| < \frac{\varepsilon |s|^2}{2}.\]

Let $N=\max\{N_1, N_2\}$, then $n \geq N$ implies that
\[\left\lvert \frac{1}{s_n} - \frac{1}{s} \right\rvert = |s-s_n|\frac{1} {|s||s_n|} < \frac{\varepsilon |s|^2}{2} \frac{1}{|s|\frac{|s|}{2}} = \varepsilon.\]



\end{proof}

\begin{theorem} Extending our previous results to sequences in $\mathbb{R}^k$

(a) Suppose $x_n \in \mathbb{R}^k$ $(n=1,2,3,\ldots)$ and
\[x_n = (\alpha_{1,n}, \ldots, \alpha_{k,n})\]
Then $\{x_n\}$ converges to $x = (\alpha_1, \ldots \alpha_k)$ iff 
\[\lim_{n\to\infty}\alpha_{j,n} = \alpha_j \quad (1\leq j \leq k)\]

(b) Suppose $\{x_n\}, \{y_n\}$ are sequences in $\mathbb{R}^k$, $\{\beta_n\}$ is a sequence of real numbers, and $x_n \to x$, $y_n \to y$, $\beta_n \to \beta$. Then
\[\lim_{n\to\infty} (x_n + y_n ) = x+y, \quad \lim_{n\to\infty} (x_n \cdot y_n ) = x\cdot y,  \quad \lim_{n\to\infty} \beta_n x_n  = \beta x\]
\end{theorem}

\begin{proof}

\end{proof}

The above theorem verifies that the relationship between algebraic combinations of sequences and the limiting process is as trouble-free as we could hope for.  Limits can be computed from the individual component sequences provided that each component limit exists.  The limiting process is also well-behaved with respect to the order operation.

\begin{theorem}[Order Limit Theorem]  Assume $(a_{n}) \to a$ and $(b_{n})\to b$
\begin{enumerate}[(i)]
	\item If $a_{n} \geq 0$ for all $n \in \mathbb{N}$, then $a\geq0$.
	\item If $a_{n}\leq b_{n}$ for all $n \in \mathbb{N}$, then $a\leq b$.
	\item If there exists $c \in \mathbb{R}$ for which $c \leq b_{n}$ for all $n \in \mathbb{N}$, then $c \leq b$.  Similarly, if $a_{n} \leq c$ for all $n \in \mathbb{N}$, then $a\leq c$.
\end{enumerate}
\end{theorem}

\begin{proof}
\begin{enumerate}[(i)]
	\item We will begin this proof by contradiction; that is, assume that $a<0$.  The idea is to produce a term in the sequence $(a_{n})$ that is also less than zero.  To do this, we consider a particular value $\varepsilon_{0}=|a|$. The definition of convergence guarantees that we can find an $N$ such that $|a_{n} - a|<|a|$ for all $n \geq N$.  In particular this would mean that $|a_{N} - a| < |a|$ which implies that $a_{N}<0$.  This contradicts our hypothesis that $a_{N} \geq 0$.  We therefore conclude that $a \geq 0$.
	
	\begin{marginfigure}
  \includegraphics{OrderLimitTH.jpg}
  \caption{If $a_{n} \geq 0$ for all $n \in \mathbb{N}$, then $a\geq0$.}
  %\label{fig:textfig1}
\end{marginfigure}
	\item  The Algebraic Limit Theorem ensures that the sequence $(b_{n} - a_{n})$  converges to $b-a$.  Because $b_{n} - a_{n} \geq 0$, we can apply the results from part (i) to get that $b-a \geq 0$, that is, $a \leq b$.
	\item  Take $a_{n} = c$ (or $b_{n} = c$) for all $n \in \mathbb{N}$ and apply our results from part (ii).

\end{enumerate}
\end{proof}


Loosely speaking, limits and their properties do not depend at all on what happens at the beginning of the sequence but are strictly determined by what happens as $n$ gets larger.  Changing the values of the first ten (or ten thousand) terms of the sequence has no effect on its eventual limit.


\medskip


\subsection{Exercises}
\noindent
\emph{Exercise 2.3.1.} Show that the constant sequence $(a,a,a,a,...)$ converges to $a$.
\begin{proof}
Let $\varepsilon > 0$.  Since the sequence is constantly $a$, we know that for all $n$:
	$$ |a_n - a| = |a-a| = 0 < \varepsilon $$
	Therefore we can choose $N$ completely arbitrarily since any choice will ensure that $|a_n -a| = 0 < \varepsilon$
\end{proof}

\smallskip

\noindent
\emph{Exercise 2.3.2.}
Let $x_n \geq 0$ for all $n \in \mathbb{N}$. 
\begin{enumerate}[(a)]
\item If $(x_n) \to 0$, show that $(\sqrt{x_n}) \to 0$. 
\item If $(x_n) \to x$, show that $(\sqrt{x_n}) \to \sqrt{x}$.
\end{enumerate}

\begin{proof}
\begin{enumerate}[(a)]
\item Let $\varepsilon> 0$.  (We want to find an $N$ such that $n \geq N$ implies $|\sqrt{x_n}| < \varepsilon$.)

Because $(x_n) \to 0$, there exists $N$ such that $n \geq N$ implies $|x_n|= x_n < \varepsilon^2$.  Then for the same $N$ and all $n \geq N$, $|\sqrt{x_n}| = \sqrt{x_n} < \sqrt{\varepsilon^2} = \varepsilon$.

As a result, for all $\varepsilon> 0$, there exists $N$ such that $n \geq N$ implies that $|\sqrt{x_n}- 0| < \varepsilon$ and $(\sqrt{x_n}) \to 0$.



\item The case $x = 0$ is part (a), so assume that $x > 0$.

Let $\varepsilon > 0$.  We wish to prove that there exists $N$ such that $n \geq N$ implies that $|\sqrt{x_n} - \sqrt{x}| < \varepsilon$.  

$\begin{array} {lcl} 
|\sqrt{x_n} -  \sqrt{x}|& = &|\sqrt{x_n} -  \sqrt{x}| \frac{\sqrt{x_n} +  \sqrt{x}}{\sqrt{x_n} +  \sqrt{x}}\\ 
	& = & \frac{|x_n -  x|}{\sqrt{x_n} +  \sqrt{x}}\\
&\leq	& \frac{|x_n -  x|}{\sqrt{x}}
    \end{array}$
    
    Because $(x_n ) \to x$ and $x>0$, we can choose $N$ such that $|x_n -x | <\varepsilon \sqrt{x}$ whenever $n \geq N$.  This implies that $\forall n \geq N$,
    \[		|\sqrt{x_n} - \sqrt{x}| \leq \frac{|x_n -  x|}{\sqrt{x}} < \frac{\varepsilon \sqrt{x}}{\sqrt{x}} = \varepsilon	\]

\end{enumerate} \end{proof}

\smallskip

\noindent
\emph{Exercise 2.3.3.} (Squeeze Theorem). Show that if $x_n \leq y_n \leq z_n$ for all $n \in \mathbb{N}$, and if $\lim x_n = \lim z_n = l$, then $\lim y_n = l$ as well.
\begin{proof}
Let $\varepsilon > 0$.  Since $(x_n) \to l$, there exists $N_{1}$ such that $n \geq N_{1} \implies x_n \in (l- \varepsilon, l + \varepsilon)$.

Since $(z_n) \to l$, there exists $N_{2}$ such that $n \geq N_{2} \implies z_n \in (l- \varepsilon, l + \varepsilon)$.

Since $x_n \leq y_n \leq z_n$ for all $n \in \mathbb{N}$, setting $N = \max \{N_1, N_2\}$ it follows that $y_n \in (l-\varepsilon, l + \varepsilon)$ for all $n \geq N$.

\end{proof}

\smallskip


\noindent
\emph{Exercise 2.3.4.}  Show that limits, if they exist, must be unique. In other words, assume $\lim a_n = l_1$ and $\lim a_n = l_2$, and prove that $l_1 = l_2$.
\begin{proof}
\begin{enumerate}[(i)]
	\item    Using the definition of convergence
	\noindent \marginnote{ \begin{proof} $\Rightarrow$ trivial.
		
		$\Leftarrow$ Suppose $l_1 \neq l_2$ and set $\varepsilon_{0} = |l_1 - l_2| > 0$.  Then the inequality $|l_1 - l_2| < \varepsilon$ fails for the particular choice $\varepsilon = \varepsilon_{0}$.  Contradiction.
		\end{proof}}\begin{lemma}  Two numbers $l_1$ and $l_2$ are equal iff $|l_1 - l_2 |< \varepsilon$ for all $\varepsilon > 0$.
		\end{lemma} 
	
	Let $\varepsilon > 0$.  Since $(a_n) \to l_1$ and $(a_n) \to l_2$, there exist $N_1, N_2$ such that $n \geq N_1$ and $n \geq N_2$ implies that $|a_n - l_1| < \varepsilon/2$ and $|a_n - l_2| < \varepsilon/2$  respectively.
	
	Therefore choosing $N = \max\{N_1, N_2\}, \quad n \geq N$ implies:


$\begin{array} {lcl} 
|l_1 - l_2|& = &|l_1 + a_n - a_n - l_2 | \\ 
	& \leq & |l_1 - a_n | +  |a_n - l_2 |\\
	& = & |a_n - l_1 | + |a_n - l_2 |\\
	& < & \frac{\varepsilon}{2} + \frac{\varepsilon}{2} \\
	&=& \varepsilon
    \end{array}$	 
    
    Since $\varepsilon$ was arbitrary, the lemma gives $l_1 = l_2$.

	\item  By the Algebraic Limit Theorem
	
	$\begin{array} {lclcl} 
\lim (a_n - a_n)& = &\lim (a_n) -\lim (a_n)  &=& l_{1} - l_{2}\\ 
			& = &\lim 0 			&=& 0 \\
    \end{array}	$
    
    Therefore $l_1 - l_2 = 0$ implying that $l_1 = l_2$.
	
\end{enumerate}
\end{proof}

\smallskip


\noindent
\emph{Exercise 2.3.6.} 
\begin{enumerate}[(a)]
\item Show that if $(b_n) \to b$, then the sequence of absolute values $(|b_n|)$ converges to $|b|$.


\item Is the converse of part (a) true? If we know that $|b_n| \to |b|$, can we deduce that $(b_n) \to b$?
\end{enumerate}

\begin{proof}
\begin{enumerate}[(a)]
\item We wish to show that $\forall \varepsilon > 0 \quad |b_n - b| < \varepsilon  \implies ||b_n| - |b|| < \varepsilon$.

$\begin{array} {rcrcl} 
|b_n|& = &|b_n - b + b|  &\leq & |b_n - b| + |b|\\ 
	&&	|b_n| - |b|	& \leq & |b_n - b|	\\
||b_n| - |b|| &\leq&    ||b_n - b|| &=& |b_n - b|\\
    \end{array}	$
    
    (The last line uses the same argument with the roles of $b_n$ and $b$ exchanged, which gives $|b| - |b_n| \leq |b_n - b|$ as well.)  Now let $\varepsilon > 0$.  Since $(b_n) \to b$, there exists $N$ such that $n \geq N \Rightarrow |b_n - b| < \varepsilon$.
    
    Therefore, for the same $N$ and all $n \geq N$,
    $$ ||b_n| - |b|| \leq     |b_n - b| < \varepsilon $$
    and $(|b_n|)$ converges to $|b|$.
    
\item No, the converse is not true.  Take $b_n = (-1)^n$: then $(|b_n|)$ converges to $|1|$, but $(b_n)$ does not converge to $1$ (or to anything else).
\end{enumerate}
\end{proof}

\smallskip


\noindent
\emph{Exercise 2.3.7.} 
\begin{enumerate}[(a)]
	\item Let $(a_n)$ be a bounded (not necessarily convergent) sequence, and assume $\lim b_n = 0$. Show that $\lim(a_n b_n ) = 0$. Why are we not allowed to use the Algebraic Limit Theorem to prove this?

\item Can we conclude anything about the convergence of $(a_n b_n)$ if we assume that $(b_n)$ converges to some nonzero limit $b$?
%\item Use $(a)$ to prove Theorem 2.3.3, part (iii), for the case when $ a = 0$.
\end{enumerate}

\begin{proof}

\begin{enumerate}[(a)]
	\item   Let $\varepsilon > 0$.  Because $(a_n)$ is bounded, there exists $M > 0$ such that $|a_n| \leq M$ for all $n \in \mathbb{N}$.
	
	Furthermore since $(b_n) \to 0$, there exists $N$ such that $n \geq N$ implies that $|b_n| < \frac{\varepsilon}{M}$.
	
	Then the sequence $(a_n b_n) $ is such that for $n \geq N$
	$$ |a_n b_n | \leq |a_n| |b_n| < M \frac{\varepsilon}{M} = \varepsilon $$
	
	
	and $a_n b_n  \to 0$
	
	
	We cannot use the Algebraic Limit Theorem to prove this since $(a_n)$ is not assumed to converge.
	
	
	

\item No, we cannot conclude anything about the sequence.  With $b_n = 1$ for every $n$ the product is just $(a_n)$ itself, which may converge or diverge. 
\end{enumerate}
\end{proof}

\noindent
\emph{Exercise 2.3.10.}
If $(a_n) \to 0$ and $|b_n - b| \leq a_n$ for all $n$, then $(b_n) \to b$.

\begin{proof}
	Since $0 \leq |b_n - b| \leq a_n$, every $a_n$ is nonnegative.  Therefore $|a_n| = a_n$.
	
	Let $\varepsilon > 0$.  Since $(a_n) \to 0$, there exists $N$ such that $n \geq N$ implies $|a_n | = a_n < \varepsilon$.
	
	Then the inequality above gives $|b_n - b| \leq a_n < \varepsilon$ whenever $n \geq N$.   Therefore $(b_n) \to b$.
\end{proof}


\noindent
\emph{Exercise 14.}

  If $\{s_{n}\}$ is a complex sequence, define its arithmetic means $\sigma_{n}$ by
 $$  \sigma_{n} = \frac{s_{0} + s_{1} + \cdots + s_{n}}{n+1} \qquad (n=0,1,2,\ldots) $$
 
 \begin{enumerate}[(a)]
 	\item  If $\lim s_{n} = s$, prove that $\lim \sigma_{n} = s$.
	
	\begin{proof}
			Let $\varepsilon >0$.
			
			If $\lim s_{n} = s$  then the sequence is bounded, that is, there exists an $M >0$ such that $$ |s_n - s| < M \quad \text{for all } n. $$
			
			Furthermore, there exists an $N_{1} \in \mathbb{N}$ such that $n \geq N_{1}$ implies that $$  |s_n - s|< \frac{\varepsilon}{2} $$
			
			Therefore for $n \geq N_{1}$ we have that 
			
			\begin{eqnarray*}
|\sigma_n - s|  &=& \big|\frac{s_0 + s_1 + s_2 + \cdots + s_n}{n+1}  - s\big|  \\
   	&=& \big|\frac{s_0 + s_1 + s_2 + \cdots + s_n - (n+1)s}{n+1}\big| \\
   &=& \big|\frac{(s_0-s) + (s_1 -s) + (s_2 -s)  + \cdots + (s_n - s)}{n+1}\big| \\
   &\leq & \big|\frac{(s_0-s) + (s_1 -s) +\cdots + (s_{N_{1} - 1} -s) }{n+1}\big|  + \big|\frac{(s_{N_{1}} -s) + \cdots + (s_n - s)}{n+1}\big| \\
   &<& \frac{N_{1} M}{n+1}  + \frac{\varepsilon}{2}\frac{(n - N_{1} + 1)}{n+1}\\
   &<& \frac{N_{1} M}{n+1}  + \frac{\varepsilon}{2}
\end{eqnarray*}

Note that in the last line we used the fact that for $n \geq N_1$, $\frac{(n - N_{1} + 1)}{n+1} < 1$. 


But $N_{1}$ and $M$ are fixed meaning they are simply constants.  Since the sequence $\{\frac{1}{n+1}\} \rightarrow 0$ we can make the quantity $\frac{N_{1} M}{n+1}$ to be arbitrarily small.

Choose $N_2$ such that $n \geq N_2$ implies that $$ \frac{N_{1} M}{n+1} < \frac{\varepsilon}{2} $$




Finally, let $N = \max\{N_1, N_2\}$.   Then $n \geq N$ implies that 


\begin{eqnarray*}
|\sigma_n - s|  &<& \frac{N_{1} M}{n+1}  + \frac{\varepsilon}{2}\\
		&<& \frac{\varepsilon}{2} + \frac{\varepsilon}{2} \\
		&=&  \varepsilon
\end{eqnarray*}

as desired.	
	\end{proof}
	\item  Construct a sequence $\{s_{n}\}$ which does not converge although $\lim \sigma_{n} = 0$.
	
	\noindent
	\emph{Solution:}   For any $a \neq 0$, the sequence $\{s_{n}\} = (a, -a, a, -a, \ldots )$ does not converge, but the sequence $\{\sigma_{n}\}$ converges to $0$.
 \end{enumerate}

 


\smallskip

\section{The Monotone Convergence Theorem }
We showed previously that convergent sequences are bounded.  While the converse is not necessarily true, if a sequence is bounded as well as monotonic, then it converges.


\begin{definition}Increasing sequence:  A sequence $\{a_{n}\}$ is \emph{increasing} if $a_{n} \leq a_{n+1}$ for all $n \in \mathbb{N}$ and \emph{decreasing} if $a_{n} \geq a_{n+1}$ for all $n \in \mathbb{N}$.  A sequence is \emph{monotone} if it is either increasing or decreasing. \end{definition}

\begin{theorem}  [The Monotone Convergence Theorem] If a sequence is monotone and bounded, then it converges.
\end{theorem}

\begin{proof}
	Let $(a_n)$ be monotone and bounded.
	WLOG assume $(a_n)$  is increasing and consider the set of points $\{a_n : n \in \mathbb{N}\}$.
	
	By assumption, the set is bounded.  By the axiom of completeness, the set has a supremum.  Let $s = \sup\{a_n : n \in \mathbb{N}\}$.
	
	We claim that $\lim (a_n) = s$.
	
	Let $\varepsilon > 0$.  Because $s$ is the least upper bound of $\{a_n : n \in \mathbb{N}\}$, $s- \varepsilon$ is not an upper bound, so there exists  a point in the sequence call it $a_N$ such that $s-\varepsilon < a_N$.
	
	Now the fact that $(a_n)$ is increasing implies that if $n \geq N$ then $a_N \leq a_n$.
	
	Hence,
	$$ s- \varepsilon < a_N \leq a_n \leq s < s+\varepsilon $$
	which implies that $|a_n  - s| < \varepsilon$ for all $n \geq N$.
\end{proof}

Note that the Monotone Convergence Theorem is extremely useful since it asserts the convergence of a sequence without explicit mention of the actual limit.

Also note that the Monotone Convergence Theorem follows from the Axiom of Completeness.

\subsection{Exercises}

	\emph{Exercise 2.4.2.}
		\begin{enumerate}[(a)]
			\item Prove that the sequence defined by $x_1 = 3$ and 
				\[		x_{n+1} = \frac{1}{4-x_n}\]
				converges.
			\item Now that we know $\lim x_n$ exists, explain why $\lim x_{n+1}$ must also exist and equal the same value.
			\item Take the limit of each side of the recursive equation in part (a) of this exercise to explicitly compute  $\lim x_n$.
		\end{enumerate}
		
		\begin{proof}
			\begin{enumerate}[(a)]
				\item  
				\begin{enumerate}[(i)]
					\item The first few terms are $3, 1, \frac{1}{3}, \frac{3}{11}, \ldots$.  Then we can guess that the sequence is decreasing ($x_{n+1}< x_{n} \quad \forall n$)  and we can prove it by induction.
				
				For $n=1$,  $x_{2}< x_{1}$ since $1 <3$.
				
				For $n=k$, assume that $x_{k+1} < x_{k}$.
				
				Then when $n=k+1$, we want to be able to show that $x_{k+2} =\frac{1}{4-x_{k+1}} < x_{k+1} = \frac{1}{4-x_{k}}  $.   Using the inductive hypothesis, we have that
				
				
				\begin{eqnarray*}
x_{k+1}  &< & x_{k} \\
   	-x_{k+1}  &>& -x_{k} \\
 4-x_{k+1}  &>&4 -x_{k}\\
    \frac{1}{4-x_{k+1}}  &<&\frac{1}{4 -x_{k}}\\
   x_{k+2}  &< & x_{k+1} \\
\end{eqnarray*}

as desired.
				\item We can see that $(x_n)$ is bounded above by $3$ but we would like to prove that it is bounded below by $0$.  Again we can prove this by induction, showing that $x_n > 0 \quad \forall n$.
				
				Clearly for $n=1$, $x_1 = 3 >0$.
				Then assume for $n=k$, that $x_k >0$.
				
				To show this for $n = k+1$, note that we have that $3>x_k>0$.  This implies that $x_{k+1} = \frac{1}{4-x_k}$  is positive.  Therefore $x_{k+1}>0$ as desired.
				
				
				Since $(x_n)$ is decreasing and bounded below by $0$, we have by the MCT that the sequence is convergent.
				
								\end{enumerate}
				
				\item Since $ (x_n)$ is bounded and decreasing, $(x_{n+1})$ must also be bounded and decreasing, then by the MCT it must also be convergent.
				
				Furthermore since $(x_n) \to L$, it must be that for every $\varepsilon >0$ there exists $N \in \mathbb{N}$ such that $n \geq N$ implies that $|x_n - L| <\varepsilon$.  However, notice that for $(x_{n+1})$ this same $N$ works since $n+1> n \geq N$ and therefore $|x_{n+1} - L| <\varepsilon$ and both sequences converge to the same limit.
				
							\item By part (b), we have that  $\lim x_n = \lim x_{n+1} = L$. Furthermore, by the Algebraic Limit Theorem we can deduce that $L = \lim x_{n+1} = \lim \frac{1}{4-x_n} = \frac{1}{4-L}$ and we have that $L$ satisfies the equation $L^2 - 4L +1 =0$.  Solving for $L$ we have that $L = 2 \pm \sqrt{3}$, however since $3$ is an upper bound for the sequence, we must have that $L= 2 - \sqrt{3}$.
			\end{enumerate}
		
		\end{proof}
		
\noindent		
\emph{Exercise 2.4.3.} Using the same methods, prove that the following sequence is convergent and find its limit.
	$$y_{n+1} = 4 - \frac{1}{y_n}	\qquad y_1 = 1$$
	
	\begin{proof}
		\begin{enumerate}[(i)]
			\item Monotone.  Looking at the first few terms of the sequence again, we can see that it is increasing: $y_n \leq y_{n+1}$ for all $n$.  To prove this we use induction (together with the fact, shown in part (ii), that every term is positive).
				
				For $n=1$, $y_1 \leq y_2$ since $1 < 3$.
				
				Assume that for $n=k$, it is true that $y_{k} \leq y_{k+1}$.
				
				Then for $n = k+1$, 
				\begin{eqnarray*}
y_{k}  &\leq & y_{k+1} \\
   	\frac{1}{y_{k}}  &\geq& \frac{1}{y_{k+1}} \\
    -\frac{1}{y_{k}}  &\leq& -\frac{1}{y_{k+1}}\\
      4- \frac{1}{y_{k}}  &\leq&4- \frac{1}{y_{k+1}}\\
   y_{k+1}  &\leq& y_{k+2} \\
\end{eqnarray*}

as desired.

\item Bounded.  We show by induction that $1 \leq y_n < 4$ for every $n$.  For $n=1$ this is trivial.  Assume then that $1 \leq y_k < 4$.  Then $0 < \frac{1}{y_k}\leq 1$, and so $3 \leq y_{k+1} = 4 - \frac{1}{y_k} < 4$ as well.

\item Limit.  By the MCT the sequence converges; call its limit $y$, so that $\lim y_n = \lim y_{n+1} = y$.  Then $y = \lim y_{n+1} =  \lim (4 - \frac{1}{y_n}) = 4 - \frac{1}{y}$.  By this relation we have that $y^2 - 4y + 1 =0$ with roots $y = 2 \pm \sqrt{3}$.  However, since the sequence is bounded below by 1, we have that $y = 2 + \sqrt{3}$.

				
		\end{enumerate}
	\end{proof}
	
\emph{Exercise 2.4.4.}  Prove that the sequence $x_n = \sqrt{2}, \sqrt{2\sqrt{2}}, \sqrt{2\sqrt{2\sqrt{2}}}, \ldots$
	converges and find the limit.
	
	\begin{proof}
	First, define the sequence recursively:
	$$  x_{n+1} = \sqrt{2 x_n} \qquad x_1 = \sqrt{2} $$
		\begin{enumerate}[(i)]
		
			\item Monotone Increasing  ($x_{n} < x_{n+1}\quad \forall n$).  For $n=1, \quad \sqrt{2} < \sqrt{2\sqrt{2}}$. Then assume that for $n=k$, $x_k < x_{k+1}$.
			
			Using the Inductive hypothesis
			\begin{eqnarray*}
x_{k}  & <& x_{k+1} \\
   	2x_{k}  &<& 2x_{k+1} \\
    \sqrt{2x_{k}}  &<&\sqrt{ 2 x_{k+1}}\\
    x_{k+1}  & <& x_{k+2}
\end{eqnarray*} as desired.

\item Bounded.  We wish to prove that the sequence is bounded above by 2.
For $x_1$ this is trivial.  Assume that $x_k < 2$.
Then 
\marginnote{Notice that here we take the limit under a square root which is not covered in the Algebraic Limit Theorem.  We therefore use Exercise 2.3.2 to support our argument.}
\begin{eqnarray*}
x_{k}  & <& 2 \\
   	2x_{k}  &<& 4 \\
    \sqrt{2x_{k}}  &<&\sqrt{ 4}\\
 & <& 2 \\
 x_{k+1}  & <& 2
\end{eqnarray*} as desired.   


\item Limit.  Note that $x = \lim x_n = \lim x_{n+1}$.  Therefore we have that $x = \lim x_{n+1} = \sqrt{2x}$ and we have that $x^2 - 2x = 0$, so $x = 0$ or $x = 2$.  Since the sequence increases from $x_1 = \sqrt{2} > 0$, the limit cannot be $0$.  Therefore $x = 2$.
		\end{enumerate}
	\end{proof}
	
	\emph{Exercise 2.4.6.} Let $(a_n)$ be a bounded sequence.
	
	\begin{enumerate}[(i)]
	\item Define the sequence $y_n = \sup \{a_k : k \geq n\}$.  Prove that $\lim y_n$ exists.  This limit is called the \emph{limit superior} of $(a_n)$ and is written $\limsup a_n$.
	\item Define the sequence $z_n = \inf \{a_k : k \geq n\}$ and prove that it converges.  Its limit is the \emph{limit inferior}, written $\liminf a_n$.
	\item Show that $\limsup a_n = \liminf a_n$ if and only if $\lim a_n$ exists, in which case all three values are equal.
	\end{enumerate} % COMPLETE


\section{Subsequences and Bolzano-Weierstrass Theorem}

\begin{definition}Subsequences:  Given a sequence $\{p_n\}$ consider a sequence $\{n_k\}$ of positive integers such that $n_1 < n_2 < n_3 < \cdots$.  Then the sequence $\{p_{n_k}\}$ is called a \emph{subsequence} of $\{p_n\}$.  If $\{p_{n_k}\}$ converges, its limit is called a \emph{subsequential limit} of $\{p_{n}\}$.
\end{definition}
\medskip
Note:  It is clear that $\{p_{n}\}$ converges to $p$ if and only if every subsequence of $\{p_{n}\}$ converges to $p$. 

Notice also that the order of the terms in the subsequence is the same as in the original sequence and that repetitions are not allowed.

\begin{theorem} Subsequences of a convergent sequence converge to the same limit as the original sequence.\end{theorem}

\begin{proof} 
Assume the sequence $(x_n) \to L$ and let $(x_{n_j})$ be a subsequence of this sequence.  
Let $\varepsilon > 0$.  Then there exists an $N$ such that $n \geq N$ implies that $|x_n - L| < \varepsilon$.  But since $n_j \geq j$ this same $N$ works for the subsequence as well.  To be more precise, $j\geq N$ implies that $n_j \geq N$ and that $|x_{n_j} - L| < \varepsilon$ as desired.


\end{proof}

This theorem can sometimes provide a very useful proof for divergence of a sequence.  If we can find two subsequences that converge to two different limits, then we can rigorously conclude that the original sequence diverges.

Furthermore, this theorem can be applied to infinite series to show that infinite sums are associative which we will see later.

\begin{example} Let $0 < b< 1$.  Because $b >b^2 > b^3 > b^4 > \ldots > 0$ the sequence $(b^n)$ is decreasing and bounded below.  By the MCT the sequence converges to some limit $l$ satisfying $b > l \geq 0$.   Notice however that $(b^{2n})$ is a subsequence, so by the above theorem $(b^{2n}) \to l$. But $(b^{2n}) = (b^n)( b^n)$ so by the Algebraic limit theorem, $(b^{2n}) \to l \cdot l = l^2$.  Because limits are unique, $l^2 = l$, and since $l \leq b < 1$ rules out $l = 1$, we conclude that $l=0$.
\end{example}

\begin{example}
Consider the sequence 
$$\big( 1, - \frac{1}{2}, \frac{1}{3}, -\frac{1}{4}, \frac{1}{5}, -\frac{1}{5}, \frac{1}{5}, -\frac{1}{5}, \frac{1}{5}, -\frac{1}{5},\ldots	\big)$$

We can show that the sequence is divergent since we can provide two subsequences, namely:
$$ \big( \frac{1}{5}, \frac{1}{5}, \frac{1}{5}, \frac{1}{5},  \ldots \big) \quad \text{and} \quad \big( -\frac{1}{5},-\frac{1}{5},-\frac{1}{5},-\frac{1}{5}, \ldots \big) $$

that converge to two different limits.
\end{example}

\begin{theorem}\newthought{Bolzano-Weierstrass Theorem:} Every bounded sequence contains a convergent subsequence.\end{theorem}

\begin{proof}
Let $(a_n)$ be a bounded sequence such that $|a_n| \leq M$ for all $n$.  Bisect the closed interval $[-M, M]$ into $[-M, 0]$ and $[0,M]$ (the midpoint belongs to both halves).  Now it must be that at least one of these closed intervals contains infinitely many terms of the sequence $(a_n)$.  Select the half for which this is the case and label it $I_1$.  Then let $a_{n_1}$ be some point in the sequence satisfying $a_{n_1} \in I_1$.

Now bisect $I_1$ and again select the half with infinitely many terms of $(a_n)$ and call it $I_2$.  Then again select a point $a_{n_2} \in I_2$ such that $n_2> n_1$.

Continue in this way, constructing $I_{k}$ by halving $I_{k-1}$ and choosing the half with infinitely many terms of the sequence $(a_n)$ and then selecting $a_{n_k} \in I_{k}$ such that
\[	n_k  > n_{k-1} > n_{k-2} > \cdots > n_{2} > n_{1}	\]


We want to show that $(a_{n_k})$ converges but we first need to find a candidate for the limit.

Note that 
\[	I_1 \supseteq I_2 \supseteq I_3 \supseteq I_4 \supseteq	\cdots \]
and that each $I_i$ is closed.

By the NIP, there exists $x \in \mathbb{R}$ such that $x \in I_k$ for every $k$.


Let $\varepsilon >0$.  By construction, the length of $I_k$ is $M (\frac{1}{2})^{k-1}$, which tends to $0$ as $k \to \infty$.

Therefore there exists $N$ such that $k \geq N$ implies $M (\frac{1}{2})^{k-1} < \varepsilon$.  But $x$ and $a_{n_k}$ both lie in $I_k$, so the distance between them is at most the length of $I_k$.  As a result, $|x-a_{n_k}|< \varepsilon$ for all $k \geq N$, and $(a_{n_k}) \to x$.

\end{proof}

\section{The Cauchy Criterion}  
\newthought{Cauchy Sequence:}  A sequence $(a_n)$ is called a \emph{Cauchy sequence} if, for every $\varepsilon >0$ there exists an $N \in \mathbb{N}$ such that whenever $m,n \geq N$ it follows that $|a_n - a_m| < \varepsilon$

That is, a sequence is Cauchy if, for every $\varepsilon$, there is a point in the sequence after which all the terms are closer to each other than the given $\varepsilon$.  To spoil the surprise, we will argue that Cauchy sequences are convergent and that convergent sequences are Cauchy.  This, like the Monotone Convergence Theorem, allows us to prove that a sequence converges without having explicit knowledge of what the limit might be. 


\begin{lemma} Cauchy Sequences are bounded.\end{lemma}

\begin{proof}
	Given $\varepsilon = 1$, there exists an $N$ such that $|x_m - x_n| <1$ for all $m,n \geq N$.  
	As a result \marginnote{Note that $|x_n| = |x_n + x_N - x_N| \leq |x_n -x_N| + |x_N| < 1 + |x_N|$}  $|x_n| < |x_N| + 1$ for all $n \geq N$. Note that there were finitely many terms before $x_N$ in the sequence. Therefore let
	
	
	\[	M	= \max \{|x_1|,|x_2|, \ldots, |x_{N-1}|, |x_N| + 1 \}\]
	And this value is a bound for the sequence.
\end{proof}


\begin{theorem}\newthought{Cauchy Criterion:} A sequence converges iff it is a Cauchy sequence. \end{theorem}

\begin{proof}
	$\Longrightarrow$.  This we can prove by applying the triangle inequality. Assume that a sequence $(x_n)$ converges to a limit $x$.   Then $\forall \varepsilon >0$, there exists an $N$ such that $n,m \geq N$ implies that 
	\[ |x_n -x | < \varepsilon/2	\qquad 	|x_m - x| < \varepsilon/2	\]
Applying the triangle inequality, we can state that 

$$
\begin{array}{rcl}
|x_n - x_m| &=& |x_n - x  + x - x_m| ,\\
&\leq & |x_n - x| + |x - x_m|,\\
& < & \varepsilon\\
\end{array}
$$

as desired.   

	$\Longleftarrow$.  So the other direction is a lot harder to prove and we have to invoke Bolzano-Weierstrass to get our result.
	
	Assume we have a Cauchy sequence $(x_n)$.  By the lemma above we know that it is bounded and so by BW we know that it contains a convergent subsequence.   Let us call this subsequence $(x_{n_k})$ and call its limit $x$.  
	
	The idea here is that because the sequence is Cauchy, we know that elements in the ``tail'' end of the sequence are getting closer together.  But because the terms in the subsequence are getting closer to the limit $x$, we can force each value to be less than $\varepsilon/2$ and apply the triangle inequality.
	
	So let $\varepsilon >0$.
	
		Because $(x_n)$ is Cauchy there is an $N$ such that $m,n \geq N$ implies that \[|x_n - x_m|<\varepsilon/2.\]  We also have that the subsequence converges so we can choose an index, let us say $n_K$, such that $n_K \geq N$ and 
		\[	|x_{n_K} - x|<\varepsilon/2 \]
		
		
		Now take any $n \geq N$.  Since $n_K \geq N$ as well, the Cauchy condition applies to the pair $n$, $n_K$, and so 
		
				$$
				\begin{array}{rcl}
|x_n - x| &=& |x_n - x_{n_K}  + x_{n_K} - x| ,\\
&\leq & |x_n - x_{n_K}| + |x_{n_K} - x|,\\
& < & \varepsilon
\end{array}
$$
	
\end{proof}



One important thing to note is that many of our significant results up to this point rely heavily on the Axiom of Completeness.

The AoC was what allowed us to prove the NIP which allowed us to prove Bolzano-Weierstrass which allowed us to prove the Cauchy Criterion.  Furthermore the AoC also gave us the Monotone Convergence Theorem.

$$
\begin{array}{rcl}
AoC &\implies & NIP \implies BW \implies CC ,\\
&\implies & MCT,\\
\end{array}
$$

Also note that we could have begun this chain with \emph{any} of these statements and arrived at the others as a result.\marginnote{A subtlety: starting from the Nested Interval Property or the Cauchy Criterion, one must also assume the Archimedean Property in order to recover the Axiom of Completeness.}  This is pretty remarkable.  In other words all of these statements are equivalent and that assuming any one of them to be true, means that we can derive the others.   This is because all of these statements essentially say the same thing : that the field of real numbers is \emph{complete}.   Remember how we could not show the AoC in the field of rational numbers.  This is because the rationals did not contain a certain property, independent of the field and order axioms, that forced this result to be true.  In the next section we try to define more explicitly what this ``completeness'' property is and how it relates to the topology of the field in question.


\section{Infinite Series}


\newthought{Infinite Series:}  Let $(b_n)$ be a sequence.  An infinite series is an expression of the form
		\[\sum_{n=1}^{\infty} b_n = b_1 + b_2 + b_3 + \cdots \]
		
		The sequence of \emph{Partial sums} of the series $(s_m)$ is defined as 
			\[s_m = b_1 + b_2 + b_3 + \cdots + b_m \]
			
		We say that the series $\sum_{n=1}^{\infty} b_n$ \textsc{converges} if its sequence of partial sums converges to some limit B and we denote this limit as $\sum_{n=1}^{\infty} b_n  = B$.
		Specifically,
		\[\sum_{n=1}^{\infty} b_n = B \qquad \text{means that} \qquad \lim s_n = B\]
		
		
		\begin{dangerousbend}
			It is important to keep a distinction between 
				\begin{itemize}
					\item The sequence of \emph{terms}: $(b_1, b_2, b_3, \ldots)$ and
					\item The sequence of \emph{partial sums}:  $(s_1, s_2, s_3, \ldots)$ where $s_n = b_1 +  \cdots + b_n$
				\end{itemize}
		\end{dangerousbend}
		
		Nonetheless it is this close relationship between series and sequences that allows us to apply much of our previous knowledge to them.
		
		
		\begin{theorem}\newthought{Theorem: Algebraic Limit Theorem for Series}
		If $\sum_{n=1}^{\infty} a_n = A$ and $\sum_{n=1}^{\infty} b_n = B$ then,
			\begin{enumerate}[(1)]
				\item $\sum_{n=1}^{\infty} c a_n = cA$
				\item $\sum_{n=1}^{\infty} (a_n + b_n) = A + B$
			\end{enumerate}
		\end{theorem}
		
		\begin{proof}
		\end{proof}
		
		\begin{theorem} \newthought{Cauchy Criterion for Series}
		The series $\sum_{n=1}^{\infty} a_n$ converges 
		
		\[\text{if and only if}\]
		
		
		 for every $\varepsilon > 0$ there exists an $N \in \mathbb{N}$ such that $n>m\geq N$ implies that $|a_{m+1} + a_{m+2} + \cdots + a_n| < \varepsilon$
		\end{theorem}
		
		
	\begin{proof}  \end{proof}
	
	This can help us prove some basic facts about series.
	
	
	\begin{theorem} If the series $\sum_{n=1}^{\infty} a_n$ converges then $(a_k) \to 0$.  \marginnote{The converse statement is \emph{not} true! Take for example the harmonic series.}
	\end{theorem}
	
	\begin{proof}  \end{proof}
	
	\begin{theorem}\newthought{The Comparison Test}  Assume that $(a_n)$ and $(b_n)$ are sequences where $0 \leq a_n \leq b_n$  for all $n \in \mathbb{N}$
		\begin{enumerate}
			\item If $\sum_{n=1}^{\infty} b_n$ converges, then $\sum_{n=1}^{\infty} a_n$ converges.
			\item If $\sum_{n=1}^{\infty} a_n$ diverges, then $\sum_{n=1}^{\infty} b_n$ diverges.
		\end{enumerate}
	\end{theorem}
	
	
	
	\begin{proof}  \end{proof}
	
	Remember also that statements about the convergence of a sequence (or series) only care about the tail ends of the sequence.  In other words, we can change around a few things about some finite number of initial terms without affecting the infinite behavior of the sequence (or series).
	
	Because the Comparison Test is about deducing the behavior of a series in comparison to another, it might be a good idea to have an arsenal of basic series to which we can compare another. 
	
	\newthought{Absolute Convergence:}  If $\sum_{n=1}^{\infty} |a_n|$ converges then we say that $\sum_{n=1}^{\infty} a_n$ converges absolutely.  On the other hand if $\sum_{n=1}^{\infty} a_n$ converges and $\sum_{n=1}^{\infty} |a_n|$ diverges we say that the series \emph{converges conditionally}.
	
	
	\begin{theorem}\newthought{Absolute Convergence Test}
	If the series $\sum_{n=1}^{\infty} |a_n|$ converges, then $\sum_{n=1}^{\infty} a_n$ converges as well.
	\end{theorem}  \marginnote{The converse of this theorem is also false!}
	\begin{proof}  \end{proof}
	
	
	\begin{theorem}\newthought{Alternating Series Test}  Let $(a_n)$ be a sequence satisfying,
	
	\begin{enumerate}
		\item The sequence is monotone decreasing:  $a_1 \geq a_2 \geq a_3 \geq \cdots a_n \geq a_{n+1} \geq \cdots $ and
		\item $(a_n) \to 0$
	\end{enumerate}
	
	
	Then the alternating series $\sum_{n=1}^{\infty} (-1)^{n+1} a_n$ converges.
	\end{theorem}
	\begin{proof}  \end{proof}
	
	\begin{theorem}\newthought{Ratio Test}
	Given a series $\sum_{n=1}^{\infty} a_n$ with $a_n \neq 0$, the ratio test states that if $(a_n)$ satisfies 
	\[		\lim \left|\frac{a_{n+1}}{a_n}	\right| = r < 1\]
	then the series converges absolutely (and therefore converges).
	\end{theorem} 
	\begin{proof}  \end{proof}
	
	
	\subsection{Rearrangements}
		Imagine that we were to permute the terms in a sum into some other order such that all of the original terms eventually appeared in the new order and such that no terms were repeated.  Clearly there are an infinite number of possible rearrangements of a series.  
		
		More formally,
		
		\newthought{Rearrangement}  Let  $\sum_{n=1}^{\infty} a_n$ be a series.  Another series  $\sum_{n=1}^{\infty} b_n$ is called a rearrangement of  $\sum_{n=1}^{\infty} a_n$ if there exists a bijection $f: \mathbb{N} \to \mathbb{N}$ such that $b_{f(k)} = a_{k}$ for all $k$.
		
		It should be noted however, that rearrangements of series do \emph{not} always converge to the same limit as the original series.  This is only true for absolutely convergent series as we will now prove.
		
		
		\begin{theorem}
	If $\sum_{n=1}^{\infty} a_n$ converges absolutely, then any rearrangement of the series converges to the same limit.
	\end{theorem} 
	\begin{proof}  \end{proof}
	
	This theorem should make clear that absolute convergence is an extremely desirable property when manipulating series.  However it should be noted that conditional convergence has this amazing and not at all obvious property which we attribute to Riemann.
	
	
	\begin{theorem}\newthought{Riemann}
	If $\sum_{n=1}^{\infty} a_n$ converges conditionally, then for any real number $r$ we can find a rearrangement of the series that converges to $r$.
	\end{theorem} 
	\marginnote{This is Theorem 3.54 in Rudin's \emph{Principles of Mathematical Analysis}.}
		
		


\section{Problems}

\noindent The problems in this section are original to this text.

\begin{problem}[Three proofs on trial]
\hint{Read each proof against the template in Section~\ref{sec:sequence-examples}.  Who is supposed to speak first, $\varepsilon$ or $N$?  And for which $n$ must the final inequality hold?}
Each of the following ``proofs'' has something wrong with it.  In each case say whether the \emph{claim} is true, pinpoint the faulty step, and, where the claim is true, write a correct proof.
\begin{enumerate}[(a)]
	\item \emph{Claim.} $\lim \frac{1}{n} = 0$.

	\emph{``Proof.''}  Let $N \in \mathbb{N}$ be arbitrary and set $\varepsilon = 2/N$.  If $n \geq N$ then $\left|\frac{1}{n} - 0\right| = \frac{1}{n} \leq \frac{1}{N} < \frac{2}{N} = \varepsilon$.  Hence $\left(\frac{1}{n}\right) \to 0$.
	\item \emph{Claim.} the sequence $a_n = (-1)^n$ converges to $1$.

	\emph{``Proof.''}  Let $\varepsilon > 0$ and choose $N = 2$.  For every even $n \geq N$ we have $|a_n - 1| = |1 - 1| = 0 < \varepsilon$.
	\item \emph{Claim.} $\lim \frac{2n+1}{n} = 2$.

	\emph{``Proof.''}  Let $\varepsilon > 0$ and choose $N = 1/\varepsilon$.  If $n \geq N$ then $\left|\frac{2n+1}{n} - 2\right| = \frac{1}{n} \leq \frac{1}{N} = \varepsilon$.
\end{enumerate}
To convince yourself that (a) is a broken argument and not merely an unusual one, show that exactly the same reasoning ``proves'' that the constant sequence $(1, 1, 1, \ldots)$ converges to $0$.
\end{problem}

\begin{problem}[Sticking and visiting]
\hint{For the second half of (c), look back at Chapter One and ask how many rational numbers there are.  For (d), the Bolzano--Weierstrass Theorem will find you a second number to visit.}
Here are two invented definitions, built from the same ingredients as convergence but with the quantifiers rearranged.
\begin{itemize}
	\item A sequence $(a_n)$ \emph{sticks to} $a$ if there exists an $N \in \mathbb{N}$ such that for every $\varepsilon > 0$ and every $n \geq N$ we have $|a_n - a| < \varepsilon$.
	\item A sequence $(a_n)$ \emph{visits} $a$ if for every $\varepsilon > 0$ and every $N \in \mathbb{N}$ there exists an $n \geq N$ with $|a_n - a| < \varepsilon$.
\end{itemize}
\begin{enumerate}[(a)]
	\item Describe each definition in a plain sentence that mentions neither $\varepsilon$ nor $N$.
	\item Show that if $(a_n)$ sticks to $a$ then it converges to $a$, and that if it converges to $a$ then it visits $a$.  Give examples to show that neither implication can be reversed.
	\item Give an example of a sequence that visits exactly two numbers.  Then give an example of a single sequence that visits \emph{every} real number.
	\item Suppose $(a_n)$ is bounded, visits $a$, and visits no other number.  Prove that $(a_n) \to a$.  Show by example that the conclusion can fail if the word ``bounded'' is dropped.
\end{enumerate}
\end{problem}

\begin{problem}[Climbing toward a hole]
\hint{For (c), write $x_{n+1} - x_n$ as a single fraction and look at the sign of the numerator.}
Define a sequence by $x_1 = 1$ and
\[ x_{n+1} = \frac{2x_n + 3}{x_n + 2}. \]
\begin{enumerate}[(a)]
	\item Compute $x_2$, $x_3$ and $x_4$, both as fractions and as decimals.  Explain why every $x_n$ is a positive rational number.
	\item Verify the identity
	\[ 3 - x_{n+1}^2 = \frac{3 - x_n^2}{(x_n + 2)^2} \]
	and use it to prove that $x_n^2 < 3$ for every $n$.
	\item Show that $(x_n)$ is increasing, and conclude that it converges.  Find the square of its limit.
	\item Every term of this sequence lies in $\mathbb{Q}$, and within $\mathbb{Q}$ the sequence is increasing and bounded.  Which theorem of this chapter have you just shown to be false in $\mathbb{Q}$?
	\item Use the identity in (b) to show that $3 - x_n^2 \leq 2/9^{\,n-1}$.  How many terms are needed before $x_n^2$ is guaranteed to be within one millionth of $3$?
\end{enumerate}
Notice what part (c) accomplishes.  Without ever writing the symbol $\sqrt{3}$, you have proved that there is a positive real number whose square is $3$.
\end{problem}

\begin{problem}[Bolzano--Weierstrass by decimal digits]
\hint{Part (a) is the pigeonhole principle for infinitely many pigeons.  If infinitely many terms are shared among ten bins, the bins cannot all be finite.}
The proof of the Bolzano--Weierstrass Theorem in the text cuts an interval in half again and again.  This problem redoes it by cutting into tenths, which lets us watch the limit being written down one digit at a time.  Let $(a_n)$ be any sequence with $0 \leq a_n < 1$ for all $n$.
\begin{enumerate}[(a)]
	\item Sort the terms into ten bins $[0, 0.1)$, $[0.1, 0.2)$, \ldots, $[0.9, 1)$ according to their first decimal digit.  Explain why at least one bin receives $a_n$ for infinitely many values of $n$.  Call the digit of one such bin $d_1$.
	\item Split that bin into ten smaller bins according to the second decimal digit and repeat.  Explain how continuing in this way produces digits $d_1, d_2, d_3, \ldots$ and nested closed intervals $I_1 \supseteq I_2 \supseteq I_3 \supseteq \cdots$, where $I_k$ has length $10^{-k}$ and contains $a_n$ for infinitely many $n$.
	\item Build a subsequence with $a_{n_k} \in I_k$ for each $k$, and prove that it converges.  What is the decimal expansion of its limit?
	\item Carry out the construction for the sequence
	\[ 0.4,\quad 0.41,\quad 0.414,\quad 0.4142,\quad \ldots \]
	whose $n$th term consists of the first $n$ decimal digits of $\sqrt{2} - 1$.  Every term of this sequence is rational.  Is the limit?
	\item Halves work and tenths work.  Would it also work to cut each interval into two \emph{unequal} pieces, say in the ratio $1:9$?  What is the one property of the cutting scheme that the proof really needs?
\end{enumerate}
\end{problem}

\begin{problem}[A hiker's steps]
\hint{For (b), group the steps.  The third and fourth together cover at least $\frac{1}{2}$, as do the fifth through eighth, and so on.  For (c), let her turn around whenever she has wandered too far.}
A hiker walks along the number line.  Her position after $n$ steps is $a_n$, so her $n$th step has length $|a_{n+1} - a_n|$.
\begin{enumerate}[(a)]
	\item Suppose that $|a_{n+1} - a_n| \leq 2^{-n}$ for every $n$.  Prove that $(a_n)$ is a Cauchy sequence, so that she is certainly headed \emph{somewhere}.
	\item Suppose instead that every step is forward and the $n$th step has length exactly $1/n$, so that $a_{n+1} = a_n + \frac{1}{n}$.  Show that she eventually passes every milestone, which is to say that the sequence $(a_n)$ is unbounded.
	\item Now let the $n$th step again have length exactly $1/n$, but allow her to choose its direction.  Describe a rule for choosing directions that keeps her inside the interval $[-1, 2]$ forever and yet leaves $(a_n)$ divergent.  Prove both claims.
	\item In (b) and (c) the steps shrink to zero, and still the sequence is not Cauchy.  Explain in a sentence or two, with reference to the quantifiers in the definition, why ``consecutive terms get close together'' is a weaker statement than ``the sequence is Cauchy.''
\end{enumerate}
\end{problem}

\begin{problem}[The sceptic's series]
\hint{For (c), expand $(\sqrt{x} - \sqrt{y})^2 \geq 0$.  For (e), apply the Cauchy Criterion for series to the block $a_{n+1} + a_{n+2} + \cdots + a_{2n}$.}
Suppose that $a_n \geq 0$ for all $n$ and that $\sum_{n=1}^{\infty} a_n$ converges.  A sceptical friend proposes the statements below.  Decide which are true, prove those, and give a counterexample to each of the others.
\begin{enumerate}[(a)]
	\item $\sum_{n=1}^{\infty} a_n^2$ converges.  (Does your answer change if the terms are allowed to be negative?)
	\item $\sum_{n=1}^{\infty} \sqrt{a_n}$ converges.
	\item $\sum_{n=1}^{\infty} \sqrt{a_n a_{n+1}}$ converges.
	\item $n a_n \to 0$.
	\item If in addition $(a_n)$ is decreasing, then $n a_n \to 0$.
\end{enumerate}
\end{problem}

\bookchapter{Three}{Basic Topology of the Real Line}

Chapter Two ended with a promise.  Five statements---the Axiom of Completeness, the Nested Interval Property, the Monotone Convergence Theorem, the Bolzano--Weierstrass Theorem and the Cauchy Criterion---all turned out to be saying the same thing about $\mathbb{R}$, and we promised to say more precisely what that thing is.

To keep the promise we have to change what we look at.  So far the objects of study have been numbers, and sequences of numbers.  In this chapter the objects are \emph{sets} of real numbers, and the question we ask of a set is always some version of this one:

\begin{quote}
\emph{If I start inside the set and take a limit, where do I end up?}
\end{quote}

\noindent The study of that question is called \textsc{topology}.  It has a reputation for being abstract, but on the real line every one of its ideas can be drawn, and we will draw them.

\section{Where Can a Sequence Go?}\label{sec:three-escapes}

Here are three sets that look very much alike:
\[ A = (0,1], \qquad B = [0, \infty), \qquad K = [0,1]. \]
Let us release a sequence inside each one and watch where it goes.

\newthought{In the set $A$}, take $a_n = 1/n$.  Every term belongs to $A$, and the sequence converges.  But its limit is $0$, and $0$ is not in $A$.  The sequence walked to the edge of the set and found the door open.

\newthought{In the set $B$}, take $b_n = n$.  Every term belongs to $B$, but this time there is no limit at all, and no subsequence has one either.  The sequence did not leak out through an edge.  It ran away down an endless corridor.

\newthought{In the set $K$}, neither escape is possible.  Take \emph{any} sequence in $K$.  It is bounded, so the Bolzano--Weierstrass Theorem hands us a convergent subsequence, and since every term lies between $0$ and $1$ the Order Limit Theorem says the limit does too.  Whatever the sequence does, some part of it settles down, and it settles down \emph{inside} $K$.

\begin{figure}
\centering
\begin{tikzpicture}[x=1cm,y=1cm]
  % --- A = (0,1]
  \node[lbl,anchor=east] at (1.0,0) {$A=(0,1]$};
  \draw[numline] (1.2,0) -- (9.6,0);
  \draw[setline] (1.8,0) -- (7.8,0);
  \foreach \n in {1,2,3,4,5,6,8,11,16,25,45}
    \node[seqpt] at ({1.8+6/\n},0.28) {};
  \node[openpt] at (1.8,0) {};
  \node[closedpt] at (7.8,0) {};
  \node[lbl,below=3pt] at (1.8,0) {$0$};
  \node[lbl,below=3pt] at (7.8,0) {$1$};
  \draw[->,black!60,line width=0.4pt] (3.4,0.6) -- (2.0,0.6);
  % --- B = [0,infinity)
  \node[lbl,anchor=east] at (1.0,-1.6) {$B=[0,\infty)$};
  \draw[numline] (1.2,-1.6) -- (9.6,-1.6);
  \draw[setline] (1.8,-1.6) -- (9.45,-1.6);
  \foreach \n in {1,...,6}
    \node[seqpt] at ({1.8+1.2*\n},-1.32) {};
  \node[closedpt] at (1.8,-1.6) {};
  \node[lbl,below=3pt] at (1.8,-1.6) {$0$};
  \draw[->,black!60,line width=0.4pt] (7.6,-1.0) -- (9.2,-1.0);
  % --- K = [0,1]
  \node[lbl,anchor=east] at (1.0,-3.2) {$K=[0,1]$};
  \draw[numline] (1.2,-3.2) -- (9.6,-3.2);
  \draw[setline] (1.8,-3.2) -- (7.8,-3.2);
  \foreach \t in {0.93,0.12,0.66,0.30,0.80,0.52,0.40,0.47,0.43,0.455,0.445}
    \node[seqpt] at ({1.8+6*\t},-2.92) {};
  \node[closedpt] at (1.8,-3.2) {};
  \node[closedpt] at (7.8,-3.2) {};
  \node[lbl,below=3pt] at (1.8,-3.2) {$0$};
  \node[lbl,below=3pt] at (7.8,-3.2) {$1$};
\end{tikzpicture}
\caption{Three sets and a sequence (black dots) in each.  In $A$ the sequence leaks out through the missing endpoint, in $B$ it runs off to infinity, and in $K$ it has nowhere to go.  Throughout this book a filled endpoint belongs to the set and a hollow one does not.}
\label{fig:three-escapes}
\end{figure}

Three separate ideas are tangled together in this little experiment, and the work of the chapter is to pull them apart.
\begin{itemize}
	\item A set with no leaks at its edges, like $B$ and $K$, is called \textsc{closed}.
	\item A set with no corridor to infinity, like $A$ and $K$, is \textsc{bounded}, a word we know already.
	\item A set like $K$, from which there is no escape of either kind, will be called \textsc{compact}.  Compactness is the most important idea in the chapter, and one of the most important in all of analysis.
\end{itemize}
Before any of that, though, we need a name for the opposite of having an edge.

\section{Open Sets}

Recall from Chapter Two that the \textsc{$\varepsilon$-neighborhood} of a real number $a$ is the set of points within $\varepsilon$ of it,
\[ V_{\varepsilon}(a) = \{x \in \mathbb{R} : |x - a| < \varepsilon\} = (a - \varepsilon,\, a + \varepsilon). \]

\begin{definition}\label{def:open}
A set $O \subseteq \mathbb{R}$ is \textsc{open} if for every point $a \in O$ there exists an $\varepsilon > 0$ such that $V_{\varepsilon}(a) \subseteq O$.
\end{definition}

In words, every point of an open set is surrounded by a little cushion of other points of the set.  You can stand at any point of $O$ and take a small step in either direction without leaving.  \marginnote{The cushion is allowed to depend on the point.  Near the edge of the set it usually has to be very thin, and that is fine.  It only has to exist.}  An open set has no edge that belongs to it.

\begin{example}
Every open interval $(c,d)$ is an open set.  To see this, let $a \in (c,d)$ and take
\[ \varepsilon = \min\{a - c,\; d - a\}, \]
the distance from $a$ to the nearer endpoint.  If $x \in V_{\varepsilon}(a)$ then $x > a - \varepsilon \geq a - (a - c) = c$ and $x < a + \varepsilon \leq a + (d - a) = d$, so $x \in (c,d)$.
\end{example}

\begin{marginfigure}
\centering
\begin{tikzpicture}[x=1cm,y=1cm]
  \draw[numline] (0,0) -- (4.8,0);
  \draw[setline] (0.6,0) -- (4.2,0);
  \node[openpt] at (0.6,0) {};
  \node[openpt] at (4.2,0) {};
  \draw[nbhdline] (2.6,0.32) -- (4.2,0.32);
  \node[openpt=calm] at (2.6,0.32) {};
  \node[openpt=calm] at (4.2,0.32) {};
  \node[seqpt] at (3.4,0) {};
  \node[lbl,below=3pt] at (0.6,0) {$c$};
  \node[lbl,below=3pt] at (3.4,0) {$a$};
  \node[lbl,below=3pt] at (4.2,0) {$d$};
  \node[lbl,above=2pt,calm] at (3.4,0.32) {$V_{\varepsilon}(a)$};
\end{tikzpicture}
\caption{The cushion around a point $a$ of $(c,d)$.  Its radius is the distance to the nearer endpoint.}
\end{marginfigure}

\begin{example}
The half-open interval $[c,d)$ is \emph{not} open, and the trouble is at exactly one point.  The point $c$ belongs to the set, but it has no cushion.  However small $\varepsilon$ is, the neighborhood $V_{\varepsilon}(c)$ contains $c - \varepsilon/2$, which lies outside.  One bad point is enough to spoil openness.
\end{example}

\begin{example}
The whole line $\mathbb{R}$ is open.  So is the empty set $\emptyset$, because the definition asks something of every point of $O$, and when there are no points there is nothing to check.\marginnote{A statement of the form ``every element of $\emptyset$ has such-and-such a property'' is always true.  Logicians call it \emph{vacuously} true.}  On the other hand a set consisting of a single point is not open, and neither is $\mathbb{Q}$, because every neighborhood of a rational number contains irrational numbers.
\end{example}

How do open sets behave when we combine them?

\begin{theorem}\label{thm:open-unions}
Open sets combine as follows.
\begin{enumerate}[(i)]
	\item The union of an arbitrary collection of open sets is open.
	\item The intersection of a \emph{finite} collection of open sets is open.
\end{enumerate}
\end{theorem}

\begin{proof}
(i)  Let $\{O_{\lambda} : \lambda \in \Lambda\}$ be a collection of open sets and let $O$ be their union.\marginnote{The index set $\Lambda$ may be finite, countable or uncountable.  Part (i) does not care.}  If $a \in O$ then $a$ belongs to at least one of the sets, say $O_{\lambda'}$.  Because $O_{\lambda'}$ is open there is an $\varepsilon > 0$ with $V_{\varepsilon}(a) \subseteq O_{\lambda'}$, and $O_{\lambda'} \subseteq O$.  So $V_{\varepsilon}(a) \subseteq O$.

(ii)  Let $O_1, O_2, \ldots, O_N$ be open and let $a$ be a point of their intersection.  For each $k$ the point $a$ lies in the open set $O_k$, so there is an $\varepsilon_k > 0$ with $V_{\varepsilon_k}(a) \subseteq O_k$.  We need one cushion that fits inside all $N$ sets at once, so we take the thinnest:
\[ \varepsilon = \min\{\varepsilon_1, \varepsilon_2, \ldots, \varepsilon_N\}. \]
Then $\varepsilon > 0$, and $V_{\varepsilon}(a) \subseteq V_{\varepsilon_k}(a) \subseteq O_k$ for every $k$.  Hence $V_{\varepsilon}(a) \subseteq O_1 \cap \cdots \cap O_N$.
\end{proof}

Where exactly did the proof of (ii) use the word \emph{finite}?  In one place only, when we said that the smallest of finitely many positive numbers is positive.  For infinitely many positive numbers there may be no smallest, and their infimum can perfectly well be zero.

\begin{dangerousbend}
An infinite intersection of open sets need not be open.  Take $O_n = \left(-\frac{1}{n}, \frac{1}{n}\right)$ for $n = 1, 2, 3, \ldots$.  Each $O_n$ is open, but
\[ \bigcap_{n=1}^{\infty} O_n = \{0\}, \]
because $0$ lies in every $O_n$, while any $x \neq 0$ is excluded as soon as $\frac{1}{n} < |x|$.  The cushions around $0$ have been squeezed down to nothing.
\end{dangerousbend}

\begin{marginfigure}
\centering
\begin{tikzpicture}[x=1cm,y=1cm]
  \foreach \n in {1,...,6} {
    \pgfmathsetmacro{\w}{2.2/\n}
    \pgfmathsetmacro{\y}{-0.36*(\n-1)}
    \draw[nbhdline] ({2.4-\w},\y) -- ({2.4+\w},\y);
    \node[openpt=calm] at ({2.4-\w},\y) {};
    \node[openpt=calm] at ({2.4+\w},\y) {};
  }
  \node[lbl] at (2.4,-2.2) {$\vdots$};
  \node[closedpt=calm] at (2.4,-2.8) {};
  \node[lbl,right=4pt] at (2.4,-2.8) {$\{0\}$};
  \node[lbl,above=3pt] at (2.4,0) {$O_1 = (-1,1)$};
  \draw[black!40, dashed, line width=0.4pt] (2.4,0) -- (2.4,-1.9);
\end{tikzpicture}
\caption{The open intervals $\left(-\frac1n,\frac1n\right)$ shrink onto the single point $0$.  All that survives in the intersection is a set that is not open.}
\end{marginfigure}

\section{Limit Points and Closed Sets}

In the experiment of Section~\ref{sec:three-escapes}, the point $0$ was not a member of $A = (0,1]$ and yet it was plainly attached to $A$ in some way, because points of $A$ crowd right up against it.  The next definition gives that attachment a name.

\begin{definition}\label{def:limit-point}
A point $x \in \mathbb{R}$ is a \textsc{limit point} of a set $A$ if every $\varepsilon$-neighborhood $V_{\varepsilon}(x)$ contains a point of $A$ other than $x$ itself.
\end{definition}

Two things about this definition deserve a second look.  First, $x$ is not required to belong to $A$.  It may or may not.  Second, the phrase ``other than $x$ itself'' is doing real work.  It means that $x$ does not get to vouch for itself.  There have to be \emph{other} points of $A$ arbitrarily close by.  A limit point of $A$ is a point you can sneak up on from within $A$.\marginnote{Limit points are also called \emph{cluster points} or \emph{accumulation points}, which are perhaps better names, since these are the places where the set piles up.}

Sneaking up on a point is what sequences do, so it is no surprise that limit points can be described by sequences.

\begin{theorem}\label{thm:limit-point-sequence}
A point $x$ is a limit point of a set $A$ if and only if $x = \lim a_n$ for some sequence $(a_n)$ of points of $A$ with $a_n \neq x$ for all $n \in \mathbb{N}$.
\end{theorem}

\begin{proof}
($\Rightarrow$)  Suppose $x$ is a limit point of $A$.  We use the definition over and over, with smaller and smaller neighborhoods.  For each $n \in \mathbb{N}$ the neighborhood $V_{1/n}(x)$ contains a point of $A$ other than $x$.  Choose one and call it $a_n$.  Then $a_n \in A$, $a_n \neq x$, and $|a_n - x| < \frac{1}{n}$.  Given $\varepsilon > 0$, choose $N > 1/\varepsilon$.  For $n \geq N$ we have $|a_n - x| < \frac{1}{n} \leq \frac{1}{N} < \varepsilon$, so $(a_n) \to x$.

($\Leftarrow$)  Suppose $(a_n) \to x$ with $a_n \in A$ and $a_n \neq x$.  Let $\varepsilon > 0$.  By the definition of convergence there is an $N$ with $a_N \in V_{\varepsilon}(x)$, and $a_N$ is a point of $A$ different from $x$.  So every neighborhood of $x$ contains such a point, and $x$ is a limit point of $A$.
\end{proof}

\begin{definition}
A point $a \in A$ is an \textsc{isolated point} of $A$ if it is not a limit point of $A$, which is to say that some neighborhood $V_{\varepsilon}(a)$ contains no point of $A$ except $a$.
\end{definition}

\marginnote{Every point \emph{of} $A$ is either a limit point of $A$ or an isolated point of $A$.  But a limit point of $A$ may lie outside $A$, whereas an isolated point, by definition, never does.}

\begin{example}
Let $A = \left\{\frac{1}{n} : n \in \mathbb{N}\right\} = \left\{1, \frac12, \frac13, \frac14, \ldots\right\}$.

\emph{Every point of $A$ is isolated.}  The nearest neighbor of $\frac1n$ in $A$ is $\frac{1}{n+1}$, at distance $\frac{1}{n(n+1)}$.  A neighborhood of $\frac1n$ with that radius contains no other point of $A$.

\emph{The point $0$ is a limit point of $A$.}  Given $\varepsilon > 0$, the Archimedean Property supplies an $n$ with $\frac1n < \varepsilon$, and $\frac1n$ is a point of $A$, different from $0$, inside $V_{\varepsilon}(0)$.

\emph{There are no other limit points.}  A point $y < 0$ or $y > 1$ has a neighborhood that misses $A$ entirely.  If $y > 0$, only finitely many points of $A$ are larger than $y/2$ (those with $n < 2/y$), so by choosing $\varepsilon$ smaller than $y/2$ and smaller than the distance from $y$ to each of those finitely many points other than $y$ itself, we get a neighborhood of $y$ containing no point of $A$ except possibly $y$.

So this set $A$ has exactly one limit point, and it is not a member of $A$.
\end{example}

\begin{marginfigure}
\centering
\begin{tikzpicture}[x=1cm,y=1cm]
  \draw[numline] (0,0) -- (4.8,0);
  \draw[nbhdline] (0.15,0.3) -- (1.15,0.3);
  \node[openpt=calm] at (0.15,0.3) {};
  \node[openpt=calm] at (1.15,0.3) {};
  \foreach \n in {1,...,40}
    \node[circle, fill=accent, inner sep=0pt, minimum size=2.4pt] at ({0.65+3.8/\n},0) {};
  \node[openpt] at (0.65,0) {};
  \node[lbl,below=3pt] at (0.65,0) {$0$};
  \node[lbl,below=3pt] at (4.45,0) {$1$};
  \node[lbl,below=3pt] at (2.55,0) {$\frac12$};
  \node[lbl,below=3pt] at (1.92,0) {$\frac13$};
  \node[lbl,above=2pt,calm] at (0.65,0.3) {$V_{\varepsilon}(0)$};
\end{tikzpicture}
\caption{The set $\left\{\frac1n\right\}$.  Each of its points stands alone, yet every neighborhood of $0$ catches infinitely many of them.}
\end{marginfigure}

\begin{example}
Every real number is a limit point of $\mathbb{Q}$.  This is just the density of $\mathbb{Q}$ in $\mathbb{R}$ in new clothes, since between $x$ and $x + \varepsilon$ there is always a rational number.  In the same way every real number is a limit point of the set of irrationals.
\end{example}

\newthought{We can now say} what it means for a set to have no leaks.

\begin{definition}\label{def:closed}
A set $F \subseteq \mathbb{R}$ is \textsc{closed} if it contains all of its limit points.
\end{definition}

If you can sneak up on a point from inside a closed set, then that point was inside all along.  The following theorem says the same thing in the language of sequences, and it is the form in which closedness is most often used.

\begin{theorem}\label{thm:closed-sequences}
A set $F \subseteq \mathbb{R}$ is closed if and only if every convergent sequence of points of $F$ has its limit in $F$.
\end{theorem}
\marginnote{In $\mathbb{R}$ a sequence converges exactly when it is Cauchy.  So we could also say that $F$ is closed if and only if every Cauchy sequence in $F$ has a limit in $F$.  A closed subset of $\mathbb{R}$ inherits the completeness of $\mathbb{R}$.}

\begin{proof}
($\Rightarrow$)  Let $F$ be closed, and let $(x_n) \to x$ with $x_n \in F$ for all $n$.  If $x_n = x$ for some $n$, then $x \in F$ and we are done.  Otherwise $x_n \neq x$ for every $n$, and Theorem~\ref{thm:limit-point-sequence} says that $x$ is a limit point of $F$.  Since $F$ is closed, $x \in F$.

($\Leftarrow$)  Suppose every convergent sequence in $F$ has its limit in $F$, and let $x$ be a limit point of $F$.  By Theorem~\ref{thm:limit-point-sequence} there is a sequence in $F$ converging to $x$, and so $x \in F$.  Thus $F$ contains all its limit points.
\end{proof}


Table~\ref{tab:open-closed} collects a number of examples.  It is worth covering the right-hand columns and working each row out for yourself.

\begin{table}
\centering
\begin{tabular}{llcc}
\toprule
Set & Its limit points & Closed? & Open? \\
\midrule
$(0,1)$ & $[0,1]$ & no & yes \\
$[0,1]$ & $[0,1]$ & yes & no \\
$(0,1]$ & $[0,1]$ & no & no \\
$\left\{\frac1n : n \in \mathbb{N}\right\}$ & $\{0\}$ & no & no \\
$\left\{\frac1n : n \in \mathbb{N}\right\} \cup \{0\}$ & $\{0\}$ & yes & no \\
$\mathbb{Z}$ & none & yes & no \\
$\mathbb{Q}$ & all of $\mathbb{R}$ & no & no \\
$\mathbb{R}$ & all of $\mathbb{R}$ & yes & yes \\
$\emptyset$ & none & yes & yes \\
\bottomrule
\end{tabular}
\caption{Open and closed are not opposites.}
\label{tab:open-closed}
\end{table}

\begin{dangerousbend}
A set is not a door.  A door must be either open or closed, but the table shows sets that are neither, such as $(0,1]$ and $\mathbb{Q}$, and two sets, $\mathbb{R}$ and $\emptyset$, that are both.  To show that a set is closed it is never enough to show that it is not open.
\end{dangerousbend}

\subsection{Closure}

If a set fails to be closed, we can repair it by throwing in the points it is missing.

\begin{definition}
Let $A \subseteq \mathbb{R}$ and let $L$ be the set of all limit points of $A$.  The \textsc{closure} of $A$ is the set $\overline{A} = A \cup L$.
\end{definition}

For example $\overline{(0,1)} = [0,1]$ and $\overline{\mathbb{Q}} = \mathbb{R}$, while $\overline{\mathbb{Z}} = \mathbb{Z}$.  It seems clear that $\overline{A}$ ought to be closed, but there is something to check.  By adding the limit points of $A$ we have made a bigger set, and a bigger set might have acquired new limit points of its own.  Happily it has not.

\begin{theorem}\label{thm:closure}
For any $A \subseteq \mathbb{R}$, the closure $\overline{A}$ is a closed set, and it is the smallest closed set containing $A$.
\end{theorem}

\begin{proof}
Let $L$ be the set of limit points of $A$, and let $x$ be a limit point of $\overline{A} = A \cup L$.  We will show that $x$ is in fact a limit point of $A$, so that $x \in L \subseteq \overline{A}$.

Let $\varepsilon > 0$.  The neighborhood $V_{\varepsilon}(x)$ contains some point $y \in \overline{A}$ with $y \neq x$.  If $y \in A$, we have what we want, a point of $A$ other than $x$ inside $V_{\varepsilon}(x)$.  If not, then $y \in L$, and we use the fact that $y$ is itself a limit point of $A$.  Put a small neighborhood $V_{\delta}(y)$ around $y$, small enough to fit inside $V_{\varepsilon}(x)$ and to leave out $x$.  The choice
\[ \delta = \min\{\varepsilon - |x - y|,\; |x - y|\} \]
does both.  Since $y$ is a limit point of $A$, the neighborhood $V_{\delta}(y)$ contains a point $a \in A$.  This $a$ lies in $V_{\varepsilon}(x)$ and is different from $x$.  In either case, then, $V_{\varepsilon}(x)$ contains a point of $A$ other than $x$, and $x$ is a limit point of $A$.

For the second claim, let $F$ be any closed set with $A \subseteq F$.  A limit point of $A$ is automatically a limit point of the larger set $F$, and $F$ contains its limit points.  Hence $L \subseteq F$, and so $\overline{A} = A \cup L \subseteq F$.
\end{proof}

\begin{marginfigure}
\centering
\begin{tikzpicture}[x=1cm,y=1cm]
  \draw[numline] (0,0) -- (4.8,0);
  \draw[bracedeco] (0.5,0.62) -- (3.9,0.62) node[midway,above=4pt,lbl] {$V_{\varepsilon}(x)$};
  \draw[bracedeco,draw=accent] (3.9,-0.3) -- (2.5,-0.3) node[midway,below=4pt,lbl,accent] {$V_{\delta}(y)$};
  \draw[nbhdline] (0.5,0) -- (3.9,0);
  \draw[setline] (2.5,0) -- (3.9,0);
  \node[seqpt] at (2.0,0) {};
  \node[seqpt] at (3.2,0) {};
  \node[seqpt] at (2.75,0) {};
  \node[lbl,below=3pt] at (2.0,0) {$x$};
  \node[lbl,above=3pt] at (3.2,0) {$y$};
  \node[lbl,above=3pt] at (2.75,0) {$a$};
\end{tikzpicture}
\caption{A neighborhood inside a neighborhood.  The point $y$ is close to $x$, and $a \in A$ is close to $y$, so $a$ is close to $x$.}
\end{marginfigure}

\subsection{Complements}

Open sets and closed sets are not opposites, as the table showed.  The true relationship between them is more interesting.  Write $A^c = \{x \in \mathbb{R} : x \notin A\}$ for the \textsc{complement} of a set $A$.

\begin{theorem}\label{thm:complements}
A set $O$ is open if and only if its complement $O^c$ is closed.  Likewise, a set $F$ is closed if and only if $F^c$ is open.
\end{theorem}

\begin{proof}
($\Rightarrow$)  Let $O$ be open and let $x$ be a limit point of $O^c$.  We must show that $x \in O^c$.  If instead $x$ were in $O$, there would be a neighborhood $V_{\varepsilon}(x)$ lying entirely inside $O$.  That neighborhood would contain no points of $O^c$ at all, which is impossible for a limit point of $O^c$.  So $x \in O^c$, and $O^c$ is closed.

($\Leftarrow$)  Let $O^c$ be closed and let $x \in O$.  Then $x$ is not in $O^c$, and since $O^c$ contains all of its limit points, $x$ is not a limit point of $O^c$.  So some neighborhood $V_{\varepsilon}(x)$ contains no point of $O^c$ other than $x$, and $x$ itself is not in $O^c$ either.  That is, $V_{\varepsilon}(x) \subseteq O$.  Hence $O$ is open.

The second statement is the first one applied to $O = F^c$, since $(F^c)^c = F$.
\end{proof}

This theorem lets us translate any fact about open sets into a fact about closed sets by taking complements.  The dictionary is supplied by De Morgan's Laws.\marginnote{\textsc{De Morgan's Laws:}
\[ \Big(\bigcup_{\lambda} E_{\lambda}\Big)^{c} = \bigcap_{\lambda} E_{\lambda}^{c}, \]
\[ \Big(\bigcap_{\lambda} E_{\lambda}\Big)^{c} = \bigcup_{\lambda} E_{\lambda}^{c}. \]
A point fails to be in \emph{some} set exactly when it is not in \emph{all} of them.}

\begin{corollary}\label{cor:closed-unions}
Closed sets combine as follows.
\begin{enumerate}[(i)]
	\item The union of a \emph{finite} collection of closed sets is closed.
	\item The intersection of an arbitrary collection of closed sets is closed.
\end{enumerate}
\end{corollary}

\begin{proof}
(i)  If $F_1, \ldots, F_N$ are closed then each $F_k^c$ is open, and by De Morgan's Laws $(F_1 \cup \cdots \cup F_N)^c = F_1^c \cap \cdots \cap F_N^c$.  This is a finite intersection of open sets, hence open by Theorem~\ref{thm:open-unions}.  So $F_1 \cup \cdots \cup F_N$ is closed.  Part (ii) is proved in the same way from part (i) of Theorem~\ref{thm:open-unions}.
\end{proof}

Notice that the roles of ``finite'' and ``arbitrary'' have changed places.  An infinite union of closed sets can fail to be closed.  Each interval $\left[\frac1n, 1\right]$ is closed, but
\[ \bigcup_{n=1}^{\infty} \left[\tfrac1n, 1\right] = (0, 1]. \]

\section{Compact Sets}

We return to the set $K = [0,1]$ of Section~\ref{sec:three-escapes}, the set from which no sequence could escape.  Here is the definition that captures what was special about it.

\begin{definition}\label{def:compact}
A set $K \subseteq \mathbb{R}$ is \textsc{compact} if every sequence in $K$ has a subsequence that converges to a limit that is also in $K$.
\end{definition}

\begin{example}
Every closed interval $[c,d]$ is compact, by the argument we gave for $[0,1]$.  The Bolzano--Weierstrass Theorem produces a convergent subsequence, and the Order Limit Theorem keeps its limit between $c$ and $d$.

The set $(0,1]$ is not compact.  The sequence $\left(\frac1n\right)$ lies in it, but every subsequence converges to $0$, which does not.

The set $[0,\infty)$ is not compact.  The sequence $(n)$ lies in it and has no convergent subsequence whatever.
\end{example}

The two failures in this example are the two escapes of Figure~\ref{fig:three-escapes}, and the next theorem says that they are the only two ways for compactness to fail.

\begin{theorem}[Characterization of Compactness in $\mathbb{R}$]\label{thm:closed-bounded}
A set $K \subseteq \mathbb{R}$ is compact if and only if it is closed and bounded.
\end{theorem}

\begin{proof}
($\Rightarrow$)  Let $K$ be compact.

\emph{$K$ is bounded.}  Suppose not.  Then for each $n \in \mathbb{N}$ there is a point $x_n \in K$ with $|x_n| > n$.  Every subsequence of $(x_n)$ is unbounded, and convergent sequences are bounded, so no subsequence of $(x_n)$ converges.  This contradicts compactness.

\emph{$K$ is closed.}  Let $x$ be a limit point of $K$.  By Theorem~\ref{thm:limit-point-sequence} there is a sequence $(x_n)$ in $K$ with $(x_n) \to x$.  Compactness gives a subsequence converging to some point of $K$.  But every subsequence of a convergent sequence converges to the same limit as the whole sequence, so that point is $x$.  Hence $x \in K$.

($\Leftarrow$)  Let $K$ be closed and bounded, and let $(x_n)$ be a sequence in $K$.  It is a bounded sequence, so by Bolzano--Weierstrass it has a convergent subsequence.  The limit of that subsequence is the limit of a convergent sequence of points of the closed set $K$, so by Theorem~\ref{thm:closed-sequences} it belongs to $K$.
\end{proof}

\marginnote{``Closed and bounded'' is the practical test for compactness in $\mathbb{R}$.  The definition by sequences is what compactness \emph{means}, and this theorem is how one usually \emph{checks} it.}

Compact sets are the right generalization of closed intervals.  As evidence, here is the Nested Interval Property again, with ``closed interval'' replaced by ``compact set.''

\begin{theorem}[Nested Compact Set Property]\label{thm:nested-compact}
If
\[ K_1 \supseteq K_2 \supseteq K_3 \supseteq K_4 \supseteq \cdots \]
is a nested sequence of nonempty compact sets, then $\bigcap_{n=1}^{\infty} K_n$ is not empty.
\end{theorem}

\begin{proof}
For each $n$ choose a point $x_n \in K_n$.  Because the sets are nested, the whole sequence $(x_n)$ lies in $K_1$, and $K_1$ is compact.  So there is a subsequence $(x_{n_k})$ converging to some $x \in K_1$.  We claim that $x$ belongs to every $K_m$.

Fix $m$.  Whenever $k \geq m$ we have $n_k \geq k \geq m$, so $x_{n_k} \in K_{n_k} \subseteq K_m$.  Thus, after its first few terms, $(x_{n_k})$ is a sequence in $K_m$ converging to $x$.  Since $K_m$ is closed, $x \in K_m$.
\end{proof}

Both hypotheses are needed.  The sets $\left(0, \frac1n\right)$ are nested, bounded and nonempty, but their intersection is empty.  They are not closed.  The sets $[n, \infty)$ are nested, closed and nonempty, but their intersection is empty.  They are not bounded.

\subsection{Open Covers}

There is a second way of describing compactness, one that makes no mention of sequences.  It is stranger than the first, and at first sight it seems to have nothing to do with it.  It is also the description that mathematicians reach for most often, so it repays some patience.

Imagine a road, and street lamps along it.  Each lamp lights up an open stretch of the road.  Suppose there are enough lamps, possibly infinitely many, that every point of some set $A$ on the road is lit.  Could we switch off all but \emph{finitely many} of the lamps and still have every point of $A$ lit?

\begin{definition}
Let $A \subseteq \mathbb{R}$.  An \textsc{open cover} of $A$ is a collection $\{O_{\lambda} : \lambda \in \Lambda\}$ of open sets whose union contains $A$.  A \textsc{finite subcover} is a finite subcollection of the original open cover whose union still contains $A$.
\end{definition}

\begin{example}\label{ex:cover-open-interval}
Let $A = (0,1)$, and for $n = 2, 3, 4, \ldots$ let $O_n = \left(\frac1n, 1\right)$.

This is an open cover of $A$.  If $x \in (0,1)$, choose $n > 1/x$.  Then $\frac1n < x < 1$, so $x \in O_n$.

But it has no finite subcover.  Take any finitely many of the sets, $O_{n_1}, \ldots, O_{n_m}$, and let $N$ be the largest of the indices.  The sets are nested, so their union is just $O_N = \left(\frac1N, 1\right)$, and this misses the point $\frac{1}{N}$ of $A$.  However many lamps we leave on, there is a dark stretch of road next to $0$.
\end{example}

\begin{figure}
\centering
\begin{tikzpicture}[x=1cm,y=1cm]
  \draw[numline] (0.4,0) -- (9.7,0);
  \draw[setline] (1,0) -- (9,0);
  \node[openpt] at (1,0) {};
  \node[openpt] at (9,0) {};
  \node[lbl,above=3pt] at (1,0) {$0$};
  \node[lbl,above=3pt] at (9,0) {$1$};
  \node[lbl,above=3pt,accent] at (5,0) {$A = (0,1)$};
  \foreach \n in {2,...,8} {
    \pgfmathsetmacro{\y}{-0.42*(\n-1)}
    \pgfmathsetmacro{\l}{1+8/\n}
    \draw[nbhdline] (\l,\y) -- (9,\y);
    \node[openpt=calm] at (\l,\y) {};
    \node[openpt=calm] at (9,\y) {};
    \node[lbl,anchor=west] at (9.25,\y) {$O_{\n}$};
  }
  \node[lbl] at (5,-3.5) {$\vdots$};
  \draw[black!45, dashed, line width=0.4pt] (1,0) -- (1,-3.7);
\end{tikzpicture}
\caption{The open cover $O_n = \left(\frac1n, 1\right)$ of the interval $(0,1)$.  The sets creep toward $0$ without ever arriving, so no finite number of them will do.}
\label{fig:open-cover}
\end{figure}

\begin{example}\label{ex:cover-closed-interval}
Now try the same thing with the closed interval $[0,1]$.  The sets $O_n$ no longer form a cover, because nothing covers $0$ or $1$.  To get a cover we must add at least one open set $O_0$ containing $0$ and one open set $O_1'$ containing $1$.  But look at what $O_0$ is forced to do.  It is open, so it contains a whole neighborhood $V_{\varepsilon}(0)$, and that neighborhood swallows the dark stretch of road, since it covers all of $[0, \varepsilon)$.  Now choose $n$ with $\frac1n < \varepsilon$.  The three sets $O_0$, $O_n$ and $O_1'$ already cover $[0,1]$.

Adding the two endpoints to the set changed everything.  Whatever covers the point $0$ must cover its neighbors too, and only finitely many more sets are then needed.
\end{example}

That is the phenomenon in miniature.  Here is the theorem.

\begin{theorem}[Heine--Borel Theorem]\label{thm:heine-borel}
For a set $K \subseteq \mathbb{R}$, the following three statements are equivalent.
\begin{enumerate}[(i)]
	\item $K$ is compact, meaning that every sequence in $K$ has a subsequence converging to a point of $K$.
	\item $K$ is closed and bounded.
	\item Every open cover of $K$ has a finite subcover.
\end{enumerate}
\end{theorem}

\begin{proof}
We know from Theorem~\ref{thm:closed-bounded} that (i) and (ii) are equivalent, so it is enough to prove that (iii) implies (ii) and that (ii) implies (iii).

\newthought{(iii) implies (ii).}  Assume every open cover of $K$ has a finite subcover.  The trick in both halves is to design an open cover whose finite subcovers tell us what we want to know.

\emph{$K$ is bounded.}  The open intervals $O_n = (-n, n)$, $n \in \mathbb{N}$, cover all of $\mathbb{R}$ and so certainly cover $K$.  Finitely many of them cover $K$, and if $N$ is the largest index used, then $K \subseteq (-N, N)$.

\emph{$K$ is closed.}  Let $y$ be a limit point of $K$ and suppose, hoping for a contradiction, that $y \notin K$.  For each $n \in \mathbb{N}$ let
\[ O_n = \left\{x \in \mathbb{R} : |x - y| > \tfrac1n\right\} = \left(-\infty,\, y - \tfrac1n\right) \cup \left(y + \tfrac1n,\, \infty\right), \]
an open set.  The union of all the $O_n$ is every real number except $y$, and $y \notin K$, so these sets cover $K$.  Finitely many of them cover $K$, and since $O_1 \subseteq O_2 \subseteq \cdots$ this means $K \subseteq O_N$ for a single $N$.  But then no point of $K$ lies within $\frac1N$ of $y$, and $y$ is not a limit point of $K$ after all.  This contradiction shows that $y \in K$.

\newthought{(ii) implies (iii).}  This is the heart of the theorem.  Let $K$ be closed and bounded, let $\{O_{\lambda} : \lambda \in \Lambda\}$ be an open cover of $K$, and suppose for contradiction that no finite subcollection covers $K$.

Since $K$ is bounded, it lies in some closed interval $I_0 = [-M, M]$.  Cut $I_0$ in half.  If the part of $K$ in the left half could be covered by finitely many of the $O_{\lambda}$, and likewise the part in the right half, then putting the two finite collections together would cover $K$.  So for at least one of the halves, call it $I_1$, the set $K \cap I_1$ cannot be covered by finitely many of the $O_{\lambda}$.  Cut $I_1$ in half and repeat.  We obtain closed intervals
\[ I_0 \supseteq I_1 \supseteq I_2 \supseteq \cdots, \]
where $I_n$ has length $2M/2^n$ and $K \cap I_n$ cannot be covered by finitely many of the $O_{\lambda}$.  In particular $K \cap I_n$ is never empty.

By the Nested Interval Property there is a point $x$ in every $I_n$.  We claim $x \in K$.  Choose $x_n \in K \cap I_n$ for each $n$.  Since $x$ and $x_n$ both lie in $I_n$, we have $|x_n - x| \leq 2M/2^n$, so $(x_n) \to x$.  Since $K$ is closed, $x \in K$ by Theorem~\ref{thm:closed-sequences}.

Now the cover comes back into play.  The point $x$ is in $K$, so it lies in one of the open sets, say $O_{\lambda_0}$, and therefore so does a whole neighborhood $V_{\varepsilon}(x)$.  Choose $n$ large enough that $2M/2^n < \varepsilon$.  The interval $I_n$ contains $x$ and is shorter than $\varepsilon$, so
\[ I_n \subseteq V_{\varepsilon}(x) \subseteq O_{\lambda_0}. \]
But this says that $K \cap I_n$ is covered by the \emph{single} set $O_{\lambda_0}$, when we built $I_n$ so that no finite number of the $O_{\lambda}$ would do.  This contradiction completes the proof.
\end{proof}

\begin{marginfigure}
\centering
\begin{tikzpicture}[x=1cm,y=1cm]
  \draw[numline] (0,0) -- (4.8,0);
  \draw[bracedeco] (0.3,0.62) -- (4.5,0.62) node[midway,above=4pt,lbl] {$O_{\lambda_0}$};
  \draw[nbhdline] (1.1,0.3) -- (3.7,0.3);
  \node[openpt=calm] at (1.1,0.3) {};
  \node[openpt=calm] at (3.7,0.3) {};
  \node[lbl,calm,anchor=west] at (3.8,0.3) {$V_{\varepsilon}(x)$};
  \draw[setline] (1.9,0) -- (2.9,0);
  \node[closedpt] at (1.9,0) {};
  \node[closedpt] at (2.9,0) {};
  \node[seqpt] at (2.4,0) {};
  \node[lbl,below=3pt] at (2.4,0) {$x$};
  \draw[bracedeco,draw=accent] (2.9,-0.45) -- (1.9,-0.45) node[midway,below=4pt,lbl,accent] {$I_n$};
\end{tikzpicture}
\caption{The end of the proof.  The intervals $I_n$ shrink onto $x$, so eventually one of them fits inside the neighborhood of $x$ that a single set of the cover provides.}
\end{marginfigure}

Compare this proof with the proof of the Bolzano--Weierstrass Theorem in Chapter Two.  They are the same argument.  Bisect, keep the ``bad'' half, and let the Nested Interval Property find the point where the trouble is concentrated.  And once again, underneath everything, is the Axiom of Completeness.

\subsection{What Is Compactness For?}

The open-cover description may seem like a peculiar way of saying ``closed and bounded.''  Its value is that it converts \emph{local} information into \emph{global} information.  Very often in analysis we know something about a small neighborhood of each point.  Near each point a function is bounded, or two quantities are close, or some estimate holds.  Each point has its own neighborhood and its own estimate, and there are infinitely many points.  If the set is compact, finitely many of those neighborhoods already cover it, and finitely many estimates can always be combined into one.

\begin{example}
Let $K$ be compact and let $f : K \to \mathbb{R}$ be a function that is \emph{locally bounded}, meaning that for each $x \in K$ there are a neighborhood $V_{\varepsilon_x}(x)$ and a number $M_x$ such that $|f(t)| \leq M_x$ for all $t \in K \cap V_{\varepsilon_x}(x)$.  Then $f$ is bounded on all of $K$.

Indeed, the neighborhoods $V_{\varepsilon_x}(x)$, one for each $x \in K$, form an open cover of $K$.  By the Heine--Borel Theorem, finitely many of them, centered at $x_1, \ldots, x_m$ say, cover $K$.  Put $M = \max\{M_{x_1}, \ldots, M_{x_m}\}$.  Every $t \in K$ lies in one of the $m$ neighborhoods, so $|f(t)| \leq M$.

Without compactness the conclusion is false.  The function $f(x) = 1/x$ on $(0,1)$ is locally bounded, since near any point $x$ it is at most $2/x$, but it is certainly not bounded.
\end{example}

We will lean on this local-to-global principle again and again in the chapters on continuity and integration.

\section{The Cantor Set}

Our intuition about subsets of the line is built on intervals, and intervals are misleadingly tame.  The following set, constructed by Georg Cantor in 1883, is a good corrective.  It is a compact set that manages to be very small and very large at the same time.

Start with the closed interval $C_0 = [0,1]$ and remove its open middle third, leaving
\[ C_1 = \left[0, \tfrac13\right] \cup \left[\tfrac23, 1\right]. \]
Now remove the open middle third of each of the two pieces, leaving
\[ C_2 = \left[0, \tfrac19\right] \cup \left[\tfrac29, \tfrac13\right] \cup \left[\tfrac23, \tfrac79\right] \cup \left[\tfrac89, 1\right], \]
and carry on in this way forever.  At stage $n$ we have a set $C_n$ consisting of $2^n$ closed intervals, each of length $1/3^n$.  The \textsc{Cantor set} is what remains when all the removals have been made:
\[ C = \bigcap_{n=0}^{\infty} C_n. \]

\begin{figure}
\centering
\begin{tikzpicture}[x=9cm,y=1cm]
  \def\h{0.2}
  \fill[accent] (0,0) rectangle (1,\h);
  \foreach \a in {0,2}
    \fill[accent] ({\a/3},-0.65) rectangle ++(1/3,\h);
  \foreach \a in {0,2} \foreach \b in {0,2}
    \fill[accent] ({\a/3+\b/9},-1.3) rectangle ++(1/9,\h);
  \foreach \a in {0,2} \foreach \b in {0,2} \foreach \c in {0,2}
    \fill[accent] ({\a/3+\b/9+\c/27},-1.95) rectangle ++(1/27,\h);
  \foreach \a in {0,2} \foreach \b in {0,2} \foreach \c in {0,2} \foreach \d in {0,2}
    \fill[accent] ({\a/3+\b/9+\c/27+\d/81},-2.6) rectangle ++(1/81,\h);
  \foreach \k/\y in {0/0,1/-0.65,2/-1.3,3/-1.95,4/-2.6}
    \node[lbl,anchor=east] at (-0.025,{\y+0.1}) {$C_{\k}$};
  \node[lbl] at (0,0.45) {$0$};
  \node[lbl] at (1,0.45) {$1$};
  \node[lbl] at (1/3,-0.24) {$\frac13$};
  \node[lbl] at (2/3,-0.24) {$\frac23$};
  \draw[calm, line width=0.6pt] (0.25,0.4) -- (0.25,-2.8);
  \node[lbl,calm] at (0.25,-3.05) {$\frac14$};
\end{tikzpicture}
\caption{The first five stages in the construction of the Cantor set.  The vertical line marks the number $\frac14$, which is never an endpoint and is never removed.}
\label{fig:cantor}
\end{figure}

\newthought{How much is left?}  One way to answer is to add up what we threw away.  At the first stage we removed one interval of length $\frac13$, at the second stage two intervals of length $\frac19$, and at the $n$th stage $2^{n-1}$ intervals of length $1/3^n$.  The total length removed is
\[ \frac13 + \frac29 + \frac{4}{27} + \cdots = \frac13 \sum_{k=0}^{\infty} \left(\frac23\right)^{k} = \frac13 \cdot \frac{1}{1 - \frac23} = 1. \]
We started with an interval of length $1$ and removed a total length of $1$.  In this sense the Cantor set has length zero.\marginnote{Put another way, $C$ is contained in $C_n$, which has total length $\left(\frac23\right)^n$, and that tends to zero.}  It certainly contains no interval, however short, because an interval of length $\ell$ cannot fit inside $C_n$ once $1/3^n < \ell$.

\newthought{Is anything left at all?}  Yes.  The endpoints $0$ and $1$ are never removed, and neither are $\frac13$, $\frac23$, $\frac19$, $\frac29$, or any other endpoint of any interval of any $C_n$, since we only ever remove \emph{open} middle thirds.  So $C$ contains infinitely many points.  There are countably many of these endpoints, and it would be natural to guess that they are all there is.

\newthought{The guess is wrong.}  To see what else is in $C$, give every point of $C$ an address.  A point $x \in C$ lies in $C_1$, so it is in either the \textbf{L}eft piece or the \textbf{R}ight piece.  Whichever piece that is, it splits into two at the next stage, and $x$ is in the left or the right part of it, and so on forever.  In this way $x$ determines an infinite sequence of L's and R's.  For instance, Figure~\ref{fig:cantor} shows that the address of $\frac14$ begins L, R, L, R, and in fact it continues in the same pattern forever.

Conversely, every infinite sequence of L's and R's is the address of exactly one point of $C$.  Following the instructions picks out a nested sequence of closed intervals $I_1 \supseteq I_2 \supseteq I_3 \supseteq \cdots$, one from each $C_n$, with lengths $1/3^n$.  By the Nested Interval Property they have a point in common.  Because their lengths tend to zero they have only one.  That point lies in every $C_n$, so it lies in $C$.  Two different addresses lead to different points, because at the first place where the addresses disagree the two points fall into disjoint intervals.

So the points of the Cantor set correspond one-to-one with the infinite sequences of two symbols.  We proved in Chapter One that the set of all such sequences is uncountable.  The conclusion is hard to believe at first.

\begin{quote}
\emph{The Cantor set is uncountable.}  It has as many points as the whole interval $[0,1]$ from which it was carved, even though the carving removed the entire length of that interval.
\end{quote}

\marginnote{Readers who know base-3 notation can describe $C$ in one line.  It is the set of numbers in $[0,1]$ that can be written in base 3 using only the digits $0$ and $2$.  An L is a digit $0$ and an R is a digit $2$, so that $\frac14 = 0.020202\ldots$ in base 3.}

\newthought{What does topology say about $C$?}  Each $C_n$ is a finite union of closed intervals and is therefore closed.  So $C$, an intersection of closed sets, is closed (Corollary~\ref{cor:closed-unions}).  It is bounded.  Hence \emph{the Cantor set is compact}.

Moreover, \emph{the Cantor set has no isolated points.}  Let $x \in C$ and $\varepsilon > 0$, and choose $n$ with $1/3^n < \varepsilon$.  The point $x$ lies in one of the intervals that make up $C_n$.  Both endpoints of that interval belong to $C$, both are within $1/3^n < \varepsilon$ of $x$, and at least one of them is different from $x$.  So $V_{\varepsilon}(x)$ contains a point of $C$ other than $x$.

Sets with these two properties have a name.

\begin{definition}
A set $P \subseteq \mathbb{R}$ is \textsc{perfect} if it is closed and has no isolated points.
\end{definition}

Closed intervals $[c,d]$ with $c < d$ are perfect, and so are $\mathbb{R}$ and the Cantor set.  The set $\left\{\frac1n\right\} \cup \{0\}$ is closed but not perfect, since most of its points are isolated.  The set $\mathbb{Q}$ has no isolated points but is not perfect, since it is not closed.  The address argument showed that one particular perfect set is uncountable.  Remarkably, they all are.

\begin{theorem}\label{thm:perfect-uncountable}
A nonempty perfect set is uncountable.
\end{theorem}

\begin{proof}
Let $P$ be perfect and nonempty.  Then $P$ is infinite, because every point of a finite set is isolated.  Suppose, for contradiction, that $P$ is countable, so that its points can be listed:
\[ P = \{x_1, x_2, x_3, \ldots\}. \]
We will build a nested sequence of closed intervals $I_1 \supseteq I_2 \supseteq I_3 \supseteq \cdots$ so that
\begin{enumerate}[(a)]
	\item each $I_n$ contains a point of $P$ in its interior, that is, strictly between its endpoints, and
	\item $x_n \notin I_{n+1}$ for every $n$.
\end{enumerate}
Begin with $I_1 = [x_1 - 1, x_1 + 1]$, which has the point $x_1$ of $P$ in its interior.

Suppose $I_n$ has been constructed, and let $y$ be a point of $P$ in its interior.  We first find a point $y'$ of $P$ in the interior of $I_n$ with $y' \neq x_n$.  If $y \neq x_n$ we simply take $y' = y$.  If $y = x_n$, we use the fact that $y$ is not isolated.  It is a limit point of $P$, so every neighborhood of $y$, including one small enough to lie in the interior of $I_n$, contains a point $y' \in P$ different from $y$.  Now let $I_{n+1}$ be a closed interval centered at $y'$, short enough to lie inside $I_n$ and to leave out $x_n$.  Then (a) holds for $I_{n+1}$, and (b) holds for $n$.

Finally let $K_n = I_n \cap P$.  Each $K_n$ is closed and bounded, hence compact.  It is nonempty by (a), and the sets are nested.  By the Nested Compact Set Property (Theorem~\ref{thm:nested-compact}) there is a point $x$ in every $K_n$.  This $x$ is a point of $P$, so it appears on the list, say as $x = x_m$.  But $x \in I_{m+1}$ and, by (b), $x_m \notin I_{m+1}$.  This contradiction shows that $P$ cannot be countable.
\end{proof}

\begin{marginfigure}
\centering
\begin{tikzpicture}[x=1cm,y=1cm]
  \draw[numline] (0,0) -- (4.8,0);
  \draw[bracedeco] (0.4,0.6) -- (4.4,0.6) node[midway,above=4pt,lbl] {$I_n$};
  \draw[bracedeco,draw=accent] (3.8,-0.5) -- (2.6,-0.5) node[midway,below=4pt,lbl,accent] {$I_{n+1}$};
  \draw[tickmark] (0.4,-0.12) -- (0.4,0.12);
  \draw[tickmark] (4.4,-0.12) -- (4.4,0.12);
  \draw[setline] (2.6,0) -- (3.8,0);
  \node[closedpt] at (2.6,0) {};
  \node[closedpt] at (3.8,0) {};
  \node[seqpt] at (1.5,0) {};
  \node[seqpt] at (3.2,0) {};
  \node[lbl,below=3pt] at (1.5,0) {$x_n$};
  \node[lbl,above=3pt] at (3.2,0) {$y'$};
\end{tikzpicture}
\caption{One step of the construction.  The new interval is centered on a point $y'$ of $P$ and is short enough to avoid $x_n$.}
\end{marginfigure}

Since $\mathbb{R}$ is itself a nonempty perfect set, this gives one more proof that $\mathbb{R}$ is uncountable.  In fact the Nested Interval proof of that fact in Chapter One is precisely this argument in the special case $P = \mathbb{R}$, where finding the point $y'$ takes no effort at all.

\section{Connected Sets}

An interval is all in one piece.  The set $[0,1] \cup [2,3]$ is in two pieces.  It is surprisingly delicate to say precisely what ``in one piece'' means.

As a first attempt, we might call a set disconnected if it can be split into two nonempty disjoint parts.  This is useless, since \emph{every} set with at least two points can be split in that way.  The interval $(1,3)$, for instance, is the union of the disjoint sets $(1,2]$ and $(2,3)$.  What is wrong with this splitting is that the two parts are still touching.  The point $2$ belongs to the first part and is a limit point of the second.  A genuine break in a set should leave the two parts not merely disjoint but out of reach of each other.

\begin{definition}
Two nonempty sets $A, B \subseteq \mathbb{R}$ are \textsc{separated} if $\overline{A} \cap B$ and $A \cap \overline{B}$ are both empty.  In other words, neither set contains a point of the other, or a limit point of the other.

A set $E \subseteq \mathbb{R}$ is \textsc{disconnected} if it can be written as $E = A \cup B$ where $A$ and $B$ are nonempty separated sets.  A set that is not disconnected is \textsc{connected}.
\end{definition}

\begin{marginfigure}
\centering
\begin{tikzpicture}[x=1cm,y=1cm]
  % separated
  \draw[numline] (0,0) -- (4.8,0);
  \draw[setline] (0.6,0) -- (2.4,0);
  \draw[nbhdline] (2.4,0) -- (4.2,0);
  \node[openpt] at (0.6,0) {};
  \node[openpt=calm] at (4.2,0) {};
  \node[openpt=black] at (2.4,0) {};
  \node[lbl,below=3pt] at (0.6,0) {$1$};
  \node[lbl,below=3pt] at (2.4,0) {$2$};
  \node[lbl,below=3pt] at (4.2,0) {$3$};
  \node[lbl,above=3pt,accent] at (1.5,0) {$A$};
  \node[lbl,above=3pt,calm] at (3.3,0) {$B$};
  % not separated
  \draw[numline] (0,-1.5) -- (4.8,-1.5);
  \draw[setline] (0.6,-1.5) -- (2.4,-1.5);
  \draw[nbhdline] (2.4,-1.5) -- (4.2,-1.5);
  \node[openpt] at (0.6,-1.5) {};
  \node[openpt=calm] at (4.2,-1.5) {};
  \node[closedpt] at (2.4,-1.5) {};
  \node[lbl,below=3pt] at (0.6,-1.5) {$1$};
  \node[lbl,below=3pt] at (2.4,-1.5) {$2$};
  \node[lbl,below=3pt] at (4.2,-1.5) {$3$};
  \node[lbl,above=3pt,accent] at (1.5,-1.5) {$A$};
  \node[lbl,above=3pt,calm] at (3.3,-1.5) {$B$};
\end{tikzpicture}
\caption{In the upper picture $A = (1,2)$ and $B = (2,3)$ are separated, and the missing point $2$ keeps them apart.  In the lower picture $A = (1,2]$ and $B = (2,3)$ are disjoint but not separated, because $2 \in A$ is a limit point of $B$.}
\end{marginfigure}

\begin{example}
The sets $A = (1,2)$ and $B = (2,3)$ are separated, since $\overline{A} = [1,2]$ does not meet $B$, and $\overline{B} = [2,3]$ does not meet $A$.  So $(1,2) \cup (2,3)$ is disconnected, as it should be.

The sets $(1,2]$ and $(2,3)$ are not separated, because the closure of the second contains the point $2$ of the first.  This splitting does not show that $(1,3)$ is disconnected, and we will see in a moment that no splitting can.
\end{example}

\begin{example}
The set $\mathbb{Q}$ is disconnected.  Let $A = \mathbb{Q} \cap (-\infty, \sqrt2)$ and $B = \mathbb{Q} \cap (\sqrt2, \infty)$.  Their union is $\mathbb{Q}$, because $\sqrt2$ is irrational.  Since $\overline{A} = (-\infty, \sqrt2]$ contains no point of $B$, and $\overline{B} = [\sqrt2, \infty)$ contains no point of $A$, the two sets are separated.  The rational numbers fall apart at every one of the irrational gaps between them.
\end{example}

On the real line there is a simple test for connectedness.

\begin{theorem}\label{thm:connected-intervals}
A set $E \subseteq \mathbb{R}$ is connected if and only if, whenever $a < c < b$ with $a \in E$ and $b \in E$, it follows that $c \in E$ as well.
\end{theorem}

\begin{proof}
($\Rightarrow$)  We prove the contrapositive.  Suppose there are points $a, b \in E$ and a point $c \notin E$ with $a < c < b$.  Let
\[ A = E \cap (-\infty, c), \qquad B = E \cap (c, \infty). \]
Then $a \in A$ and $b \in B$, so both are nonempty, and $E = A \cup B$ because $c \notin E$.  Every limit point of $A$ is at most $c$, so $\overline{A} \subseteq (-\infty, c]$ does not meet $B$, and similarly $\overline{B}$ does not meet $A$.  Thus $E$ is disconnected.

($\Leftarrow$)  Suppose $E$ has the property in the theorem, and suppose for contradiction that $E = A \cup B$ with $A$ and $B$ nonempty and separated.  Pick $a \in A$ and $b \in B$, where we may assume $a < b$.  By our hypothesis the whole interval $[a,b]$ lies in $E$.  Let
\[ s = \sup\big(A \cap [a,b]\big). \]
The supremum of a set is either a member of the set or a limit point of it, so $s \in \overline{A}$.\marginnote{For every $\varepsilon > 0$ there is an element of the set in $(s - \varepsilon, s]$.  If none of these elements is $s$ itself, then $s$ is a limit point.}  Because the sets are separated, $s \notin B$.  But $s \in [a,b] \subseteq E = A \cup B$, so $s \in A$.

In particular $s \neq b$, so $s < b$.  Now consider the points of $(s, b]$.  They all lie in $E$, and none of them is in $A$, since they are in $[a,b]$ and larger than the supremum of $A \cap [a,b]$.  So they are all in $B$.  But then $s$ is a limit point of $B$, and $s \in A \cap \overline{B}$, contradicting the assumption that $A$ and $B$ are separated.
\end{proof}

\begin{marginfigure}
\centering
\begin{tikzpicture}[x=1cm,y=1cm]
  \draw[numline] (0,0) -- (4.8,0);
  \draw[setline] (0.7,0) -- (2.6,0);
  \draw[nbhdline] (2.6,0) -- (4.1,0);
  \node[closedpt] at (0.7,0) {};
  \node[closedpt] at (2.6,0) {};
  \node[closedpt=calm] at (4.1,0) {};
  \node[lbl,below=3pt] at (0.7,0) {$a$};
  \node[lbl,below=3pt] at (2.6,0) {$s$};
  \node[lbl,below=3pt] at (4.1,0) {$b$};
  \node[lbl,above=3pt,calm] at (3.4,0) {all in $B$};
\end{tikzpicture}
\caption{The proof looks for the last point of $A$ before $b$.  That point $s$ belongs to $A$, yet the points of $B$ press right up against it.}
\end{marginfigure}

The sets with the property in Theorem~\ref{thm:connected-intervals} are exactly the intervals, if we allow the word to include single points, rays such as $[c, \infty)$, the whole line and the empty set.  So \emph{the connected subsets of $\mathbb{R}$ are precisely the intervals.}\marginnote{This theorem is the engine behind the Intermediate Value Theorem, which we will meet in Chapter Four.}

At the other extreme stands the Cantor set.  Between any two of its points there is a point that is not in it, since it contains no interval.  So not only is $C$ disconnected, but every subset of $C$ with more than one point is disconnected.  Such a set is called \textsc{totally disconnected}.  The Cantor set is as far from an interval as a set can be, and yet, like an interval, it is compact, perfect and uncountable.

\section{Beyond the Real Line: Metric Spaces}\label{sec:metric-spaces}

Look back over the definitions of open set, limit point, closed set and compact set.  Each was built from $\varepsilon$-neighborhoods, and an $\varepsilon$-neighborhood was built from a single ingredient, the distance $|x - y|$ between two real numbers.  We never added two elements of a set, and never asked which of two points was larger.  This suggests that the same definitions will make sense wherever there is a reasonable notion of distance.

\begin{definition}
A set $X$, whose elements we shall call points, is a \textsc{metric space} if with any two points $p$ and $q$ of $X$ there is associated a real number $d(p,q)$, called the \textsc{distance} from $p$ to $q$, such that
\begin{enumerate}[(i)]
	\item $d(p,q) > 0$ if $p \neq q$, and $d(p,p) = 0$;
	\item $d(p,q) = d(q,p)$;
	\item $d(p,q) \leq d(p,r) + d(r,q)$ for any $r \in X$.
\end{enumerate}
\end{definition}

Property (iii) is the \textsc{triangle inequality}.  It says that going from $p$ to $q$ by way of $r$ is never shorter than going directly.

\newthought{Examples of metric spaces.}
\begin{enumerate}
	\item The real line $\mathbb{R}$ with $d(p,q) = |p - q|$.  This is the metric space we have been studying all along.
	\item The rational numbers $\mathbb{Q}$ with the same distance.  More generally, any subset of a metric space is itself a metric space under the same distance function.
	\item The plane $\mathbb{R}^2 = \{(x_1, x_2) : x_1, x_2 \in \mathbb{R}\}$ with
	\[ d(\mathbf{x}, \mathbf{y}) = \sqrt{(x_1 - y_1)^2 + (x_2 - y_2)^2}, \]
	and in the same way the Euclidean space $\mathbb{R}^k$ for any $k$.  The properties of norms listed at the end of Chapter One show that the three conditions are satisfied.
	\item Any set $X$ at all, with the \textsc{discrete metric}
	\begin{displaymath}
	   d(p,q) = \left\{
	     \begin{array}{ll}
	       1 & \text{if } p \neq q, \\
	       0 & \text{if } p = q.
	     \end{array}
	   \right.
	\end{displaymath}
	In this space every point is at distance $1$ from every other.
\end{enumerate}

In a metric space the $\varepsilon$-neighborhood of a point $p$ is $V_{\varepsilon}(p) = \{q \in X : d(p,q) < \varepsilon\}$.  In the plane this is a disc without its rim, which is why neighborhoods are also called \emph{open balls}.  With this one change, Definitions~\ref{def:open}, \ref{def:limit-point}, \ref{def:closed} and \ref{def:compact} can be repeated word for word.  Do the theorems survive as well?  Many do, and with the same proofs.  Here is one whose proof shows the triangle inequality at work.

\begin{theorem}
In any metric space, every neighborhood $V_r(p)$ is an open set.
\end{theorem}

\begin{proof}
Let $q \in V_r(p)$, so that $d(p,q) < r$, and put $h = r - d(p,q) > 0$.  We claim that $V_h(q) \subseteq V_r(p)$.  If $s \in V_h(q)$, then by the triangle inequality
\[ d(p,s) \leq d(p,q) + d(q,s) < d(p,q) + h = r, \]
so $s \in V_r(p)$.  Thus every point $q$ of $V_r(p)$ has a neighborhood contained in $V_r(p)$.
\end{proof}

\begin{marginfigure}
\centering
\begin{tikzpicture}[x=1cm,y=1cm]
  \draw[calm, line width=0.8pt, dashed, fill=calm!8] (0,0) circle (2.1);
  \draw[accent, line width=0.8pt, dashed, fill=accent!14] (1.1,0.6) circle (0.84);
  \draw[black!60, line width=0.4pt] (0,0) -- (1.1,0.6) -- (1.844,1.005);
  \node[seqpt] at (0,0) {};
  \node[seqpt] at (1.1,0.6) {};
  \node[seqpt] at (1.45,0.2) {};
  \node[lbl,below left] at (0,0) {$p$};
  \node[lbl,above left] at (1.1,0.6) {$q$};
  \node[lbl,right=1pt] at (1.45,0.2) {$s$};
  \node[lbl,above=1pt] at (1.5,0.82) {$h$};
  \node[lbl,calm] at (-1.0,-1.1) {$V_r(p)$};
\end{tikzpicture}
\caption{An open ball is open.  The point $q$ is at distance $h$ from the rim, so the ball of radius $h$ about $q$ stays inside.}
\end{marginfigure}

\newthought{What survives and what does not.}  Theorems~\ref{thm:open-unions}, \ref{thm:limit-point-sequence}, \ref{thm:closed-sequences}, \ref{thm:closure} and \ref{thm:complements} hold in every metric space, and their proofs need no change beyond writing $d(x,y)$ in place of $|x - y|$.  It is also true in every metric space that a compact set is closed and bounded, and that a set is compact if and only if each of its open covers has a finite subcover, although the proof of the latter requires new ideas.

One thing does \emph{not} survive.

\begin{dangerousbend}
In a general metric space, a closed and bounded set need not be compact.
\begin{itemize}
	\item In the metric space $\mathbb{Q}$, the set $K = \{x \in \mathbb{Q} : 0 \leq x \leq 2\}$ is closed and bounded.  But the sequence $1,\ 1.4,\ 1.41,\ 1.414,\ \ldots$ of decimal approximations to $\sqrt2$ lies in $K$, and no subsequence of it converges to a point of $\mathbb{Q}$.
	\item Let $X$ be an infinite set with the discrete metric.  The whole space $X$ is closed and bounded, since no two points are more than $1$ apart.  But a sequence of distinct points of $X$ has no convergent subsequence, because its terms stay at distance $1$ from one another.
\end{itemize}
\end{dangerousbend}

So the implication ``closed and bounded implies compact'' in the Heine--Borel Theorem is really a theorem about $\mathbb{R}$, and about the Euclidean spaces $\mathbb{R}^k$, where it is also true.  Its proof rested on the Bolzano--Weierstrass Theorem, and so, in the end, on the Axiom of Completeness.

This brings us back to the promise we began with.  Of the five equivalent statements at the end of Chapter Two, three are phrased in terms of the order on $\mathbb{R}$.  They speak of suprema, of intervals, of monotone sequences.  The Bolzano--Weierstrass Theorem can be stated using distance alone, but the second example in the caution above shows that it asks for more than the absence of holes.  That leaves the Cauchy Criterion, which speaks of nothing but distance, and it is taken as the definition.  A metric space is \textsc{complete} if every Cauchy sequence in it converges to a point of the space.  The real line is complete, and $\mathbb{Q}$ is not.  That is the precise sense in which $\mathbb{R}$ has no holes, and the first example in the caution above shows what goes wrong when a space has them.

\section{Problems}

\noindent The problems in this section are original to this text.

\begin{problem}[Sorting sets]
\hint{Draw each set first.  For (f), think carefully about the point $0$.}
For each of the following sets, find its limit points and its isolated points, decide whether it is open, closed, both or neither, and write down its closure.
\begin{enumerate}[(a)]
	\item $\mathbb{N}$
	\item $\{x \in \mathbb{R} : x^2 < 2\}$
	\item $\{x \in \mathbb{Q} : x^2 < 2\}$
	\item $\left\{(-1)^n + \frac1n : n \in \mathbb{N}\right\}$
	\item $\displaystyle\bigcup_{n=1}^{\infty} \left[\frac{1}{2n}, \frac{1}{2n-1}\right]$
	\item $\left\{x \in \mathbb{R} : x \neq \frac1n \text{ for every } n \in \mathbb{N}\right\}$
\end{enumerate}
\end{problem}

\begin{problem}[Measuring the cushion]
\hint{For (b), if a ball of radius $\varepsilon$ around $a$ fits inside $O$, how large a ball around a nearby point $b$ is guaranteed to fit?}
Let $O$ be an open set other than $\mathbb{R}$ itself.  For $a \in O$ let $c(a)$ be the radius of the largest cushion around $a$:
\[ c(a) = \sup\{\varepsilon > 0 : V_{\varepsilon}(a) \subseteq O\}. \]
\begin{enumerate}[(a)]
	\item For $O = (0,1) \cup (1,3)$, find a formula for $c(a)$ and sketch its graph.
	\item For a general $O$, prove that $|c(a) - c(b)| \leq |a - b|$ for all $a, b \in O$.
	\item What does (b) say about the steepness of the graph you drew in (a)?
\end{enumerate}
\end{problem}

\begin{problem}[True or false]
\hint{The pair $A = (0,1)$, $B = (1,2)$ is a useful test case.}
Decide whether each statement is true for all sets $A, B \subseteq \mathbb{R}$.  Prove the true ones and give counterexamples to the others.
\begin{enumerate}[(a)]
	\item If $A$ is open and $B$ is closed, then $A \setminus B = \{x \in A : x \notin B\}$ is open.
	\item Every limit point of $A \cup B$ is a limit point of $A$ or a limit point of $B$.
	\item Every limit point of $A \cap B$ is a limit point of both $A$ and $B$.
	\item $\overline{A \cup B} = \overline{A} \cup \overline{B}$.
	\item $\overline{A \cap B} = \overline{A} \cap \overline{B}$.
	\item Every infinite set has a limit point.
	\item Every bounded infinite set has a limit point.
\end{enumerate}
\end{problem}

\begin{problem}[Lamps along the road]
\hint{In (a), three open intervals of length $\frac23$ have total length $2$.  Why is that not quite enough for a \emph{closed} interval of length $2$?}
For each set and each open cover, either find a finite subcover or prove that none exists.
\begin{enumerate}[(a)]
	\item The set $[0,2]$, covered by all the neighborhoods $V_{1/3}(x)$ with $x \in [0,2]$.  Find a finite subcover using as few sets as you can, and explain why no smaller number will do.
	\item The set $(0,1]$, covered by the intervals $\left(\frac{x}{2}, 2x\right)$ for $x \in (0,1]$.
	\item The set $[0,1]$, covered by the intervals in (b) together with one more open set $O_0$ that contains $0$.
	\item The set $\mathbb{Z}$, covered by the neighborhoods $V_{1/2}(n)$ for $n \in \mathbb{Z}$.
\end{enumerate}
\end{problem}

\begin{problem}[Arithmetic with compact sets]
\hint{For (c), take a sequence $x_n + y_n$ and choose a subsequence twice, once to make the $x$'s behave, and again to make the $y$'s behave.  For (d), try $A = \mathbb{N}$ and a set $B$ of negative numbers that are not quite integers.}
For sets $A, B \subseteq \mathbb{R}$ define $A + B = \{a + b : a \in A,\ b \in B\}$.  Let $K$ and $L$ be compact.
\begin{enumerate}[(a)]
	\item Prove that $K \cup L$ is compact.
	\item Prove that $K \cap L$ is compact.
	\item Prove that $K + L$ is compact.
	\item Show by example that if $A$ and $B$ are merely closed, then $A + B$ need not be closed.
\end{enumerate}
\end{problem}

\begin{problem}[Addresses in the Cantor set]
\hint{For (d), what do the two maps do to an address?}
Every point of the Cantor set has an address, an infinite sequence of L's and R's, as described in the text.
\begin{enumerate}[(a)]
	\item Find the address of $\frac34$.
	\item Which points have the addresses L, R, R, R, R, \ldots\ and R, L, L, L, L, \ldots?
	\item Show that a point of $C$ is an endpoint of one of the intervals making up some $C_n$ if and only if its address is eventually constant.  Deduce that the set of all such endpoints is countable, so that ``most'' points of $C$ are not endpoints.
	\item Show that if $x \in C$ then $\frac{x}{3} \in C$ and $\frac{x+2}{3} \in C$.  Conclude that the Cantor set is the union of two copies of itself, each one third the size.
\end{enumerate}
\end{problem}

\begin{problem}[A fat Cantor set]
\hint{For (b), how long can each of the $2^n$ intervals remaining after stage $n$ be?}
Repeat the construction of the Cantor set, but remove less.  Start with $[0,1]$.  At stage $n$, from the middle of each of the $2^{n-1}$ closed intervals that remain, remove an open interval of length $1/4^n$.  Let $F$ be the set of points that are never removed.
\begin{enumerate}[(a)]
	\item Show that the total length removed is $\frac12$.
	\item Show that $F$ nevertheless contains no interval.
	\item Is $F$ compact?  Is it perfect?
	\item The Cantor set contains no interval, and we said that it has ``length zero.''  What does $F$ suggest about the relationship between those two facts?
\end{enumerate}
\end{problem}

\begin{problem}[One piece or several?]
\hint{Theorem~\ref{thm:connected-intervals} turns every part of this problem into a question about intervals.}
\begin{enumerate}[(a)]
	\item Decide whether each of these three sets is connected.  They are the set of irrational numbers, the set $\{x \in \mathbb{R} : |x - 1| + |x + 1| \leq 3\}$, and the set $\{x \in \mathbb{R} : x^2 - 3x + 2 > 0\}$.
	\item Let $A$ and $B$ be connected subsets of $\mathbb{R}$ with $A \cap B \neq \emptyset$.  Prove that $A \cup B$ and $A \cap B$ are both connected.
	\item Prove that the closure of a connected set is connected.
	\item One of the two conclusions in (b) is special to the real line.  Draw two ``one-piece'' regions in the plane, overlapping, whose intersection is in two pieces.
\end{enumerate}
\end{problem}

\begin{problem}[From local to global]
\hint{This is the street-lamp argument, with one lamp for each point of $K$.}
Let $K$ be a compact set, and suppose that every point $x \in K$ has a neighborhood $V_{\varepsilon_x}(x)$ that contains only finitely many points of $K$.
\begin{enumerate}[(a)]
	\item Prove that $K$ is finite.
	\item Show by example that the conclusion can fail for a set that is closed but not bounded, and for a set that is bounded but not closed.
	\item Deduce that every infinite compact set has a limit point, and that this limit point belongs to the set.
\end{enumerate}
\end{problem}

\bookchapter{Four}{Functional Limits and Continuity}

In Chapter Two we made the phrase ``gets close to'' precise for sequences.  In this chapter we do the same for functions.  Loosely speaking, a function $f$ is continuous at a point $c$ if choosing $x$ close to $c$ forces $f(x)$ to be close to $f(c)$.  Our first job is to turn that sentence into a definition we can prove things with.

The reward for doing so is a trio of theorems about continuous functions.  A continuous function on a closed interval reaches a highest and a lowest value.  It cannot get from one value to another without passing through everything in between.  And its continuity can be measured by a single standard that works at every point at once.  Each of these facts turns out to be a statement about one of the kinds of set we met in Chapter Three, and the topology we built there will do most of the work.

\section{A Question of Tolerance}\label{sec:tolerance}

A tiler has been asked for a square tile with an area of exactly $9$ square centimeters.  She knows the side should be $3$ centimeters, but no saw cuts a length of exactly $3$.  Her client is reasonable about this and says that any area within $0.1$ of $9$ will do.  How accurately does she have to cut?

Suppose the side comes out as $x$ instead of $3$.  The error in the area is
\[ |x^2 - 9| = |x - 3|\,|x + 3|. \]
If she can at least promise to be within $1$ of the target, then $x$ is between $2$ and $4$, so $|x + 3|$ is less than $7$ and the error in the area is less than $7|x - 3|$.  To keep this below $0.1$ it is enough to cut the side to within $0.1/7$, which is about $0.014$ centimeters.

Now the client changes his mind and asks for an area within $0.001$ of $9$.  The tiler is not troubled.  The same calculation tells her to cut to within $0.001/7$.  Whatever tolerance the client names for the \emph{area}, she can answer with a tolerance for the \emph{side} that guarantees it.  That ability to answer every challenge is what we will mean by saying that the function $f(x) = x^2$ is continuous at $3$.

\newthought{Here is a function} for which the game cannot be won.  A post office charges $1.00$ to send a letter weighing up to $20$ grams and $1.50$ for anything heavier.  Suppose you are aiming for a weight of $20$ grams and want the price to be within $0.25$ of $1.00$.  No scale is accurate enough to promise this.  However small the error in the weight, a letter that comes out the slightest bit over $20$ grams costs $1.50$.  The challenge $0.25$ has no answer, and the price is not a continuous function of the weight at $20$.

\begin{figure}
\centering
\begin{tikzpicture}
  % ---- left panel: the tile
  \begin{scope}[x=1.8cm,y=0.3cm]
    \fill[epsband] (0,3.6) rectangle (1.95,6.4);
    \fill[delband] (0.79,0) rectangle (1.21,10.2);
    \draw[graphaxis] (0,0) -- (2.05,0) node[lbl,right] {side};
    \draw[graphaxis] (0,0) -- (0,10.8) node[lbl,above] {area};
    \draw[curve] plot[domain=2.2:3.72, samples=40] ({\x-2},{\x*\x-4});
    \draw[guide] (1,0) -- (1,5) -- (0,5);
    \node[seqpt] at (1,5) {};
    \node[lbl,below=2pt] at (1,0) {$3$};
    \node[lbl,left=2pt] at (0,5) {$9$};
    \node[lbl,left=2pt,calm] at (0,6.4) {$9+\varepsilon$};
    \node[lbl,left=2pt,calm] at (0,3.6) {$9-\varepsilon$};
  \end{scope}
  % ---- right panel: the postage
  \begin{scope}[shift={(5.5,0)},x=1.8cm,y=0.3cm]
    \fill[epsband] (0,1.6) rectangle (1.95,4.4);
    \fill[delband] (0.79,0) rectangle (1.21,10.2);
    \draw[graphaxis] (0,0) -- (2.05,0) node[lbl,right] {weight};
    \draw[graphaxis] (0,0) -- (0,10.8) node[lbl,above] {price};
    \draw[curve] (0.1,3) -- (1,3);
    \draw[curve] (1,7.5) -- (1.9,7.5);
    \node[closedpt] at (1,3) {};
    \node[openpt] at (1,7.5) {};
    \node[lbl,below=2pt] at (1,0) {$20$};
    \node[lbl,left=2pt] at (0,3) {$1.00$};
    \node[lbl,left=2pt] at (0,7.5) {$1.50$};
  \end{scope}
\end{tikzpicture}
\caption{Two tolerance games.  The blue band is the tolerance the client allows on the output and the red band is the tolerance we choose for the input.  On the left, the part of the graph above the red band stays inside the blue band, so we win.  On the right, no red band is narrow enough.}
\label{fig:tolerance}
\end{figure}

Notice the order of play.  The client speaks first and names a tolerance for the output.  We respond with a tolerance for the input.  This is the same challenge and response that we met in the definition of a convergent sequence, where the challenge was an $\varepsilon$ and the response was an $N$.  The challenge will again be called $\varepsilon$.  The response is traditionally called $\delta$.

\newthought{You may have been told} that a function is continuous when its graph can be drawn without lifting the pen from the paper.  That picture agrees with what we have just said about the tile and the postage.  But it is no help at all with a function whose graph cannot be drawn.  Consider the function that Dirichlet proposed in 1829,
\begin{displaymath}
   g(x) = \left\{
     \begin{array}{ll}
       1 & \text{if } x \in \mathbb{Q}, \\
       0 & \text{if } x \notin \mathbb{Q}.
     \end{array}
   \right.
\end{displaymath}
Between any two rational numbers there is an irrational one, and between any two irrational numbers a rational one, so the graph of $g$ is a pair of dotted lines with no two dots adjacent.  Is $g$ continuous anywhere?  What about the function $h(x) = x\,g(x)$?  The pen has no opinion.  The definition will.

\section{Functional Limits}

Before we define continuity we isolate the idea underneath it, which is the limit of a function at a point.  There is one subtlety to absorb first.  When we ask what $f(x)$ approaches as $x$ approaches $c$, we are asking about the values of $f$ \emph{near} $c$.  The value of $f$ \emph{at} $c$ is beside the point, and $f$ need not even be defined there.  What we do need is that there are points of the domain arbitrarily close to $c$ and different from $c$.  In the language of Chapter Three, $c$ must be a limit point of the domain (Definition~\ref{def:limit-point}).

\begin{definition}\label{def:functional-limit}
Let $f : A \to \mathbb{R}$ and let $c$ be a limit point of the domain $A$.  We say that $\lim_{x \to c} f(x) = L$ if, for every $\varepsilon > 0$, there exists a $\delta > 0$ such that whenever $0 < |x - c| < \delta$ and $x \in A$, it follows that $|f(x) - L| < \varepsilon$.
\end{definition}

\marginnote{The inequality $0 < |x - c|$ is only a compact way of saying $x \neq c$.  It is how the definition ignores what happens at $c$ itself.}

As with sequences, the definition can be recast in the language of neighborhoods, and the picture is easier to hold in the mind.

\begin{definition}[Topological Version]
Let $c$ be a limit point of the domain of $f : A \to \mathbb{R}$.  We say that $\lim_{x \to c} f(x) = L$ if, for every $\varepsilon$-neighborhood $V_{\varepsilon}(L)$ of $L$, there exists a $\delta$-neighborhood $V_{\delta}(c)$ of $c$ such that every $x \in V_{\delta}(c)$ other than $c$ that lies in $A$ satisfies $f(x) \in V_{\varepsilon}(L)$.
\end{definition}

\begin{marginfigure}
\centering
\begin{tikzpicture}[x=1.2cm,y=0.7cm]
  \fill[epsband] (0,1.6) rectangle (3.55,2.8);
  \fill[delband] (1.65,0) rectangle (2.35,4.9);
  \draw[graphaxis] (0,0) -- (3.75,0);
  \draw[graphaxis] (0,0) -- (0,5.2);
  \draw[curve] plot[domain=0.3:3.3, samples=40] (\x,{0.8+0.35*\x*\x});
  \node[openpt] at (2,2.2) {};
  \node[closedpt] at (2,3.9) {};
  \draw[guide] (2,0) -- (2,2.05);
  \node[lbl,below=2pt] at (2,0) {$c$};
  \node[lbl,left=2pt] at (0,2.2) {$L$};
  \node[lbl,calm] at (0.72,2.25) {$V_{\varepsilon}(L)$};
  \node[lbl,above=1pt,accent] at (2,4.9) {$V_{\delta}(c)$};
\end{tikzpicture}
\caption{The limit of $f$ at $c$ is $L$.  Over the red band, the graph stays inside the blue band.  The value of $f$ at $c$ itself, shown by the solid dot, plays no part.}
\end{marginfigure}

\begin{example}
We show that $\lim_{x \to 2} (3x + 1) = 7$.  Let $\varepsilon > 0$.  Since
\[ |(3x + 1) - 7| = 3\,|x - 2|, \]
the response $\delta = \varepsilon/3$ works.  If $0 < |x - 2| < \delta$ then $|(3x+1) - 7| = 3|x - 2| < 3\delta = \varepsilon$.
\end{example}

\begin{example}\label{ex:tile-limit}
We show that $\lim_{x \to 3} x^2 = 9$.  This is the tiler's calculation, written as a proof.  Let $\varepsilon > 0$ and choose
\[ \delta = \min\left\{1, \frac{\varepsilon}{7}\right\}. \]
Suppose $0 < |x - 3| < \delta$.  Because $\delta \leq 1$, we have $2 < x < 4$ and therefore $|x + 3| < 7$.  Because $\delta \leq \varepsilon/7$, we get
\[ |x^2 - 9| = |x - 3|\,|x + 3| < \frac{\varepsilon}{7} \cdot 7 = \varepsilon. \]
\end{example}

The device used here is worth remembering.  The factor $|x + 3|$ is not constant, but it cannot be large when $x$ is near $3$.  Agreeing in advance that $\delta$ will be at most $1$ gives us a fixed bound for it, and after that the problem is as easy as the previous one.

\begin{example}
Let $f(x) = \dfrac{x^2 - 9}{x - 3}$, a function whose domain is every real number except $3$.  The point $3$ is a limit point of the domain, so we may ask for the limit there.  For $x \neq 3$ we have $f(x) = x + 3$, so $|f(x) - 6| = |x - 3|$, and the response $\delta = \varepsilon$ shows that $\lim_{x \to 3} f(x) = 6$.  The function has a limit at a point where it has no value.
\end{example}

\newthought{Template for a proof that $\lim_{x \to c} f(x) = L$.}
\begin{enumerate}
	\item Let $\varepsilon > 0$ be arbitrary.
	\item Name a $\delta > 0$.  Finding it usually takes some scratch work, which is best done before the proof is written.
	\item Assume $0 < |x - c| < \delta$ and $x \in A$.
	\item Derive the inequality $|f(x) - L| < \varepsilon$.
\end{enumerate}

\subsection{Limits of Functions and Limits of Sequences}

We spent a whole chapter on the limits of sequences.  It would be a pity to prove everything a second time for functions, and we do not have to.  The next theorem is a bridge between the two kinds of limit.

\begin{theorem}[Sequential Criterion for Functional Limits]\label{thm:sequential-criterion}
Let $f : A \to \mathbb{R}$ and let $c$ be a limit point of $A$.  The following two statements are equivalent.
\begin{enumerate}[(i)]
	\item $\lim_{x \to c} f(x) = L$.
	\item For every sequence $(x_n)$ in $A$ with $x_n \neq c$ and $(x_n) \to c$, it follows that $f(x_n) \to L$.
\end{enumerate}
\end{theorem}

\begin{proof}
(i) $\Rightarrow$ (ii).  Assume $\lim_{x \to c} f(x) = L$, and let $(x_n)$ be a sequence in $A$ with $x_n \neq c$ and $(x_n) \to c$.  Let $\varepsilon > 0$.  The definition of the functional limit supplies a $\delta > 0$ such that $0 < |x - c| < \delta$ and $x \in A$ imply $|f(x) - L| < \varepsilon$.  Since $(x_n) \to c$, there is an $N$ such that $|x_n - c| < \delta$ for all $n \geq N$, and $x_n \neq c$ gives $0 < |x_n - c|$.  So $|f(x_n) - L| < \varepsilon$ for all $n \geq N$.

(ii) $\Rightarrow$ (i).  We prove the contrapositive.  Suppose that (i) is false.  Then there is some $\varepsilon_0 > 0$ for which no $\delta$ works.  In particular $\delta = \frac1n$ does not work, so for each $n \in \mathbb{N}$ there is a point $x_n \in A$ with
\[ 0 < |x_n - c| < \frac1n \qquad \text{and yet} \qquad |f(x_n) - L| \geq \varepsilon_0. \]
The sequence $(x_n)$ lies in $A$, avoids $c$, and converges to $c$, but $f(x_n)$ does not converge to $L$.  So (ii) is false as well.
\end{proof}

\marginnote{The second half of this proof is the same manoeuvre as in Theorem~\ref{thm:limit-point-sequence}.  When a statement about every $\delta$ fails, try $\delta = \frac1n$ and collect the offenders into a sequence.}

With this bridge in place, the Algebraic Limit Theorem crosses over from sequences to functions at no cost.

\begin{corollary}[Algebraic Limit Theorem for Functional Limits]\label{cor:alt-functions}
Let $f$ and $g$ be functions defined on $A$, and suppose $\lim_{x \to c} f(x) = L$ and $\lim_{x \to c} g(x) = M$ at a limit point $c$ of $A$.  Then
\begin{enumerate}[(i)]
	\item $\lim_{x \to c} k f(x) = kL$ for every $k \in \mathbb{R}$,
	\item $\lim_{x \to c} \big(f(x) + g(x)\big) = L + M$,
	\item $\lim_{x \to c} f(x) g(x) = LM$, and
	\item $\lim_{x \to c} f(x)/g(x) = L/M$, provided $M \neq 0$ and $g$ is never zero on $A$.
\end{enumerate}
\end{corollary}

\begin{proof}
Take any sequence $(x_n)$ in $A$ with $x_n \neq c$ and $(x_n) \to c$.  By Theorem~\ref{thm:sequential-criterion}, $f(x_n) \to L$ and $g(x_n) \to M$.  The Algebraic Limit Theorem for sequences then gives $f(x_n) + g(x_n) \to L + M$, and since the sequence was arbitrary, Theorem~\ref{thm:sequential-criterion} applied in the other direction gives (ii).  The other parts are proved in the same way.
\end{proof}

The bridge is just as useful for showing that a limit does \emph{not} exist.

\begin{corollary}[Divergence Criterion for Functional Limits]\label{cor:divergence-criterion}
Let $f : A \to \mathbb{R}$ and let $c$ be a limit point of $A$.  If there are two sequences $(x_n)$ and $(y_n)$ in $A$, with $x_n \neq c$ and $y_n \neq c$, such that
\[ \lim x_n = \lim y_n = c \qquad \text{but} \qquad \lim f(x_n) \neq \lim f(y_n), \]
then $\lim_{x \to c} f(x)$ does not exist.
\end{corollary}

\begin{example}
Dirichlet's function $g$ has no limit at any point.  Let $c \in \mathbb{R}$.  Because $\mathbb{Q}$ is dense in $\mathbb{R}$ there is a sequence of rational numbers $r_n \neq c$ converging to $c$, and because the irrationals are dense there is a sequence of irrational numbers $t_n \neq c$ converging to $c$.  Then $g(r_n) = 1$ for every $n$ and $g(t_n) = 0$ for every $n$.  The two sequences of values have different limits, so $\lim_{x \to c} g(x)$ does not exist.
\end{example}

\begin{example}
For an example you can draw, borrow the sine function from calculus\marginnote{We have not yet defined the sine function rigorously, so examples that use it are for illustration only.  Nothing we prove will depend on them.} and consider $f(x) = \sin(1/x)$ for $x \neq 0$.  As $x$ approaches $0$, the quantity $1/x$ runs off to infinity and the sine oscillates between $-1$ and $1$ faster and faster.  Taking $x_n = \frac{1}{2n\pi}$ and $y_n = \frac{1}{2n\pi + \pi/2}$ we have $f(x_n) = 0$ and $f(y_n) = 1$ for every $n$, so $\lim_{x \to 0} \sin(1/x)$ does not exist.
\end{example}

\begin{marginfigure}
\centering
\begin{tikzpicture}[x=4.0cm,y=1.25cm]
  \draw[graphaxis] (0,0) -- (1.08,0);
  \draw[graphaxis] (0,-1.25) -- (0,1.35);
  \draw[curve, line width=0.6pt] plot[domain=1:75, samples=1100, variable=\t] ({1/\t},{sin(deg(\t))});
  \draw[tickmark] (-0.02,1) -- (0.02,1);
  \draw[tickmark] (-0.02,-1) -- (0.02,-1);
  \node[lbl,left=2pt] at (0,1) {$1$};
  \node[lbl,left=2pt] at (0,-1) {$-1$};
  \node[lbl,below=2pt] at (1,0) {$1$};
  \draw[tickmark] (1,-0.04) -- (1,0.04);
\end{tikzpicture}
\caption{The graph of $\sin(1/x)$.  Near $0$ it visits every height between $-1$ and $1$ infinitely often and settles on none of them.}
\end{marginfigure}

\section{Continuous Functions}

We are ready for the main definition.

\begin{definition}\label{def:continuity}
A function $f : A \to \mathbb{R}$ is \textsc{continuous at a point} $c \in A$ if, for every $\varepsilon > 0$, there exists a $\delta > 0$ such that whenever $|x - c| < \delta$ and $x \in A$, it follows that $|f(x) - f(c)| < \varepsilon$.

If $f$ is continuous at every point of $A$, we say that $f$ is \textsc{continuous on} $A$.
\end{definition}

Set this beside the definition of a functional limit and two differences stand out.  The point $c$ must now belong to the domain, and the number $L$ has been replaced by $f(c)$.  The condition $0 < |x - c|$ has quietly disappeared, because $x = c$ now does no harm.  It only says that $|f(c) - f(c)| < \varepsilon$.  When $c$ is a limit point of $A$, the definition says precisely that
\[ \lim_{x \to c} f(x) = f(c). \]
A function is continuous at $c$ when its limit at $c$ exists and is the value it actually takes there.\marginnote{If $c$ is an isolated point of $A$, then $f$ is automatically continuous at $c$, since a small enough $\delta$ leaves no $x$ to test except $c$ itself.  This is harmless, if a little odd.}

The next theorem gathers the different ways of saying this.

\begin{theorem}[Characterizations of Continuity]\label{thm:continuity-characterizations}
Let $f : A \to \mathbb{R}$ and let $c \in A$.  The following statements are equivalent.
\begin{enumerate}[(i)]
	\item For every $\varepsilon > 0$ there exists a $\delta > 0$ such that $|x - c| < \delta$ and $x \in A$ imply $|f(x) - f(c)| < \varepsilon$.
	\item For every neighborhood $V_{\varepsilon}(f(c))$ there exists a neighborhood $V_{\delta}(c)$ such that $x \in V_{\delta}(c) \cap A$ implies $f(x) \in V_{\varepsilon}(f(c))$.
	\item For every sequence $(x_n)$ in $A$ with $(x_n) \to c$, it follows that $f(x_n) \to f(c)$.
\end{enumerate}
\end{theorem}

\begin{proof}
Statement (ii) is statement (i) written with neighborhoods, so these two are equivalent.  The equivalence of (i) and (iii) is proved exactly as in Theorem~\ref{thm:sequential-criterion}, with $f(c)$ in place of $L$.  The only change is that the terms $x_n$ are now allowed to equal $c$, and such terms cause no difficulty because $|f(c) - f(c)| = 0$.
\end{proof}

Statement (iii) is the one to carry around.  It says that continuous functions are the functions that respect limits.  If $f$ is continuous and $(x_n)$ converges, then
\[ \lim f(x_n) = f\big(\lim x_n\big), \]
so a continuous function may be moved through a limit sign.  We have in fact done this already.  In Exercise 2.4.4 we took a limit inside a square root, and the justification we gave there amounts to the continuity of the square root function.

\begin{corollary}[Criterion for Discontinuity]\label{cor:discontinuity-criterion}
Let $f : A \to \mathbb{R}$ and $c \in A$.  If there is a sequence $(x_n)$ in $A$ with $(x_n) \to c$ such that $f(x_n)$ does not converge to $f(c)$, then $f$ is not continuous at $c$.
\end{corollary}

\begin{theorem}[Algebraic Continuity Theorem]\label{thm:algebraic-continuity}
Suppose $f : A \to \mathbb{R}$ and $g : A \to \mathbb{R}$ are continuous at a point $c \in A$.  Then $kf$ is continuous at $c$ for every $k \in \mathbb{R}$, and so are $f + g$ and $fg$.  The quotient $f/g$ is continuous at $c$ as well, provided $g$ is never zero on $A$.
\end{theorem}

\begin{proof}
Let $(x_n)$ be a sequence in $A$ converging to $c$.  Then $f(x_n) \to f(c)$ and $g(x_n) \to g(c)$, and the Algebraic Limit Theorem for sequences gives $f(x_n) + g(x_n) \to f(c) + g(c)$, and similarly for the other combinations.  Now apply part (iii) of Theorem~\ref{thm:continuity-characterizations}.
\end{proof}

\begin{example}
A constant function is continuous everywhere, since any $\delta$ works.  The function $f(x) = x$ is continuous everywhere, since $\delta = \varepsilon$ works.  By the Algebraic Continuity Theorem, every polynomial
\[ p(x) = a_0 + a_1 x + a_2 x^2 + \cdots + a_n x^n \]
is continuous on $\mathbb{R}$, being built from these two by finitely many products and sums.  A quotient of two polynomials is continuous wherever its denominator is not zero.
\end{example}

\begin{example}\label{ex:sqrt-continuous}
The function $f(x) = \sqrt{x}$ is continuous on $[0, \infty)$.  At $c = 0$, given $\varepsilon > 0$ take $\delta = \varepsilon^2$.  Then $0 \leq x < \delta$ gives $\sqrt{x} < \varepsilon$.  At a point $c > 0$ we use the identity
\[ |\sqrt{x} - \sqrt{c}| = \frac{|x - c|}{\sqrt{x} + \sqrt{c}} \leq \frac{|x - c|}{\sqrt{c}}, \]
which shows that $\delta = \varepsilon\sqrt{c}$ works.
\end{example}

\begin{example}
Dirichlet's function $g$ is continuous nowhere.  We saw that $g$ has no limit at any point $c$, so its limit at $c$ certainly cannot equal $g(c)$.
\end{example}

\begin{example}\label{ex:one-point}
Now consider
\begin{displaymath}
   h(x) = x\,g(x) = \left\{
     \begin{array}{ll}
       x & \text{if } x \in \mathbb{Q}, \\
       0 & \text{if } x \notin \mathbb{Q}.
     \end{array}
   \right.
\end{displaymath}
At $c = 0$ this function is continuous.  Whether $x$ is rational or not, $|h(x) - h(0)| = |h(x)| \leq |x|$, so $\delta = \varepsilon$ works.  At any other point $c \neq 0$ it is not.  A sequence of rationals $r_n \to c$ has $h(r_n) = r_n \to c$, while a sequence of irrationals $t_n \to c$ has $h(t_n) = 0 \to 0$, and $c \neq 0$.  One of the two limits must differ from $h(c)$, and the Criterion for Discontinuity applies.

So $h$ is continuous at exactly one point.  The idea of drawing a graph without lifting the pen could never have told us that.
\end{example}

\begin{marginfigure}
\centering
\begin{tikzpicture}[x=1.15cm,y=1.15cm]
  \fill[epsband] (-1.9,-0.5) rectangle (1.9,0.5);
  \fill[delband] (-0.5,-1.7) rectangle (0.5,1.7);
  \draw[graphaxis] (-2,0) -- (2.05,0);
  \draw[graphaxis] (0,-1.8) -- (0,1.9);
  \draw[accent, line width=1.1pt, dash pattern=on 1pt off 2.2pt] (-1.7,-1.7) -- (1.7,1.7);
  \draw[accent, line width=1.1pt, dash pattern=on 1pt off 2.2pt] (-1.9,0) -- (1.9,0);
  \node[seqpt] at (0,0) {};
  \node[lbl,anchor=east] at (1.42,1.62) {$x \in \mathbb{Q}$};
  \node[lbl,below=2pt] at (1.5,0) {$x \notin \mathbb{Q}$};
  \node[lbl,left=2pt,calm] at (0,0.5) {$\varepsilon$};
  \node[lbl,below=1pt,accent] at (0.5,-1.7) {$\delta$};
\end{tikzpicture}
\caption{The function $h$ lives on two dotted lines.  Near the origin both lines are inside the blue band, and nowhere else do they come together.}
\end{marginfigure}

Functions are combined by composition as well as by arithmetic, and continuity survives that too.

\begin{theorem}[Composition of Continuous Functions]\label{thm:composition}
Let $f : A \to \mathbb{R}$ and $g : B \to \mathbb{R}$, where the range of $f$ is contained in $B$, so that the composition $g \circ f(x) = g(f(x))$ is defined on $A$.  If $f$ is continuous at $c \in A$ and $g$ is continuous at $f(c)$, then $g \circ f$ is continuous at $c$.
\end{theorem}

\begin{proof}
Let $(x_n)$ be a sequence in $A$ with $(x_n) \to c$.  Since $f$ is continuous at $c$, we have $f(x_n) \to f(c)$.  This is a sequence in $B$ converging to $f(c)$, and $g$ is continuous at $f(c)$, so $g(f(x_n)) \to g(f(c))$.
\end{proof}

For instance, $\sqrt{x^2 + 1}$ is continuous on $\mathbb{R}$, being the composition of a polynomial with the square root.

\section{Continuous Functions on Compact Sets}

A function $f : A \to \mathbb{R}$ carries each subset $B$ of its domain to a set of values, $f(B) = \{f(x) : x \in B\}$.  If $f$ is continuous, which properties of $B$ are passed on to $f(B)$?  The answer is fewer than one might hope.
\begin{itemize}
	\item \emph{Open sets need not go to open sets.}  The function $f(x) = x^2$ takes the open interval $(-1,1)$ to $[0,1)$.
	\item \emph{Closed sets need not go to closed sets.}  The function $f(x) = \frac{1}{1+x^2}$ takes the closed set $[0, \infty)$ to $(0,1]$.
	\item \emph{Bounded sets need not go to bounded sets.}  The function $f(x) = \frac1x$ takes the bounded set $(0,1)$ to $(1, \infty)$.
\end{itemize}
In the second example the set was closed but not bounded, and in the third it was bounded but not closed.  When a set is both, the story changes completely.

\begin{theorem}[Preservation of Compact Sets]\label{thm:preserve-compact}
Let $f : A \to \mathbb{R}$ be continuous on $A$.  If $K \subseteq A$ is compact, then $f(K)$ is compact.
\end{theorem}

\begin{proof}
Let $(y_n)$ be any sequence in $f(K)$.  We must find a subsequence converging to a point of $f(K)$.  Each $y_n$ is a value of $f$ on $K$, so we can choose $x_n \in K$ with $f(x_n) = y_n$.  Since $K$ is compact, $(x_n)$ has a subsequence $(x_{n_k})$ converging to a point $x \in K$.  By continuity,
\[ y_{n_k} = f(x_{n_k}) \to f(x), \]
and $f(x) \in f(K)$.
\end{proof}

The proof is short because the definitions were made for each other.  Compactness is about sequences having convergent subsequences, and continuity is about functions preserving convergence.  An immediate consequence is one of the most useful theorems in analysis.

\begin{theorem}[Extreme Value Theorem]\label{thm:evt}
If $f : K \to \mathbb{R}$ is continuous on a nonempty compact set $K$, then $f$ attains a maximum and a minimum value.  That is, there are points $x_0, x_1 \in K$ such that
\[ f(x_0) \leq f(x) \leq f(x_1) \qquad \text{for all } x \in K. \]
\end{theorem}

\begin{proof}
The set $f(K)$ is compact, so it is bounded and closed.  Because it is bounded and nonempty, $\alpha = \sup f(K)$ exists.  A supremum is either a member of its set or a limit point of it, as we observed in the proof of Theorem~\ref{thm:connected-intervals}, and $f(K)$ is closed, so in either case $\alpha \in f(K)$.  This means $\alpha = f(x_1)$ for some $x_1 \in K$, and $f(x) \leq \alpha = f(x_1)$ for all $x \in K$.  The minimum is found in the same way from $\inf f(K)$.
\end{proof}

Every hypothesis of the Extreme Value Theorem is needed, as Figure~\ref{fig:evt-failures} shows.

\begin{figure}
\centering
\begin{tikzpicture}
  % (a) domain not closed
  \begin{scope}[x=2.3cm,y=2.3cm]
    \draw[graphaxis] (0,0) -- (1.2,0);
    \draw[graphaxis] (0,0) -- (0,1.2);
    \draw[guide] (0,1) -- (1,1) -- (1,0);
    \draw[curve] (0,0) -- (1,1);
    \node[openpt] at (0,0) {};
    \node[openpt] at (1,1) {};
    \node[lbl,below=2pt] at (1,0) {$1$};
    \node[lbl,left=2pt] at (0,1) {$1$};
    \node[lbl] at (0.6,-0.42) {$f(x) = x$ on $(0,1)$};
    \node[lbl] at (0.6,-0.66) {domain not closed};
  \end{scope}
  % (b) domain not bounded
  \begin{scope}[shift={(3.75,0)},x=0.42cm,y=2.3cm]
    \draw[graphaxis] (0,0) -- (6.6,0);
    \draw[graphaxis] (0,0) -- (0,1.2);
    \draw[guide] (0,1) -- (6.4,1);
    \draw[curve] plot[domain=0:6.3, samples=50] (\x,{\x/(1+\x)});
    \node[closedpt] at (0,0) {};
    \node[lbl,left=2pt] at (0,1) {$1$};
    \node[lbl] at (3.3,-0.42) {$f(x) = \frac{x}{1+x}$ on $[0,\infty)$};
    \node[lbl] at (3.3,-0.66) {domain not bounded};
  \end{scope}
  % (c) not continuous
  \begin{scope}[shift={(7.5,0)},x=2.3cm,y=2.3cm]
    \draw[graphaxis] (0,0) -- (1.2,0);
    \draw[graphaxis] (0,0) -- (0,1.2);
    \draw[guide] (0,1) -- (1,1) -- (1,0);
    \draw[curve] (0,0) -- (1,1);
    \node[closedpt] at (0,0) {};
    \node[openpt] at (1,1) {};
    \node[closedpt] at (1,0) {};
    \node[lbl,below=3pt] at (1,0) {$1$};
    \node[lbl,left=2pt] at (0,1) {$1$};
    \node[lbl] at (0.6,-0.42) {$f(1)$ moved to $0$};
    \node[lbl] at (0.6,-0.66) {not continuous};
  \end{scope}
\end{tikzpicture}
\caption{Three functions that never reach the value $1$, although they come arbitrarily close to it.  Each one fails exactly one hypothesis of the Extreme Value Theorem.}
\label{fig:evt-failures}
\end{figure}

\subsection{Uniform Continuity}

Let us go back to the tiler.  For a tile of side $3$ she found that an area tolerance of $\varepsilon$ could be met by cutting the side to within $\varepsilon/7$.  Suppose she is now asked for a much larger tile, of side $30$.  The same calculation gives
\[ |x^2 - 900| = |x - 30|\,|x + 30| < 61\,|x - 30| \]
when $x$ is within $1$ of $30$, so now she must cut to within $\varepsilon/61$.  The function is the same and the tolerance on the area is the same, but the answer has become nearly nine times more demanding.  Where the graph of $x^2$ is steep, a small error in the side makes a large error in the area.

This is a feature of the definition of continuity that is easy to overlook.  When we say $f$ is continuous on $A$, the response $\delta$ is allowed to depend on the point $c$ as well as on the challenge $\varepsilon$.  For some functions there is a single $\delta$ that answers a given $\varepsilon$ at every point of the domain simultaneously.

\begin{definition}\label{def:uniform-continuity}
A function $f : A \to \mathbb{R}$ is \textsc{uniformly continuous} on $A$ if, for every $\varepsilon > 0$, there exists a $\delta > 0$ such that for all $x, y \in A$, $|x - y| < \delta$ implies $|f(x) - f(y)| < \varepsilon$.
\end{definition}

The difference from ordinary continuity is entirely in the order of the quantifiers.  For continuity on $A$, we are told the point $c$ and the tolerance $\varepsilon$ and then choose $\delta$.  For uniform continuity, we are told only $\varepsilon$, and our $\delta$ must be chosen before we know where it will be used.

\begin{figure}
\centering
\begin{tikzpicture}[x=3.3cm,y=0.78cm]
  \fill[epsband] (0.273,-0.14) rectangle (0.927,0.86);
  \fill[epsband] (1.773,3.11) rectangle (2.027,4.11);
  \draw[graphaxis] (0,0) -- (2.42,0);
  \draw[graphaxis] (0,-0.2) -- (0,5.6);
  \draw[curve] plot[domain=0:2.3, samples=50] (\x,{\x*\x});
  \draw[guide] (0.273,-0.14) -- (0.273,-0.55);
  \draw[guide] (0.927,-0.14) -- (0.927,-0.55);
  \draw[guide] (1.773,3.11) -- (1.773,-0.55);
  \draw[guide] (2.027,3.11) -- (2.027,-0.55);
  \draw[bracedeco] (0.927,-0.62) -- (0.273,-0.62) node[midway,below=4pt,lbl] {wide};
  \draw[bracedeco] (2.027,-0.62) -- (1.773,-0.62) node[midway,below=4pt,lbl] {narrow};
  \node[seqpt] at (0.6,0.36) {};
  \node[seqpt] at (1.9,3.61) {};
  \draw[bracedeco] (0.927,0.86) -- (0.927,-0.14) node[midway,right=4pt,lbl,calm] {$2\varepsilon$};
  \draw[bracedeco] (2.027,4.11) -- (2.027,3.11) node[midway,right=4pt,lbl,calm] {$2\varepsilon$};
\end{tikzpicture}
\caption{The function $x^2$ at two different points.  The two boxes have the same height, which is the tolerance on the output.  Where the graph is steeper, the box that keeps the graph inside must be narrower.  On the whole real line no single width will do.}
\label{fig:uniform}
\end{figure}

\begin{example}
The function $f(x) = 3x + 1$ is uniformly continuous on $\mathbb{R}$.  Since $|f(x) - f(y)| = 3|x - y|$, the response $\delta = \varepsilon/3$ works for every pair of points.  A straight line is equally steep everywhere.
\end{example}

To show that a function is \emph{not} uniformly continuous we need to negate the definition, and it is convenient to have the negation in the form of sequences.

\begin{theorem}[Sequential Criterion for Absence of Uniform Continuity]\label{thm:not-uniform}
A function $f : A \to \mathbb{R}$ fails to be uniformly continuous on $A$ if and only if there exist an $\varepsilon_0 > 0$ and two sequences $(x_n)$ and $(y_n)$ in $A$ such that
\[ |x_n - y_n| \to 0 \qquad \text{but} \qquad |f(x_n) - f(y_n)| \geq \varepsilon_0 \quad \text{for all } n. \]
\end{theorem}

\begin{proof}
To say that $f$ is not uniformly continuous is to say that there is an $\varepsilon_0 > 0$ for which no $\delta$ works.  If so, then for each $n$ the choice $\delta = \frac1n$ fails, which means there are points $x_n, y_n \in A$ with $|x_n - y_n| < \frac1n$ and $|f(x_n) - f(y_n)| \geq \varepsilon_0$.  These are the required sequences.

Conversely, if such sequences exist, then no $\delta$ can answer the challenge $\varepsilon_0$.  Given any $\delta > 0$ there is an $n$ with $|x_n - y_n| < \delta$, and yet $|f(x_n) - f(y_n)| \geq \varepsilon_0$.
\end{proof}

\begin{example}
The function $f(x) = x^2$ is not uniformly continuous on $\mathbb{R}$.  Take $x_n = n$ and $y_n = n + \frac1n$.  Then $|x_n - y_n| = \frac1n \to 0$, but
\[ |f(x_n) - f(y_n)| = \left(n + \frac1n\right)^2 - n^2 = 2 + \frac{1}{n^2} \geq 2. \]
On the other hand, the same function \emph{is} uniformly continuous on the interval $[0,10]$.  For $x, y$ in that interval, $|x^2 - y^2| = |x - y|\,|x + y| \leq 20\,|x - y|$, so $\delta = \varepsilon/20$ works.  Uniform continuity is a property of a function together with its domain.
\end{example}

\begin{example}
Unboundedness of the domain is not the only obstacle.  The function $f(x) = \frac1x$ is continuous on the bounded interval $(0,1)$ but not uniformly continuous there.  Take $x_n = \frac1n$ and $y_n = \frac{1}{n+1}$.  Then $|x_n - y_n| = \frac{1}{n(n+1)} \to 0$, but $|f(x_n) - f(y_n)| = 1$.
\end{example}

In the first example the trouble came from a domain that was not bounded, and in the second from a domain that was not closed.  By now this pattern should be familiar.

\begin{theorem}[Uniform Continuity on Compact Sets]\label{thm:uniform-compact}
A function that is continuous on a compact set $K$ is uniformly continuous on $K$.
\end{theorem}

\begin{proof}
This is the local-to-global principle of Chapter Three in action.  Continuity gives us a $\delta$ at each point, and compactness will let us get by with finitely many of them.

Let $\varepsilon > 0$.  For each point $c \in K$, continuity at $c$ gives a $\delta_c > 0$ such that
\[ |x - c| < \delta_c \ \text{ and } \ x \in K \quad \text{imply} \quad |f(x) - f(c)| < \frac{\varepsilon}{2}. \]
The neighborhoods $V_{\delta_c/2}(c)$, one for each $c \in K$, form an open cover of $K$.  Note that we have used only half of each radius.  By the Heine--Borel Theorem (Theorem~\ref{thm:heine-borel}), finitely many of these neighborhoods cover $K$.  Call their centers $c_1, \ldots, c_m$, and let
\[ \delta = \min\left\{\frac{\delta_{c_1}}{2}, \ldots, \frac{\delta_{c_m}}{2}\right\}. \]
Now suppose $x, y \in K$ with $|x - y| < \delta$.  The point $x$ lies in one of the finitely many neighborhoods, say $|x - c_i| < \delta_{c_i}/2$.  Then $y$ is close to the same center, because
\[ |y - c_i| \leq |y - x| + |x - c_i| < \delta + \frac{\delta_{c_i}}{2} \leq \delta_{c_i}. \]
So both $x$ and $y$ are within $\delta_{c_i}$ of $c_i$, and therefore
\[ |f(x) - f(y)| \leq |f(x) - f(c_i)| + |f(c_i) - f(y)| < \frac{\varepsilon}{2} + \frac{\varepsilon}{2} = \varepsilon. \qedhere \]
\end{proof}

\marginnote{Why cover $K$ with neighborhoods of \emph{half} the radius?  To leave room.  Knowing that $x$ is within $\delta_{c_i}/2$ of the center, we can still move a little way to $y$ without leaving the full neighborhood.}

\section{The Intermediate Value Theorem}

\begin{theorem}[Intermediate Value Theorem]\label{thm:ivt}
Let $f : [a,b] \to \mathbb{R}$ be continuous.  If $L$ is a real number strictly between $f(a)$ and $f(b)$, then there is a point $c \in (a,b)$ with $f(c) = L$.
\end{theorem}

\begin{marginfigure}
\centering
\begin{tikzpicture}[x=0.88cm,y=1cm]
  \draw[graphaxis] (0,0) -- (4.75,0);
  \draw[graphaxis] (0,0) -- (0,4);
  \draw[guide] (0,1.9) -- (4.5,1.9);
  \draw[guide] (0.5,0) -- (0.5,0.127) -- (0,0.127);
  \draw[guide] (4.2,0) -- (4.2,3.5) -- (0,3.5);
  \draw[guide] (2.4,0) -- (2.4,1.9);
  \draw[curve] plot[domain=0.5:4.2, samples=50] (\x,{1.9+0.5*(\x-2.4)+0.12*(\x-2.4)^3});
  \node[closedpt] at (0.5,0.127) {};
  \node[closedpt] at (4.2,3.5) {};
  \node[seqpt] at (2.4,1.9) {};
  \node[lbl,below=2pt] at (0.5,0) {$a$};
  \node[lbl,below=2pt] at (2.4,0) {$c$};
  \node[lbl,below=2pt] at (4.2,0) {$b$};
  \node[lbl,left=2pt] at (0,0.3) {$f(a)$};
  \node[lbl,left=2pt] at (0,1.9) {$L$};
  \node[lbl,left=2pt] at (0,3.5) {$f(b)$};
\end{tikzpicture}
\caption{To get from the height $f(a)$ to the height $f(b)$, a continuous graph must cross every level $L$ in between.}
\end{marginfigure}

Of all the theorems in this chapter, this one looks the most obvious.  A pen that starts below a line and ends above it must cross the line.  But the statement is not as innocent as it looks.  Consider the function $f(x) = x^2 - 2$, with only rational numbers allowed as inputs.  Then $f(1) = -1$ and $f(2) = 2$, and yet there is no rational $c$ with $f(c) = 0$, as we proved at the start of Chapter One.  Among the rationals the theorem is false.  Any proof of it must therefore use the completeness of $\mathbb{R}$ in an essential way.

Our proof will use completeness in the form of Theorem~\ref{thm:connected-intervals}, which says that the connected subsets of $\mathbb{R}$ are the intervals.  What we need in addition is that continuous functions do not tear connected sets apart.

\begin{theorem}[Preservation of Connected Sets]\label{thm:preserve-connected}
Let $f : A \to \mathbb{R}$ be continuous.  If $E \subseteq A$ is connected, then $f(E)$ is connected.
\end{theorem}

\begin{proof}
We prove the contrapositive.  Suppose $f(E)$ is disconnected, so that $f(E) = P \cup Q$ where $P$ and $Q$ are nonempty separated sets.  Pull these two pieces back to the domain by setting
\[ S = \{x \in E : f(x) \in P\}, \qquad T = \{x \in E : f(x) \in Q\}. \]
Both are nonempty, and $E = S \cup T$.  We claim that $S$ and $T$ are separated, which will show that $E$ is disconnected.

Suppose, for contradiction, that some point $x \in T$ belongs to $\overline{S}$.  The sets $S$ and $T$ have no point in common, because $P$ and $Q$ have none.  So $x$ must be a limit point of $S$, and there is a sequence $(s_n)$ in $S$ with $(s_n) \to x$.  Since $f$ is continuous at $x$, we have $f(s_n) \to f(x)$.  Each $f(s_n)$ lies in $P$, so the limit $f(x)$ lies in the closed set $\overline{P}$.  But $x \in T$ means $f(x) \in Q$.  This puts $f(x)$ in $\overline{P} \cap Q$, which is empty.  The contradiction shows that $\overline{S} \cap T = \emptyset$, and the same argument with the roles exchanged gives $S \cap \overline{T} = \emptyset$.
\end{proof}

\begin{proof}[Proof of the Intermediate Value Theorem]
The interval $[a,b]$ is connected, so $f([a,b])$ is connected by Theorem~\ref{thm:preserve-connected}.  It contains the numbers $f(a)$ and $f(b)$, and $L$ lies strictly between them.  By Theorem~\ref{thm:connected-intervals}, $L$ belongs to $f([a,b])$ too.  So $L = f(c)$ for some $c \in [a,b]$, and $c$ is neither $a$ nor $b$ because $L$ is neither $f(a)$ nor $f(b)$.
\end{proof}

Here are two things the theorem can do.

\begin{example}[Roots exist]\label{ex:roots}
Let $y > 0$ and $n \in \mathbb{N}$.  We claim there is a positive real number $x$ with $x^n = y$.  The function $f(x) = x^n$ is a polynomial, so it is continuous on the interval $[0, 1 + y]$.  At the left endpoint, $f(0) = 0 < y$.  At the right endpoint, $f(1 + y) = (1+y)^n \geq 1 + y > y$.  By the Intermediate Value Theorem there is a point $x$ strictly between $0$ and $1 + y$ where $f(x) = y$.

In particular there is a real number whose square is $2$.  We have been writing $\sqrt2$ since Chapter One on trust, and that debt is now paid.
\end{example}

\begin{example}[Fixed points]
Let $f : [0,1] \to [0,1]$ be continuous.  Then $f$ has a \emph{fixed point}, a number $x$ with $f(x) = x$.  To see this, let $g(x) = f(x) - x$, which is continuous on $[0,1]$.  Since the values of $f$ lie in $[0,1]$, we have $g(0) = f(0) \geq 0$ and $g(1) = f(1) - 1 \leq 0$.  If either of these is zero we have found a fixed point.  Otherwise $g(0) > 0 > g(1)$, and the Intermediate Value Theorem gives a point $x$ with $g(x) = 0$.
\end{example}

\begin{marginfigure}
\centering
\begin{tikzpicture}[x=3.6cm,y=3.6cm]
  \draw[tickmark] (0,0) rectangle (1,1);
  \draw[guide] (0,0) -- (1,1);
  \draw[curve] plot[domain=0:1, samples=40] (\x,{0.75-0.55*\x*\x});
  \node[closedpt] at (0,0.75) {};
  \node[closedpt] at (1,0.2) {};
  \node[seqpt] at (0.5708,0.5708) {};
  \node[lbl,below=2pt] at (0,0) {$0$};
  \node[lbl,below=2pt] at (1,0) {$1$};
  \node[lbl,left=2pt] at (0,1) {$1$};
  \node[lbl,rotate=45,above=1pt] at (0.25,0.25) {$y = x$};
\end{tikzpicture}
\caption{A continuous graph that starts on the left side of the square and ends on the right side must cross the diagonal.  The crossing is a fixed point.}
\end{marginfigure}

\begin{dangerousbend}
The converse of the Intermediate Value Theorem is false.  A function can take every intermediate value on every interval without being continuous.  The function equal to $\sin(1/x)$ for $x \neq 0$ and to $0$ at $x = 0$ is an example.  Taking all the values in between is a consequence of continuity, not a definition of it.
\end{dangerousbend}

\section{How Discontinuous Can a Function Be?}

We have seen a function that is discontinuous at every point, and another that is continuous at just one.  It is natural to ask which sets can arise as the set of points where a function fails to be continuous.  Could a function be continuous at every irrational number and discontinuous at every rational one?  It sounds impossible, since the two kinds of number are so thoroughly mixed together.  Here is such a function, discovered by Carl Thomae in 1875.
\begin{displaymath}
   t(x) = \left\{
     \begin{array}{ll}
       \dfrac1n & \text{if } x = \dfrac{m}{n} \text{ is rational, in lowest terms with } n > 0, \\[2ex]
       0 & \text{if } x \text{ is irrational.}
     \end{array}
   \right.
\end{displaymath}
So $t\left(\frac12\right) = \frac12$, $t\left(\frac34\right) = \frac14$ and $t(\sqrt2) = 0$, while $t(0) = t(1) = 1$ because the integers have denominator $1$.

\begin{figure}
\centering
\begin{tikzpicture}[x=9cm,y=6.2cm]
  \draw[graphaxis] (0,0) -- (1.06,0);
  \draw[graphaxis] (0,0) -- (0,1.12);
  \foreach \n in {2,...,26} {
    \foreach \m in {1,...,\n} {
      \pgfmathtruncatemacro{\g}{gcd(\m,\n)}
      \ifnum\g=1
        \fill[accent] ({\m/\n},{1/\n}) circle (0.75pt);
      \fi
    }
  }
  \fill[accent] (0,1) circle (1.1pt);
  \fill[accent] (1,1) circle (1.1pt);
  \foreach \y/\l in {1/1, 0.5/\frac12, 0.3333/\frac13, 0.25/\frac14} {
    \draw[tickmark] (-0.008,\y) -- (0.008,\y);
    \node[lbl,left=2pt] at (0,\y) {$\l$};
  }
  \foreach \x/\l in {0.5/\frac12, 1/1} {
    \draw[tickmark] (\x,-0.02) -- (\x,0.02);
    \node[lbl,below=3pt] at (\x,0) {$\l$};
  }
  \node[lbl,below=3pt] at (0,0) {$0$};
\end{tikzpicture}
\caption{Thomae's function on $[0,1]$, showing the rationals with denominator up to $26$.  Its value at every irrational is $0$.  Only a few points are high up, and near any irrational number the points crowd down toward the axis.}
\label{fig:thomae}
\end{figure}

\newthought{$t$ is discontinuous at every rational number.}  Let $c$ be rational, so that $t(c) > 0$.  Choose irrational numbers $t_n \to c$.  Then $t(t_n) = 0$ for all $n$, which does not converge to $t(c)$.

\newthought{$t$ is continuous at every irrational number.}  Let $c$ be irrational, so that $t(c) = 0$, and let $\varepsilon > 0$.  Where is $t$ large?  The value $t(x)$ can be $\varepsilon$ or more only if $x$ is a rational number whose denominator is at most $1/\varepsilon$.  In the interval $(c - 1, c + 1)$ there are only finitely many such rationals, since there are finitely many denominators to consider and each one has finitely many numerators that land in the interval.  None of these finitely many points is $c$, because $c$ is irrational.  So we can choose $\delta > 0$, no larger than $1$, to be smaller than the distance from $c$ to the nearest of them.  Every $x$ with $|x - c| < \delta$ then satisfies $|t(x) - t(c)| = t(x) < \varepsilon$.

So the answer to our question is yes.  The picture makes the proof visible.  Draw a horizontal line at height $\varepsilon$.  Only finitely many dots lie on or above it, and an irrational $c$ is not under any of them, so there is room around $c$ where the function stays below the line.

Curiously, the reverse arrangement is impossible.  No function on $\mathbb{R}$ is continuous at every rational number and discontinuous at every irrational one.  The proof needs ideas beyond this chapter, but the reason can be stated.  The set of discontinuities of any function is a countable union of closed sets, and the set of irrational numbers is not.

\subsection{Three Kinds of Discontinuity}

To sort discontinuities into kinds we approach a point from one side at a time.

\begin{definition}
Let $f : A \to \mathbb{R}$, and let $c$ be a limit point of $A \cap (c, \infty)$.  We say that the \textsc{right-hand limit} of $f$ at $c$ is $L$, and write $\lim_{x \to c^+} f(x) = L$, if for every $\varepsilon > 0$ there exists a $\delta > 0$ such that $0 < x - c < \delta$ and $x \in A$ imply $|f(x) - L| < \varepsilon$.  The \textsc{left-hand limit} $\lim_{x \to c^-} f(x)$ is defined in the same way, with $0 < c - x < \delta$.
\end{definition}

The ordinary limit $\lim_{x \to c} f(x)$ exists exactly when both one-sided limits exist and are equal.  If $f$ is discontinuous at $c$, then one of three things has gone wrong.
\begin{itemize}
	\item The limit $\lim_{x \to c} f(x)$ exists but is not equal to $f(c)$.  This is a \textsc{removable discontinuity}, because changing the single value $f(c)$ would repair it.
	\item Both one-sided limits exist but they are different.  This is a \textsc{jump discontinuity}.  The postage function has one at $20$.
	\item At least one of the one-sided limits does not exist.  This is an \textsc{essential discontinuity}.  The function $\sin(1/x)$ has one at $0$, however we define it there, and Dirichlet's function has one at every point.
\end{itemize}

\begin{figure}
\centering
\begin{tikzpicture}
  % removable
  \begin{scope}[x=1cm,y=1cm]
    \draw[graphaxis] (0,0) -- (3,0);
    \draw[graphaxis] (0,0) -- (0,2.6);
    \draw[curve] plot[domain=0.3:2.7, samples=30] (\x,{0.5+0.22*\x*\x});
    \node[openpt] at (1.5,0.995) {};
    \node[closedpt] at (1.5,1.9) {};
    \draw[guide] (1.5,0) -- (1.5,0.85);
    \node[lbl,below=2pt] at (1.5,0) {$c$};
    \node[lbl] at (1.5,-0.75) {removable};
  \end{scope}
  % jump
  \begin{scope}[shift={(3.75,0)},x=1cm,y=1cm]
    \draw[graphaxis] (0,0) -- (3,0);
    \draw[graphaxis] (0,0) -- (0,2.6);
    \draw[curve] plot[domain=0.3:1.5, samples=20] (\x,{0.45+0.3*\x});
    \draw[curve] plot[domain=1.5:2.7, samples=20] (\x,{1.3+0.35*\x});
    \node[closedpt] at (1.5,0.9) {};
    \node[openpt] at (1.5,1.825) {};
    \draw[guide] (1.5,0) -- (1.5,0.75);
    \node[lbl,below=2pt] at (1.5,0) {$c$};
    \node[lbl] at (1.5,-0.75) {jump};
  \end{scope}
  % essential
  \begin{scope}[shift={(7.5,0)},x=1cm,y=1cm]
    \draw[graphaxis] (0,0) -- (3,0);
    \draw[graphaxis] (0,0) -- (0,2.6);
    \draw[curve] plot[domain=0.3:1.5, samples=20] (\x,{0.6+0.35*\x});
    \draw[curve, line width=0.6pt] plot[domain=1.1:90, samples=900, variable=\t] ({1.5+1.3/\t},{1.3+0.8*sin(deg(0.55*\t))});
    \node[closedpt] at (1.5,1.125) {};
    \draw[guide] (1.5,0) -- (1.5,0.98);
    \node[lbl,below=2pt] at (1.5,0) {$c$};
    \node[lbl] at (1.5,-0.75) {essential};
  \end{scope}
\end{tikzpicture}
\caption{The three ways a function can fail to be continuous at a point $c$.}
\label{fig:discontinuities}
\end{figure}

\subsection{Monotone Functions}

A function $f$ is \textsc{increasing} on a set if $x < y$ implies $f(x) \leq f(y)$, \textsc{decreasing} if $x < y$ implies $f(x) \geq f(y)$, and \textsc{monotone} if it is one or the other.  Monotone functions are the functions that never change direction, and that restriction tames their discontinuities considerably.

\begin{theorem}\label{thm:monotone-limits}
Let $f$ be increasing on an open interval $(a,b)$.  Then at every point $c \in (a,b)$ both one-sided limits exist, and
\[ \lim_{x \to c^-} f(x) \ \leq\ f(c)\ \leq\ \lim_{x \to c^+} f(x). \]
In particular, the only discontinuities an increasing function can have are jumps.  The same is true of decreasing functions, with the inequalities reversed.
\end{theorem}

\begin{proof}
Let $c \in (a,b)$.  The set $\{f(x) : a < x < c\}$ is nonempty and is bounded above by $f(c)$, because $f$ is increasing.  Let $L$ be its supremum, so that $L \leq f(c)$.  We claim that $L$ is the left-hand limit.

Let $\varepsilon > 0$.  Since $L - \varepsilon$ is not an upper bound, there is a point $x_0$ with $a < x_0 < c$ and $f(x_0) > L - \varepsilon$.  Put $\delta = c - x_0$.  If $0 < c - x < \delta$, then $x_0 < x < c$, and so
\[ L - \varepsilon < f(x_0) \leq f(x) \leq L. \]
Hence $|f(x) - L| < \varepsilon$, as required.  The right-hand limit is the infimum of $\{f(x) : c < x < b\}$, by the same argument.
\end{proof}

This is the Monotone Convergence Theorem in a new setting.  A monotone function, like a monotone sequence, has nowhere to go but toward its supremum.

\begin{theorem}\label{thm:monotone-countable}
A monotone function on an open interval has at most countably many discontinuities.
\end{theorem}

\begin{proof}
Let $f$ be increasing on $(a,b)$.  The decreasing case is the same.  At a point $c$ where $f$ is discontinuous, the two one-sided limits in Theorem~\ref{thm:monotone-limits} cannot be equal, so there is a gap between them.  Since $\mathbb{Q}$ is dense in $\mathbb{R}$ we can choose a rational number $r(c)$ inside the gap,
\[ \lim_{x \to c^-} f(x) \ <\ r(c)\ <\ \lim_{x \to c^+} f(x). \]
Different discontinuities receive different rational numbers.  Indeed, if $c < d$ are two of them, pick any point $z$ with $c < z < d$.  Then
\[ r(c) < \lim_{x \to c^+} f(x) \leq f(z) \leq \lim_{x \to d^-} f(x) < r(d). \]
So $r$ is a one-to-one map from the set of discontinuities of $f$ into $\mathbb{Q}$.  Since $\mathbb{Q}$ is countable, the set of discontinuities is at most countable.
\end{proof}

\begin{marginfigure}
\centering
\begin{tikzpicture}[x=0.88cm,y=1cm]
  \draw[graphaxis] (0,0) -- (4.7,0);
  \draw[graphaxis] (0,0) -- (0,3.5);
  \draw[curve] (0.3,0.49) -- (1.6,0.88);
  \draw[curve] (1.6,1.5) -- (3.0,1.85);
  \draw[curve] (3.0,2.6) -- (4.3,2.99);
  \node[closedpt] at (1.6,0.88) {};
  \node[openpt] at (1.6,1.5) {};
  \node[closedpt] at (3.0,1.85) {};
  \node[openpt] at (3.0,2.6) {};
  \draw[guide] (1.6,0.88) -- (0,0.88);
  \draw[guide] (1.6,1.5) -- (0,1.5);
  \draw[guide] (3.0,1.85) -- (0,1.85);
  \draw[guide] (3.0,2.6) -- (0,2.6);
  \draw[guide] (1.6,0) -- (1.6,0.88);
  \draw[guide] (3.0,0) -- (3.0,1.85);
  \draw[nbhdline] (0,0.88) -- (0,1.5);
  \draw[nbhdline] (0,1.85) -- (0,2.6);
  \node[closedpt=calm] at (0,1.2) {};
  \node[closedpt=calm] at (0,2.25) {};
  \node[lbl,left=3pt,calm] at (0,1.2) {$r(c)$};
  \node[lbl,left=3pt,calm] at (0,2.25) {$r(d)$};
  \node[lbl,below=2pt] at (1.6,0) {$c$};
  \node[lbl,below=2pt] at (3.0,0) {$d$};
\end{tikzpicture}
\caption{Each jump of an increasing function leaves a gap on the vertical axis.  The gaps do not overlap, and each one contains a rational number of its own.}
\end{marginfigure}

A monotone function, then, cannot behave like Dirichlet's.  It is continuous at all but countably many points.  This is still a great deal of freedom, since $\mathbb{Q}$ itself is countable, and a monotone function can indeed have a jump at every rational number.

\section{Problems}

\noindent The problems in this section are original to this text.

\begin{problem}[Answering the client]
\hint{In (b) and (c), first promise that $\delta$ will be at most $1$, or at most $\frac12$, and see what bound that gives you for the awkward factor.}
For each limit, find a response $\delta$ to the challenges $\varepsilon = 1$ and $\varepsilon = \frac{1}{10}$, and then to a general $\varepsilon > 0$.  Write out the proof for the general case.
\begin{enumerate}[(a)]
	\item $\lim_{x \to 4} (5 - 2x) = -3$
	\item $\lim_{x \to 1} (x^2 + x) = 2$
	\item $\lim_{x \to 2} \dfrac1x = \dfrac12$
	\item A student answers part (b) with $\delta = \varepsilon/5$ and no further conditions.  Show that this response is correct when $\varepsilon = 1$ but fails when $\varepsilon = 100$.  For which $\varepsilon$ does it work?
\end{enumerate}
\end{problem}

\begin{problem}[Scrambled definitions]
\hint{For (c), the Extreme Value Theorem is useful.}
Each of the following is the definition of continuity at $c$ with something rearranged.  Let $f : \mathbb{R} \to \mathbb{R}$ and $c \in \mathbb{R}$.
\begin{itemize}
	\item $f$ is \emph{frozen} at $c$ if there exists a $\delta > 0$ such that for every $\varepsilon > 0$, $|x - c| < \delta$ implies $|f(x) - f(c)| < \varepsilon$.
	\item $f$ is \emph{fenced} at $c$ if for every $\delta > 0$ there exists an $\varepsilon > 0$ such that $|x - c| < \delta$ implies $|f(x) - f(c)| < \varepsilon$.
	\item $f$ is \emph{revealing} at $c$ if for every $\varepsilon > 0$ there exists a $\delta > 0$ such that $|f(x) - f(c)| < \varepsilon$ implies $|x - c| < \delta$.
\end{itemize}
\begin{enumerate}[(a)]
	\item Describe each property in a plain sentence that mentions neither $\varepsilon$ nor $\delta$.
	\item Show that a function frozen at $c$ is continuous at $c$, and that $f(x) = x$ is continuous at every point but frozen at none.
	\item Show that a function continuous on all of $\mathbb{R}$ is fenced at every point.  Then show that Dirichlet's function is fenced at every point as well, so that being fenced tells us very little.
	\item Give an example of a function that is revealing at $0$, and a continuous function that is not.
\end{enumerate}
\end{problem}

\begin{problem}[Continuity by appointment]
\hint{In each part, multiply Dirichlet's function by something continuous.}
Example~\ref{ex:one-point} showed that $h(x) = x\,g(x)$ is continuous at $0$ and nowhere else, where $g$ is Dirichlet's function.
\begin{enumerate}[(a)]
	\item Let $\varphi : \mathbb{R} \to \mathbb{R}$ be continuous.  Prove that the product $\varphi\,g$ is continuous at a point $c$ if and only if $\varphi(c) = 0$.
	\item Find a function that is continuous at $1$ and at $-1$ and at no other point.
	\item Find a function that is continuous at every integer and at no other point.
\end{enumerate}
\end{problem}

\begin{problem}[What a continuous function does to a set]
\hint{Statement (c) deserves care.  Compare it with the example of $1/x$ on $(0,1)$ in the text, and ask what is different.}
Let $f : \mathbb{R} \to \mathbb{R}$ be continuous on all of $\mathbb{R}$.  For a set $B$, write $f^{-1}(B) = \{x : f(x) \in B\}$.  Decide whether each statement is always true.  Prove the true ones and give a counterexample to each of the others.
\begin{enumerate}[(a)]
	\item If $B$ is open, then $f(B)$ is open.
	\item If $B$ is closed, then $f(B)$ is closed.
	\item If $B$ is bounded, then $f(B)$ is bounded.
	\item If $B$ is an interval, then $f(B)$ is an interval.
	\item If $B$ is open, then $f^{-1}(B)$ is open.
	\item If $B$ is compact, then $f^{-1}(B)$ is compact.
\end{enumerate}
\end{problem}

\begin{problem}[The mountain path]
\hint{Consider the difference between her altitude on the way up and her altitude on the way down, as a function of the time of day.}
A hiker leaves the trailhead at 8 in the morning and reaches the summit at 6 in the evening.  She camps there, sets off down the same path at 8 the next morning, and is back at the trailhead at 6 in the evening.  She does not walk at a steady pace on either day, and she stops whenever she likes.
\begin{enumerate}[(a)]
	\item Prove that there is a time of day at which she was at exactly the same altitude on both days.
	\item Which property of walking did your proof depend on?  Describe a fanciful way of travelling for which the conclusion would be false.
\end{enumerate}
\end{problem}

\begin{problem}[The runner]
\hint{Let $d(t)$ be the distance covered after $t$ minutes and look at $d(t+5) - d(t)$.  What do the values of this function at $t = 0, 5, 10, \ldots, 45$ add up to?}
A runner covers $10$ kilometers in exactly $50$ minutes.  Her speed varies, but the distance she has covered is a continuous function of time.
\begin{enumerate}[(a)]
	\item Prove that there is some $5$-minute stretch of the run during which she covers exactly $1$ kilometer.
	\item Does the same argument show that there is a $7$-minute stretch during which she covers exactly $1.4$ kilometers?  Where does it break down?
\end{enumerate}
\end{problem}

\begin{problem}[Uniform or not]
\hint{For (c), first show that $|\sqrt{x} - \sqrt{y}| \leq \sqrt{|x - y|}$.}
Decide whether each function is uniformly continuous on the given set, and prove your answer.
\begin{enumerate}[(a)]
	\item $f(x) = \dfrac1x$ on $[1, \infty)$
	\item $f(x) = \dfrac{1}{x^2}$ on $(0,1]$
	\item $f(x) = \sqrt{x}$ on $[0, \infty)$
	\item $f(x) = x^3$ on $\mathbb{R}$
	\item Prove that a function which is uniformly continuous on a \emph{bounded} set $A$ must be a bounded function.  Which of the examples above does this explain?
\end{enumerate}
\end{problem}

\begin{problem}[A second proof by sequences]
\hint{Begin with the sequences supplied by Theorem~\ref{thm:not-uniform} and use compactness to make one of them converge.}
In the text, the theorem that a continuous function on a compact set $K$ is uniformly continuous was proved with open covers.  Give a different proof that uses sequences.
\begin{enumerate}[(a)]
	\item Suppose $f$ is continuous on $K$ but not uniformly continuous, and derive a contradiction.
	\item Run your argument for $f(x) = \frac1x$ on the set $(0,1)$, which is not compact, and identify the exact step that fails.
\end{enumerate}
\end{problem}

\begin{problem}[Popcorn and staircases]
\hint{For (a), reread the proof that $t$ is continuous at irrational points and notice how little of it used the fact that $c$ was irrational.}
\begin{enumerate}[(a)]
	\item Let $t$ be Thomae's function.  Prove that $\lim_{x \to c} t(x) = 0$ at \emph{every} real number $c$, rational or not.  What kind of discontinuity does $t$ have at a rational point?
	\item Define $s$ on $[0,1]$ by $s(0) = 0$ and $s(x) = \frac1n$ when $\frac{1}{n+1} < x \leq \frac1n$.  Sketch the graph of $s$.  At which points is $s$ discontinuous, and what kind of discontinuity does it have there?
	\item Is $s$ continuous at $0$?  Prove your answer.
\end{enumerate}
\end{problem}


\backmatter

\bibliography{sample-handout}
\bibliographystyle{plainnat}


\printindex

\end{document}

